\documentclass[11pt]{article}
\usepackage[T1]{fontenc}
\usepackage[utf8]{inputenc}
\usepackage{lmodern,amsmath,amssymb,amsthm,mathtools}
\usepackage[a4paper,margin=29mm]{geometry}
\usepackage{microtype,booktabs,array,longtable,enumitem}
\usepackage[hidelinks]{hyperref}
\usepackage{xurl}
\setlength{\emergencystretch}{2em}
% Fixed companion references, included for a standalone build.
% Static statement numbers from the companion paper; no external build file required.
\providecommand{\compref}[1]{\csname external@#1\endcsname}
\providecommand{\compeqref}[1]{(\compref{#1})}
\expandafter\def\csname external@sec:analytic-weil-sheaves\endcsname{B.8}

\setlist[enumerate]{leftmargin=*,itemsep=.35em}
\setlist[itemize]{leftmargin=*,itemsep=.25em}
\newtheorem{theorem}{Theorem}[section]
\newtheorem{proposition}[theorem]{Proposition}
\newtheorem{lemma}[theorem]{Lemma}
\newtheorem{corollary}[theorem]{Corollary}
\theoremstyle{definition}
\newtheorem{definition}[theorem]{Definition}
\newtheorem{example}[theorem]{Example}
\newtheorem{question}[theorem]{Question}
\theoremstyle{remark}
\newtheorem{remark}[theorem]{Remark}
\newcommand{\Q}{\mathbf Q}
\newcommand{\Z}{\mathbf Z}
\newcommand{\R}{\mathbf R}
\newcommand{\C}{\mathbf C}
\newcommand{\Qbar}{\overline{\mathbf Q}}
\newcommand{\A}{\mathcal A}
\newcommand{\cl}{\operatorname{cl}}
\newcommand{\Nm}{\operatorname{Nm}}
\newcommand{\Spec}{\operatorname{Spec}}
\newcommand{\Hilb}{\operatorname{Hilb}}
\newcommand{\Gal}{\operatorname{Gal}}
\newcommand{\workbenchurl}{https://github.com/KokunoYumeto/algebraic-cycles-and-motives-workbench}
\newcommand{\recordurl}{https://github.com/KokunoYumeto/mathematics-workbenches}
\title{The Hodge conjecture for abelian varieties as a uniform bound over finite fields}
\author{}
\date{}
\begin{document}
\maketitle
\begin{abstract}
Br\"unjes and Serp\'e proved an arithmetic lifting criterion for cycle
classes under a uniform complexity bound. We express the Hodge conjecture
for all complex abelian varieties as one uniform degree bound on
reductions of CM points in each special family. The equivalence uses
published inputs only; recent unrefereed CM theorems are not inputs.
The auxiliary primes and geometric places may vary, and an algebraic
pairing bounds the denominator automatically. A single Weil detector
gives the entire plane at quadratic cost. Integral isogenies reduce
split families to one principal hyperbolic family per field and
dimension, and the fourfold bounds are $48$, $192$ and $36$.
An exact dimension criterion characterizes strong semiregularity
for a pure Weil character. Higher-dimensional numerical estimates
identify their preprint inputs explicitly. No general detector
construction or proof of the Hodge conjecture is obtained.
The mathematics and text were produced by GPT-6 Astra, with separate
re-proofs by the same model and earlier checks by Claude Opus 5.5
at high effort; no human has checked them.
\end{abstract}
\section{The conjecture and the uniformity it requires}
\label{sec:introduction}

The Hodge conjecture is known for complex abelian varieties of dimension
at most five. Voisin presents this as Markman's theorem in
\cite[Th\'eor\`eme~1.3, p.~1248--03]{VoisinB}, reducing it to the
split-sixfold Weil theorem in Th\'eor\`eme~1.4 on the same page.
The reduction uses the Mumford--Tate analysis of Deligne, Tankeev and
Moonen--Zarhin; Markman's account also credits Ram\'on Mar\'i for products
of abelian surfaces \cite[Corollary~1.3, pp.~587--588]{MarkmanICM}.
Markman's earlier split-fourfold theorem, reproved by Floccari--Fu,
and Schoen's results for particular quadratic fields are the antecedents
described by Voisin on p.~1248--03.

Here a split Hermitian space of dimension $2n$ has determinant $(-1)^n$.
Voisin uses the signed discriminant, equal to one; Markman's sixfold
determinant is $-1$. The conventions describe the same case. Our
contribution in dimension four is quantitative: we bound degrees of
specified detecting cycles and keep track of the integral Weil lattice.
In every dimension we give a uniform-bound criterion for the Hodge
conjecture.


The rational Hodge conjecture says that every rational cohomology class of type $(p,p)$ on a smooth projective complex variety is a rational linear combination of codimension-$p$ algebraic cycles. Even for abelian varieties, this statement asks for algebraic cycles in situations where no construction is known. We show that, for abelian varieties, the existence question is equivalent to uniform boundedness in special families, and that this boundedness can be tested over finite fields.

The arithmetic antecedent is Br\"unjes--Serp\'e \cite[Definition 2.3(ii), p.~4; Theorem 3.4, p.~8]{BS}. They bound both the effective numerator cycles and the integer denominator, use a fixed auxiliary prime, and work with cycles over finite residue fields. The form developed here permits cycles over algebraic closures of residue fields, permits the auxiliary prime to vary, and makes the denominator bound automatic. For abelian varieties an algebraic pairing gives an explicit linear bound for that denominator. The proof below constructs the parameter spaces directly and does not invoke their Theorem~3.4. The passage from one variety to special families relates this uniformity principle to the Hodge conjecture for all complex abelian varieties.

The complex argument also has a precise antecedent. In the opening page of Cattani--Deligne--Kaplan \cite[p.~483]{CDK}, the implication from the Hodge conjecture to algebraicity of Hodge loci is explained by bounded numerator degrees, a bounded denominator, Baire's theorem and Chow parameter spaces. Their Theorem~1.1 proves algebraicity of the Hodge locus unconditionally. We use the bounded-cycle argument for a global flat class and add arithmetic specialization. Bhattacharjee--Bhattacharya \cite[Lemmas~4.2--4.3 and Theorem~4.5]{BBDense} and \cite[Theorems~16.9 and 16.17]{BBAlgebraic} also formulate closed bounded-cycle loci and dense-set criteria; those recent preprints concern the qualitative complex step. Voisin \cite[Propositions~0.2 and 0.5, Theorem~0.6]{Voisin07} studies fields of definition of Hodge loci, which is a distinct issue from the cycle parameter spaces used below.

\subsection{Conventions and the main theorem}

If $(A,\theta)$ is a polarized abelian $g$-fold and $Z$ is an effective cycle of pure dimension $d$, set
\[
 \deg_\theta Z=\int_Z\theta^d.
\]
An effective integral cycle means a finite sum of integral subvarieties with nonnegative integer coefficients; the zero cycle is allowed. The word ``integral'' here refers to coefficients and does not require the cycle to be irreducible. When an irreducible reduced subvariety is intended, we say so explicitly.

\begin{definition}\label{def:special}
A \emph{special family with a Hodge class} consists of a connected component $S$ of a Shimura variety at neat level, a map
\[
 S\longrightarrow\mathcal A_{g,N},\qquad N\geq3,
\]
induced by a morphism of Shimura data, the pulled-back principally polarized abelian scheme $f:\A\to S$ with polarization $\theta$, and a nonzero global section
\[
 t\in H^0(S,R^{2p}f_*\Q(p))
\]
of type $(0,0)$ after the Tate twist, or equivalently type $(p,p)$ before the twist, at every point. For $s\in S(\C)$ define
\[
 c_t(s)=\min\left\{\max\{M,\deg_\theta Z^+,\deg_\theta Z^-\}:
 M t_s=\cl(Z^+)-\cl(Z^-),\ M\in\Z_{>0}\right\}.
\]
The cycles in this minimum are effective integral cycles of codimension $p$. The value is $+\infty$ if the set is empty. In this article a CM point of $S$ means a special point of its Shimura datum. Its image in the Siegel variety is then a CM abelian variety.
\end{definition}

A number-field model of a fibre means a polarized abelian variety over a number field $L\subset\C$ together with a specified identification of its complex fibre. It does not mean that $t_s$ is Galois invariant. At a good place $v$ of $L$, a geometric place above $v$ and a prime $\ell\ne\operatorname{char}k(v)$ determine the usual comparison and smooth proper specialization of $t_s$ to $\ell$-adic cohomology. The hypothesis on the map of Shimura data ensures that CM points have number-field models.

\begin{theorem}[Family criterion]\label{thm:family}
For a special family $(S,\theta,t)$, the following conditions are equivalent.
\begin{enumerate}[label=\textnormal{(\roman*)}]
\item Every $t_s$, for $s\in S(\C)$, is algebraic.
\item The function $c_t$ is bounded on $S(\C)$.
\item The function $c_t$ is bounded on one Zariski-dense subset of $S(\C)$.
\item There is an integer $D\geq1$ such that, for every CM point $x$ and every number-field model of its fibre, infinitely many good places $v$ admit a geometric place above $v$, a prime $\ell\ne\operatorname{char}k(v)$, an integer $M\geq1$, and effective integral cycles $Z_v^\pm$ over $\overline{k(v)}$ satisfying
\[
 \operatorname{sp}_{v,\ell}(t_x)
 =\frac{\cl_\ell(Z_v^+)-\cl_\ell(Z_v^-)}{M},
 \qquad \deg_\theta Z_v^\pm\leq D.
\]
There is no bound on $M$ in this condition. The choices other than $D$ may depend on $x$ and $v$.
\end{enumerate}
In \textnormal{(iii)} any chosen Zariski-dense subset can be used, including a rational Hecke orbit in the connected component. In \textnormal{(iv)} the good places may be required to belong to any prescribed infinite set. Equivalently, the condition holds at all but finitely many good places of every model, with one uniform bound on the numerator degrees and on $M$, for every allowable $\ell$ and geometric place. ``Every model'' can be replaced by ``some model.''
\end{theorem}

\begin{theorem}[the Hodge conjecture as boundedness]\label{thm:global}
The rational Hodge conjecture holds for every complex abelian variety in every codimension if and only if, in every special family of Definition~\ref{def:special}, $c_t$ is bounded on the CM points. It is also equivalent to condition \textnormal{(iv)} of Theorem~\ref{thm:family} for every such family.
\end{theorem}

The arithmetic bound therefore has the order of quantifiers
\[
 \text{for every family}\quad\exists D\quad\forall x\text{ CM}
 \quad\text{for infinitely many }v\quad\exists(M,Z_v^+,Z_v^-).
\]
Replacing $\exists D\,\forall x$ by $\forall x\,\exists D_x$ gives a different statement: it tests algebraicity separately at the CM points and does not give the conjecture on the other fibres. OpenAI \cite[Theorem~1.1]{OpenAICM} and van Hoften--Sempliner \cite[Theorem~I]{vHS}, following that work, give recent unrefereed proofs of the Hodge conjecture for CM abelian varieties. If those proofs are accepted, pointwise existence at CM points is available; the uniformity in $x$ remains exactly the issue in Theorem~\ref{thm:global}. Neither CM theorem is an input to its proof.


\paragraph{Why this matters.}
If the recent unrefereed CM theorems are accepted, a finite bound exists
at each CM point. The criterion identifies the remaining issue exactly:
one bound must work for all CM points of a given family. For Weil
classes, a detector bound is independent of the chosen class, and
one detecting subvariety gives the whole plane with controlled degrees.

\paragraph{What is new here.}
The special-family formulation with varying geometric places and
auxiliary primes, the characteristic-independent multiplier estimate,
and the quantitative Weil statements are the conclusions developed
here; the closed-locus and Baire argument is already in
Cattani--Deligne--Kaplan \cite[p.~483]{CDK}, qualitative dense-set
criteria are in Bhattacharjee--Bhattacharya
\cite[Lemmas~4.2--4.3 and Theorem~4.5, pp.~10--12]{BBDense},
and the arithmetic uniform-complexity criterion is
Br\"unjes--Serp\'e \cite[Theorem~3.4, p.~8]{BS}.
This comparison makes no claim of priority.
\subsection{Further quantitative statements}

The quantitative part develops the following conclusions.

For a polarized abelian variety of Weil type, an imaginary quadratic field $K$ acts with signature $(n,n)$ on a variety of dimension $2n$. Its rational plane of Weil classes is
$W_K=\bigwedge_K^{2n}H^1(A,\Q)$. Let $\delta_W$ be the minimum degree of an irreducible subvariety with nonzero projection to $W_K$. Theorem~\ref{thm:extraction} proves, for every integral polarization compatible with the action, every nonzero $t\in W_K$, every $\delta\in\Z_{>0}$ with $\delta t$ integral, and every endomorphism $a\in K$ with $a/\bar a$ not a root of unity,
\[
 \delta_W\leq c_t\leq
 \left\lceil\frac{2\,16^n}{n!}\right\rceil
 \delta\lceil\sqrt T\rceil\Nm(a)^{n(2n+1)}\delta_W^2,
 \qquad T=(-1)^n\int t^2.
\]
Compatibility means that the Rosati involution restricts to conjugation on $K$. Infinite values have their usual extended-real interpretation. Thus one detector gives the whole Weil plane at quadratic cost.

Theorem~\ref{thm:small-detectors}, using the cycles of Schoen and van Geemen, gives detector degrees $48$ for even principal Gaussian fourfolds, $192$ for odd principal Gaussian fourfolds, and $36$ for principal Eisenstein fourfolds, all with full integral maximal-order action. In the Eisenstein case $3(W_K\cap H^4(A,\Z))$ is algebraic and
\[
 c_t\leq\max\{3\delta,36\lfloor\delta\sqrt{2T}/3\rfloor\}
 \leq17\delta\lceil\sqrt T\rceil.
\]
The even and odd Gaussian lattices, the integral normalization and the effective cycles are specified in Section~\ref{sec:small}.

Theorem~\ref{thm:stabilization}, after Schoen, says that bounded detectors in every connected split Weil family of dimension $2n+2$ over a fixed imaginary quadratic field imply bounded detectors in every Weil family of dimension $2n$ over that field, for every determinant. No class-number restriction is needed. If $\mathcal O_K=\Z[\tau]$ has class number one, acts integrally, and $\Delta=\chi(L)$, Theorem~\ref{thm:explicit-stabilization} gives
\[
 \delta_W(A,L)\leq\frac{\max\{1,\Nm(\tau)/\Delta\}}{n+1}D_{n+1}(\Delta^2),
\]
where $D_{n+1}(x)$ bounds detectors for all split Weil $(2n+2)$-folds of polarization product $x$. The premise includes every polarization type of that product. The numerator of the displayed factor is one for the Gaussian and Eisenstein fields.

Schoen's cyclic cycles and the dominance theorems of Schoen and Koike give the explicit sixfold estimates in Theorem~\ref{thm:sixfold-suppliers}: $746496$ for the even principal Gaussian lattice, $5971968$ for the odd one, and $30375$ for the principal Eisenstein lattice. Theorem~\ref{thm:principal-quotient} gives a principal hyperbolic quotient over every imaginary quadratic field. Pullback of detectors costs $t_K^nf^n\Delta$, where $\operatorname{Tr}(\mathcal O_K)=t_K\Z$, $f$ is the conductor of the acting order and $\Delta$ is the polarization volume. The ideal classes can also be removed at the additional cost $m_K^n$, where $m_K\le2\sqrt{|\operatorname{disc}K|}/\pi$. Thus one principal family with lattice $U^n$ suffices for each field and dimension; for class number one the additional factor is one. Consequently, for a compatible polarized Weil fourfold of any determinant with integral action by $\mathcal O_f=\Z+f\mathcal O_K$ and $\Delta=\chi(L)$, Theorem~\ref{thm:all-determinant-fourfolds} gives
\[
 \delta_W\le1990656 f^6\Delta^2\quad(K=\Q(i)),
 \qquad\delta_W\le10125 f^6\Delta^2\quad(K=\Q(\omega)).
\]
The proof keeps the source cycles effective and uses published geometric inputs.

Markman \cite[Theorem~1.2, p.~587]{MarkmanICM} states algebraicity for split Weil sixfolds over every imaginary quadratic field, and for all Weil fourfolds; the published account gives a proof sketch and points to the detailed preprint. In the present notation the full content of ICM Theorem~1.2 is:
for every imaginary quadratic field $K$, all $K$-Weil classes on
fourfolds of any determinant and on split sixfolds are algebraic.
The statement includes the phrase ``abelian fourfolds of Weil type
and abelian sixfolds of split Weil type''; the cited complex
multiplication means a field embedding into the rational endomorphism
algebra, not that the abelian variety is of CM type.
The detailed source is \cite[Theorem~1.5.1 and Corollary~1.6.1,
p.~9; \S9.3]{MarkmanSecant}.
The field restriction to nine class-number-one fields belongs to particular numerical estimates, not to this qualitative theorem. Assuming Perry's unrefereed equivariant semiregularity theorem, Theorem~\ref{thm:all-field-sixfolds} gives an explicit bound for every imaginary quadratic field: if $K=\Q(\sqrt{-d})$, $d$ is positive squarefree, $\mathcal O_f$ acts integrally, and the compatible polarization has volume $\Delta$, then every split Weil sixfold satisfies
\[
 \delta_W<10^{179}d^{53}f^3\Delta.
\]
For class number one the coefficient improves to $10^{178}d^{51}$. The sharper constant is an explicitly evaluated Chern-class degree. Split eightfolds and higher dimensions require further arguments.

The $G_{34}$ construction after Severn gives a different source in dimension eight. Theorem~\ref{thm:eisenstein-hurwitz-dominance} proves its dominance from explicit finite data and published inputs. Theorem~\ref{thm:eisenstein-eightfolds} writes out an explicit integer $B_0$ and proves the bounds
\[
 \begin{aligned}
 \delta_W&\le B_0f^4\Delta
 &&\text{for split Eisenstein eightfolds},\\
 \delta_W&\le\frac{B_0}{4}\frac{f^9\Delta^2}{b}
 &&\text{for all Eisenstein sixfolds},
 \end{aligned}
\]
where $b$ is the index of maximal-order saturation. Schoen's cyclic-cover cycles, Severn's exceptional-divisor correction, Brauer induction and Milne's invariant theory carry this proof. The degree calculation keeps integral orbit multiplicities and then uses the principal quotient to reduce the transport cost.


For Gaussian sixfolds of determinant class $[-3]$, Severn's separate
unrefereed $E_8$ construction \cite{SevernGaussian} gives a nonsplit
family. Theorem~\ref{thm:gaussian-volume-bound} gives an explicit
integer $\mathcal B$ such that every compatibly polarized sixfold with
an integral $\Z[i]$-action and this determinant class satisfies
$\delta_W\le(\Delta/3)\mathcal B$, where
$\Delta=\int_A\theta^6/6!$ is a positive multiple of three.
On $\Z[i]^6$ with form $U^2\perp\langle1,-3\rangle$, the bound is
$\mathcal B$.  The proof verifies the determinant suppliers, source
polarization and rank-nine period differential, then bounds degrees
using effective Fermat cycles and invariant divisors.  A classification
of maximal integral Gaussian lattices gives the exact index
$\Delta/3$ for the final pullback.

Theorem~\ref{thm:ghosts} applies to any fixed smooth projective variety over a number field with an ample line bundle and any rational Betti class. Its minimum numerator degree at good reductions is eventually the characteristic-zero minimum if that is finite; otherwise it exceeds every fixed bound outside finitely many places. The multiplier and auxiliary prime may vary. The same dichotomy holds for the cost that also bounds the multiplier.

These results concern algebraic representatives, not merely Hodge-theoretic tensors or motivated classes. Andr\'e's deformation theorem and his theorem that abelian Hodge classes are motivated \cite[Theorems~0.5 and 0.6.2, pp.~8--9]{Andre96} establish powerful propagation results in the motivated category. His CM reduction \cite[Lemma~6.3.2, p.~32]{Andre96} uses split Weil classes over general CM fields. Neither assertion supplies the uniform algebraic cycles required here, and the present criterion does not reduce all abelian Hodge classes to imaginary-quadratic Weil classes.


\paragraph{An exact semiregularity test.}
Let $E$ be a perfect complex on a balanced Weil $2n$-fold, $n\ge2$,
with $e^B\operatorname{ch}(E)=w\in W_K\setminus\{0\}$ for a rational
$(1,1)$-class $B$. Let a finite group act coherently by translations
and tensor products with degree-zero line bundles, with a chosen
equivariant structure on $E$. The ordinary Hochschild semiregularity
map is injective on invariant second self-Ext if and only if
\[
\dim_{\C}\operatorname{Ext}^2(E,E)^G=2n(2n-1).
\]
This is Theorem~\ref{thm:pure-weil-ext}. Simplicity and vanishing of
negative self-Ext are not needed for this equivalence; gluability
is a separate hypothesis in a deformation theorem.

The companion paper \cite{CompanionB} studies particular construction
obstructions and Mumford's exceptional classes. References prefixed B
refer to its numbered statements.
The separate OpenAI manuscript on split eightfolds was withdrawn on
6 October 2026 because its stabilization-trace cancellation has a sign
error; it is not an input here \cite{OpenAIEightWithdrawn}.

\section{Proof overview}\label{sec:overview}

For a fixed degree bound, integral subvarieties have only finitely many Hilbert polynomials. Products of projective Hilbert schemes therefore parametrize all pairs of cycles of those degrees. The equation specifying their cohomology class cuts out unions of connected components. Their images in the base are closed. A countable union of proper closed subvarieties cannot cover a complex variety, while one closed set containing a dense set is the whole variety. These two facts give the complex part of the theorem.

The arithmetic part uses the same finite parameter space over the ring of integers of one fibre. Its components contained in special fibres lie above finitely many places. Outside those places every point specializes from characteristic zero. The cycle class follows this specialization, so one suitable prime outside the exceptional set already suffices. Finally, an algebraic quadratic pairing is positive on all real Hodge classes of the required degree. It bounds $M$ in terms of the numerator degrees and is constant along the family. Primitivity of $t$ is unnecessary.

The proof uses the boundedness and representability theorems of Grothendieck and Kleiman, the trace, comparison and specialization theorems of SGA~4 and SGA~$4\frac12$, and the standard geometry of Shimura varieties. The particular loci below are images of spaces of algebraic cycles and are constructed directly. The argument does not use the Tate conjecture, independence of $\ell$, or a theorem on absolute Hodge classes.


For extraction, an integral polynomial in pullbacks by $a$ and $\bar a$
removes the unwanted eigenspaces of a detector. Two independent Weil
classes then give an integral Gram system. Positivity and degree
estimates bound its solution quadratically. The compatible polarization
controls pullback degree, while the non-root-of-unity condition prevents
the two directions from coinciding.

The fourfold bounds start from geometric cycles: van Geemen's quadric
gives the even Gaussian detector, a specified isogeny gives the odd
case, and Schoen's surface gives the Eisenstein case. Keeping the
principal quotient and the saturated lattice explicit yields both
degree $36$ and the integral multiplier $3$.

Stabilization adds a balanced CM surface with the cancelling determinant
and contracts its factor algebraically. The principal-family reduction
first saturates an isotropic lattice and applies an integral Hermitian
shear. An ideal representative of bounded norm then removes the class
group. Every step is an isogeny with a known polarization multiplier,
so the projection formula gives its degree cost.


The higher-dimensional estimates keep their geometric inputs separate
from the degree calculations. The all-field split-sixfold estimate uses
Markman's secant construction and Perry's equivariant semiregularity
theorem to transport a Chern-class detector. The Eisenstein eightfold
estimate proves dominance for the cover family after Severn and bounds the cost of
integral cyclic suppliers. In the nonsplit Gaussian sixfold family,
the main difficulty is matching the finite determinant identity to
the polarized source and its period map. Exact subgroup and lattice
calculations, a rigorous rank-nine period certificate, and a Fermat
supplier estimate give a numerical bound.  In determinant class
$[-3]$, every integral Gaussian lattice embeds isometrically in one
maximal reference lattice.  The dyadic pairing condition and the
one-class genus argument are essential to this step.  Source cycles
push to the reference family without increasing degree; pulling them
back to a target costs its lattice index $\Delta/3$.
Proper bounded-cycle spaces extend the source bound from an open
period image to every reference fibre.

For the exact semiregularity criterion, a pure Weil character is the
sum of volume forms on complementary subspaces. Its quadratic
annihilator is the mixed part of the exterior algebra. Injectivity of
the trace kills that ideal under evaluation; Serre duality forces the
exact dimension. The group hypotheses ensure that Hochschild
cohomology is fixed and that invariant Serre duality is perfect.
Adding a polarization polynomial changes this calculation.


\section{Proper parameter spaces for bounded cycles}
\label{sec:parameters}

This section isolates the boundedness and cycle-class facts needed in both characteristics. We keep the use of Hilbert schemes explicit because arbitrary nonreduced subschemes of bounded degree need not have bounded Hilbert polynomials.

\begin{lemma}[Cycle parameter spaces]\label{lem:parameters}
Let $T$ be noetherian and let $f:X\to T$ be smooth and projective of relative dimension $g$, with a relatively ample line bundle of class $\theta$. Fix $0\leq d\leq g$ and $D\geq1$. There is a projective $T$-scheme $H_D$ whose geometric points give all pairs of effective integral $d$-cycles of $\theta$-degree at most $D$. Each universal pair has cycle classes which, for $\ell$ invertible on the base, are sections of the pullback of $R^{2g-2d}f_*\Q_\ell(g-d)$. These sections commute with base change. The corresponding assertion holds in Betti cohomology over $\C$. The assertion also holds for a polarized abelian scheme without a chosen line bundle representing its polarization.
\end{lemma}

\begin{proof}
A sufficiently high fixed power of the relatively ample line bundle is relatively very ample. In the abelian case the polarization morphism $\lambda:\A\to\A^\vee$ and the normalized Poincar\'e bundle give, by pullback under $(1,\lambda)$, a relatively ample line bundle of class $2\theta$; one uses a fixed very ample power of this bundle instead. The choice changes the degree bound by one fixed factor and has no effect on the argument.

By SGA~6, XIII, Definition~6.9 and Corollary~6.11 \cite{SGA6}, the Hilbert polynomials of integral $d$-dimensional subvarieties of bounded degree belong to a finite set. An effective cycle of $\theta$-degree at most $D$ contains at most $D$ integral components counted with their multiplicities, since every nonzero integral subvariety has positive integer degree. Enumerate all tuples of these finitely many polynomials for which the sum of the leading degrees is at most $D$ in the $\theta$ normalization. Include the empty tuple for the zero cycle.

For each pair of tuples take the corresponding fibre product of relative Hilbert schemes. Their finite disjoint union is projective over $T$ by FGA, Expos\'e~221, Theorem~3.2 and its specialization to the Hilbert functor \cite{FGA}. At any point take the top-dimensional fundamental cycles of the universal subschemes and add them. All such cycles have the specified degree, including at nonreduced limits, and every effective cycle occurs by using its integral components with repetition. The construction does not claim to parametrize every nonreduced subscheme of bounded degree.

For a universal subscheme $Z$ flat over its parameter scheme $H$, its fibres have dimension $d$. Write $f_H:X_H\to H$ and $h:Z\to H$. Restriction followed by trace gives
\[
 R^{2d}(f_H)_*\Q_\ell(d)\longrightarrow
 R^{2d}h_*\Q_\ell(d)\xrightarrow{\operatorname{Tr}_h}\Q_\ell.
\]
The trace exists and commutes with base change by SGA~4, XVIII, Theorem~2.9 \cite{SGA4}. Poincar\'e duality for the smooth proper map $f_H$, in XVIII, Theorem~3.2.5, identifies this functional with a section of $R^{2g-2d}(f_H)_*\Q_\ell(g-d)$. On every geometric fibre it is the class of the fundamental cycle, with scheme-theoretic multiplicities. Add these sections over the tuple. Their compatibility with algebraic cycle classes is SGA~$4\frac12$, Cycle, construction~2.3.1, Definition~2.3.2, Lemma~2.3.7 and Theorem~2.3.8 \cite{SGA45}. Comparison on the smooth complex fibres is SGA~4, XI, Theorem~4.4, together with the cycle-class compatibility. This also proves the last assertion.
\end{proof}

\begin{lemma}[Closed class loci]\label{lem:closed}
Let $f:X\to S$ be a smooth projective family with a relatively ample class over a connected smooth complex algebraic variety, and let $t$ be a nonzero global rational cohomology section. For each integer $C\geq1$, the locus
\[
 S_C=\{s\in S(\C):c_t(s)\leq C\}
\]
is Zariski closed.
\end{lemma}

\begin{proof}
On the parameter scheme $H_C$ of Lemma~\ref{lem:parameters}, and for a fixed $1\leq M\leq C$, the difference of the two cycle-class sections minus $Mt$ is a section of the pullback of one fixed $\Q_\ell$ local system from $S$. Comparison on each smooth fibre identifies its vanishing with the required Betti equality. Its zero locus is a union of connected components. There are finitely many such components, and each is algebraic and both open and closed. Thus the zero locus is proper over $S$, and its image is closed. The union of these images for the finitely many possible $M$ is precisely $S_C$.
\end{proof}

\begin{lemma}[An unrestricted multiplier is automatically bounded]\label{lem:components-multiplier}
In the setting of Lemma~\ref{lem:closed}, fix $D$. Among all equations
\[
 Mt_s=\cl(Z^+)-\cl(Z^-),\qquad
 M\in\Z_{>0},\quad \deg Z^\pm\leq D,
\]
only finitely many values of $M$ occur, as $s$ and the cycles vary. In particular the locus defined by a bound on the two numerator degrees alone is Zariski closed, even without a Hodge-type assumption on $t$.
\end{lemma}

\begin{proof}
On $H_D$, let $\gamma$ be the difference of the universal cycle classes. Both $\gamma$ and the pullback of $t$ are flat sections. If $\gamma=Mt$ at one point of a connected component, the equality holds throughout that component. Two distinct values $M,M'$ on the same component would give $(M-M')t=0$, contrary to the nonvanishing of the global flat section $t$. There are finitely many connected components of $H_D$, so there are finitely many possible multipliers. The desired locus is the image of the union of the components on which one of these equalities holds; this union is proper over $S$.
\end{proof}

The preceding lemma is qualitative and applies to arbitrary smooth projective families. The pairing in the next section supplies an explicit multiplier bound in the abelian case, without needing to count components of a parameter space.

\begin{proposition}[Uniformity over a complex base]\label{prop:complex}
In Lemma~\ref{lem:closed}, assume that $t$ is a Hodge class on every fibre. Algebraicity on every fibre, boundedness of $c_t$ on all fibres, and boundedness of $c_t$ on one Zariski-dense set are equivalent.
\end{proposition}

\begin{proof}
Algebraicity says that $S(\C)=\bigcup_{C\geq1}S_C$. Since $S$ is smooth and connected, it is irreducible. A proper algebraic closed subset has empty interior in its complex manifold. A nonempty relatively compact analytic coordinate ball is a Baire space after passing to a smaller closed ball; it cannot be a countable union of such subsets. Hence some $S_C$ equals $S$. Conversely, a closed subset $S_C$ containing a Zariski-dense subset equals $S$, and finite $c_t(s)$ means algebraicity. Boundedness on all fibres implies boundedness on any subset.
\end{proof}

\section{A positive algebraic pairing bounds the denominator}
\label{sec:positive}

We now bound the multiplier using a pairing which remains available in positive characteristic. This is the part that turns a degree bound into a complexity bound.

On $A\times A$ let $\theta_1,\theta_2$ denote the pullbacks of $\theta$. For addition $m$ and difference $a$, put
\[
 H=\theta_1+\theta_2,\qquad P=m^*\theta-H.
\]
The theorem of the square gives $a^*\theta=H-P$ and $m^*\theta=H+P$. In particular $H+P$ and $H-P$ are nef.

\begin{lemma}[The Poincar\'e pairing]\label{lem:positive}
Let $t\ne0$ be a real Hodge class of codimension $p$ on a polarized complex abelian $g$-fold and put $d=g-p$. Then
\[
 q_t=\left|\int_{A\times A}(t\boxtimes t)P^{2d}\right|>0.
\]
If $t$ and $\theta$ vary as flat sections in a connected family of Hodge classes, $q_t$ is constant. For effective $d$-cycles $Z,Z'$ over any algebraically closed field,
\begin{equation}\label{eq:nef}
 \left|\int_{Z\times Z'}P^{2d}\right|
 \leq \binom{2d}{d}\deg_\theta Z\deg_\theta Z'.
\end{equation}
Consequently, if $Mt=\cl(Z^+)-\cl(Z^-)$, then
\begin{equation}\label{eq:denominator}
 M^2q_t\leq \binom{2d}{d}
       (\deg_\theta Z^++\deg_\theta Z^-)^2.
\end{equation}
The same inequality holds for a smooth proper specialization of the equation.
\end{lemma}

\begin{proof}
Let $V=H_1(A,\R)$, let $J$ be its complex structure, and let $E$ be the Riemann form of $\theta$. Choose the convention in which $G(u,v)=E(u,Jv)$ is positive definite. Let $x\in\bigwedge^{2d}V$ be Poincar\'e dual to $t$. In symplectic coordinates, the mixed part of $m^*\theta$ is the bilinear form $E$ between the two copies of $V$. Expanding its $2d$-th power gives
\[
 \left|\int_{A\times A}(t\boxtimes t)P^{2d}\right|
 =(2d)!\,\left|(\bigwedge^{2d}E)(x,x)\right|.
\]
Here the exterior-power form is defined on decomposable vectors by
\[
 (\bigwedge^r E)(u_1\wedge\cdots\wedge u_r,
 v_1\wedge\cdots\wedge v_r)=\det(E(u_i,v_j)).
\]
The factorial counts the permutations of the mixed factors; the sign disappears under the absolute value. Since $x$ has homological type $(-d,-d)$, $\bigwedge^{2d}J$ fixes $x$. Therefore
\[
 (\bigwedge^{2d}E)(x,x)
 =(\bigwedge^{2d}G)(x,x)>0.
\]
The last inequality is positive definiteness of the exterior-power metric. Constancy follows because the original integral contains only flat classes.

Write $U=H+P$ and $V'=H-P$. All mixed intersections of $U$ and $V'$ with the effective cycle $Z\times Z'$ are nonnegative. Expanding $(U-V')^{2d}$ and applying the triangle inequality, then using $U+V'=2H$, gives
\[
 \left|\int_{Z\times Z'}P^{2d}\right|
 \leq\int_{Z\times Z'}H^{2d}
 =\binom{2d}{d}\deg_\theta Z\deg_\theta Z'.
\]
Only the term of degree $d$ on each factor survives in the final expansion. This argument is intersection theory with nef divisor classes and works in every characteristic. Apply it to the four products arising from $(Z^+-Z^-)\times(Z^+-Z^-)$ to obtain \eqref{eq:denominator}. Smooth proper specialization preserves the pairing and the intersection numbers, so it preserves the same positive rational number $q_t$ and the same inequality.
\end{proof}

\begin{corollary}\label{cor:denominator}
If both numerator degrees are at most $D$ in a connected family, then
\[
 M\leq 2D\sqrt{\binom{2d}{d}/q_t}.
\]
Thus bounded numerator degrees alone are equivalent to bounded $c_t$ on any subset of that family.
\end{corollary}

\begin{example}\label{ex:elliptic-pairing}
On a principally polarized elliptic curve, take $t=r\theta$, with $r\in\Q_{>0}$, and $p=g=1$. Then $d=0$ and $q_t=r^2$. An equation $Mt=[Z^+]-[Z^-]$ says $Mr=\deg Z^+-\deg Z^-$, so the denominator inequality becomes the elementary estimate $Mr\leq\deg Z^++\deg Z^-$. The same argument in higher dimension uses the Poincar\'e pairing instead of the degree map.
\end{example}

\section{Lifting from one nonexceptional reduction}
\label{sec:lifting}

Properness turns infinitely many reductions into a characteristic-zero cycle. The exceptional set depends on the variety and the degree bound; no bound for its size or largest prime is asserted.

\begin{proposition}[Lifting bounded numerator degrees]\label{prop:lift}
Let $(X,\theta)$ be a smooth projective variety with an ample line bundle over a number field $L\subset\C$, let $t\in H^{2p}(X(\C),\Q(p))$, and fix $D\geq1$. There is a finite set $E_D$ of finite places of $L$ such that the following holds. If $v\notin E_D$ is a good place, and for one geometric place over $v$ and one allowable prime $\ell$ the specialization of $t$ equals
\[
 (\cl_\ell Z^+-\cl_\ell Z^-)/M
\]
with $M\in\Z_{>0}$ and effective cycles of degrees at most $D$, then there are effective cycles $Y^\pm$ on $X_{\Qbar}$ of degrees at most $D$ with
\[
 Mt=\cl(Y^+)-\cl(Y^-)
\]
in Betti cohomology. Neither Galois invariance nor absolute Hodgeness of $t$ is required.
\end{proposition}

\begin{proof}
Spread the smooth projective variety and its ample line bundle over $R=\mathcal O_L[1/b]$ for a nonzero integer $b$. Construct $H_D\to\Spec R$ by Lemma~\ref{lem:parameters}. It has finitely many irreducible components. Those which do not dominate $\Spec R$ are supported over finitely many closed points; put these places, and the excluded places of the model, in $E_D$.

Let $z\in H_D(\overline{k(v)})$ represent the two given cycles. It lies on the closure of a component of the generic fibre. To see the specialization with its geometric choices, pass to the strict henselization $R^{\mathrm{sh}}$ at the chosen geometric place. The reduced structure on a horizontal irreducible component of $H_D$ is torsion-free over the Dedekind base, hence flat. After this base change, every point on it is a specialization of a point in its generic fibre: flat morphisms of finite presentation are open, so every open neighbourhood meets the generic fibre. Choose a geometric generization $\eta$ of $z$ there.

Smooth proper base change with constant coefficients, SGA~4 XVI, Corollary~2.2, identifies the cohomology of the generic and special fibres over the strictly henselian trait. These coefficients satisfy the local-constancy hypothesis added in the SGA~$4\frac12$ erratum. The class $t_\ell$ extends to a constant section on this trait; it need not extend to the original arithmetic base. The universal cycle-class sections commute with specialization. Their equality with $Mt_\ell$ at $z$ therefore holds at $\eta$.

The generic parameter component containing $\eta$, after passage to an algebraically closed characteristic-zero field, descends to a component over $\Qbar$. That component has a $\Qbar$-point, and the classes of its universal cycles are constant on each connected component because the ambient variety is fixed. Choose such a point on the component of $\eta$. Its cycles have the same degrees and satisfy the same $\ell$-adic equation on $X_{\Qbar}$. Comparison with Betti cohomology, and injectivity of extension of scalars from $\Q$ to $\Q_\ell$, give the stated rational equation. All choices can be made compatibly with the original embedding $L\subset\C$ and geometric place. Inertia acts trivially on the cohomology of a good fibre, so the choice of a compatible generization does not alter the class being lifted.
\end{proof}

\begin{proposition}[Spreading an algebraic class]\label{prop:spread}
In Proposition~\ref{prop:lift}, if an equation $Mt=\cl(Y^+)-\cl(Y^-)$ with numerator degrees at most $D$ holds over $\C$, the cycles can be chosen over $\Qbar$ with the same degrees. At all but finitely many good places of any model the specialized equation holds over the residue algebraic closure, for every allowable $\ell$ and geometric place, with the same $M$ and the same degree bound.
\end{proposition}

\begin{proof}
The complex point of $H_D$ giving the cycles belongs to a connected component of $H_{D,\Qbar}$ after extension of scalars. That component has a $\Qbar$-point and has constant cycle classes, giving the descent. Its point is defined over a finite extension $L'/L$. Properness of $H_D$ extends it over every valuation ring of $L'$ above the chosen arithmetic base. Taking the fundamental cycles of the special fibres preserves the degrees. Compatibility of cycle classes with specialization gives the equation. The argument works with the fixed geometric place by choosing a place of $L'$ above it. Only finitely many places were excluded in making the model and its parameter scheme.
\end{proof}

\section{Special families and the proof of the equivalence}
\label{sec:shimura}

Special families connect the complex boundedness statement to all abelian varieties, because every rational Hodge class occurs in one and its CM points are dense.

\begin{lemma}[Every class occurs in a special family]\label{lem:shimura}
Every nonzero rational Hodge class on a complex abelian variety is, after pullback by an isogeny from a principally polarized abelian variety, the value of a section in Definition~\ref{def:special}. CM points are dense in each connected Shimura variety, and a rational Hecke orbit in a fixed connected component is dense.
\end{lemma}

\begin{proof}
A polarized complex abelian variety admits an isogeny from a principally polarized abelian variety. One elementary description is to take a symplectic basis $e_i,f_i$ of its homology lattice with $E(e_i,f_j)=d_i\delta_{ij}$, choose an integer $m$ divisible by all $d_i$, and pass to the sublattice generated by $e_i$ and $(m/d_i)f_i$. The form $E/m$ is integral and unimodular on that sublattice and remains a positive Riemann form. The corresponding torus $B$ is principally polarized and maps by an isogeny $\phi$ to the original variety. If $\phi^*t$ is algebraic, then so is $t$, because $\phi_*\phi^*t=(\deg\phi)t$.

Let $h$ describe the polarized rational Hodge structure on $H_1(B,\Q)$ and let $G$ be its Mumford--Tate group inside the corresponding group of symplectic similitudes. The pair $(G,X)$, where $X$ is the $G(\R)$-conjugacy class of $h$, is a Shimura datum of Hodge type. Indeed the Hodge types in its adjoint representation are $(-1,1),(0,0),(1,-1)$; polarization gives the Cartan involution; and minimality of the Mumford--Tate group excludes an adjoint rational factor on which $h$ is trivial. The inclusion in the symplectic similitude group is a morphism of Shimura data. At neat level it gives the required algebraic map to $\mathcal A_{g,N}$ \cite[\S\S5--7]{MilneSV}.

The tensor $\phi^*t$ with its Tate twist is fixed by $G$. It therefore gives a global section on the component containing $B$, and conjugation of $h$ by $G(\R)$ preserves its Hodge type. This proves the first assertion.

A Shimura datum has special points, namely homomorphisms factoring through rational tori. Their existence follows by approximating a regular element in a suitable real Cartan algebra by a rational regular element and taking its rational centralizer; the resulting real torus is conjugate to the original one. Real approximation for connected rational algebraic groups says that $G(\Q)$ is dense in $G(\R)$ \cite[Theorem~5.4]{MilneSV}. Conjugating one special point by these rational points therefore gives a dense set in each relevant connected component. More generally the same approximation applied to any point gives density of its rational Hecke orbit. A morphism of Shimura data sends special points to special points. In the Siegel variety these are CM abelian varieties and are defined over number fields. This is the only arithmetic property of special points used here.
\end{proof}

\begin{proof}[Proof of Theorems~\ref{thm:family} and \ref{thm:global}]
Proposition~\ref{prop:complex} proves the first three equivalences in the family theorem. If $c_t\leq C$ everywhere, Proposition~\ref{prop:spread} applied to each CM fibre gives the arithmetic condition with $D=C$, even at all but finitely many places and with $M\leq C$. This also gives the assertion for any prescribed infinite set of good places.

Conversely, assume the arithmetic condition with bound $D$. At each CM point, infinitely many places guarantee one outside the finite exceptional set of Proposition~\ref{prop:lift}. That proposition gives a characteristic-zero equation with numerator degrees at most $D$. Corollary~\ref{cor:denominator} gives
\[
 c_t(x)\leq
 \max\left\{D,\left\lceil 2D\sqrt{\binom{2g-2p}{g-p}/q_t}\right\rceil\right\}.
\]
This bound is independent of $x$ because $q_t$ is constant in the family. CM density and the closed-locus argument prove algebraicity at every point. The argument used one auxiliary prime at one suitable place for each $x$ and made no comparison between distinct auxiliary primes. It also shows that one model is enough.

If the rational Hodge conjecture holds for all abelian varieties, condition (i) holds for every special family, so all the stated boundedness conditions follow. In the other direction, Lemma~\ref{lem:shimura} places an arbitrary Hodge class, after an isogeny, in such a family. The family theorem makes that pullback algebraic, and isogeny pushforward makes the original class algebraic.
\end{proof}

\section{Minimum degrees and reductions}
\label{sec:ghosts}

The same properness argument describes exactly what bounded-degree reduction can and cannot hide for one fixed variety and one fixed class.

Let $(X,\theta)$ be a smooth projective variety with an ample line bundle over a number field and fix a rational Betti class $t$. Define $b_0$ to be the least $D$ such that $Mt=\cl(Z^+)-\cl(Z^-)$ for some $M\geq1$ and effective cycles of degrees at most $D$ over $\C$. Set $b_0=+\infty$ if there is no such equation. Define $b_v$ similarly on the geometric good reduction, allowing a choice of geometric place and auxiliary prime as in Theorem~\ref{thm:family}.

\begin{theorem}[Eventual minimum-degree dichotomy]\label{thm:ghosts}
If $b_0<\infty$, then $b_v=b_0$ for all but finitely many good places $v$. If $b_0=\infty$, then for each integer $D$ only finitely many good places satisfy $b_v\leq D$. The same statements hold when the cost also bounds $M$, as in $c_t$.
\end{theorem}

\begin{proof}
For $t=0$ all the numerator-degree minima are zero, and all the minima for the cost bounding $M$ are one, so both assertions are immediate. Suppose $t\ne0$. When $b_0<\infty$, choose a characteristic-zero equation realizing that minimum. Proposition~\ref{prop:spread} gives $b_v\leq b_0$ outside finitely many places. If $b_0\geq2$, Proposition~\ref{prop:lift} with $D=b_0-1$ gives $b_v\geq b_0$ outside another finite set. If $b_0=1$, smooth proper specialization makes $b_v=0$ impossible. If $b_0=\infty$, infinitely many places with $b_v\leq D$ would give a characteristic-zero equation by Proposition~\ref{prop:lift}, a contradiction. For the cost bounding $M$, the same lifting preserves the specific multiplier and the same spreading preserves it, so the argument is identical.
\end{proof}

Thus a class which is algebraic at every good reduction but not in characteristic zero would have reduction degrees escaping every fixed bound. This is a conditional description of such a class, not an existence assertion. Moreover, for fixed $A$ and $D$, one good place outside the finite set $E_D$ suffices for lifting. Making that finite set effective and uniform in a CM family is a separate problem.

\section{Weil classes and their detecting cycles}
\label{sec:weil}
This section studies the quantitative form of algebraicity for the two-dimensional space of Weil classes. Let $K$ be an imaginary quadratic field acting on a polarized abelian variety $A$ of dimension $2n$, with signature $(n,n)$ and with the Rosati involution restricting to complex conjugation on $K$. Its Weil space is the rational plane
\[
 W_K(A)=\bigwedge_K^{2n}H^1(A,\Q)\subset H^{2n}(A,\Q).
\]
It consists of Hodge classes. Write $e_W$ for its projector in the eigenspace decomposition for the $K$-action, and write $\delta_W(A,\theta)$ for the least degree of an irreducible reduced $n$-dimensional subvariety $Y$ with $e_W\cl(Y)\ne0$, or $+\infty$ if no such subvariety exists. Such a subvariety will be called a detector.

The principal open question for the Weil-type part of this article is whether $\delta_W$ is bounded on the CM points in every connected split Weil family, in every dimension and over every imaginary quadratic field. Here ``split'' means that the rational Hermitian space of dimension $2n$ and signature $(n,n)$ is a sum of hyperbolic planes. Its determinant is $(-1)^n$ modulo norms from $K$; the positive signed discriminant is therefore $1$ modulo norms. We use ordinary intersection degrees throughout.

\subsection{One detector gives the entire plane}

\begin{theorem}[Quadratic extraction]\label{thm:extraction}
Let $n\geq1$ and let $A$, $\theta$ and $K$ be as above, with $\theta$ integral and ample. Let $a\in K$ act by an endomorphism of $A$, assume $a/\bar a$ is not a root of unity, and put $q=\Nm_{K/\Q}(a)$. For a nonzero $t\in W_K(A)$, choose $\delta\in\Z_{>0}$ with $\delta t$ integral and set
\[
 T=(-1)^n\int_A t^2>0.
\]
Then
\begin{equation}\label{eq:sharper-extraction}
 \delta_W\leq c_t\leq
 \left\lceil\frac{2\,16^n}{n!}\right\rceil
 \delta\lceil\sqrt T\rceil\,
 q^{n(2n+1)}\delta_W^2.
\end{equation}
The assertion is valid for every polarization type. In particular it implies
\begin{equation}\label{eq:original-extraction}
 c_t\leq
 \left\lceil\frac{8\,16^n}{n!}\right\rceil
 \delta\lceil\sqrt T\rceil\,
 q^{2n(n+1)}\delta_W^2.
\end{equation}
If no detector exists, both $\delta_W$ and $c_t$ are infinite. Thus the existence of one detector is equivalent to algebraicity of the entire rational Weil plane.
\end{theorem}

The proof has two parts. A polynomial in endomorphism pullbacks removes every middle eigenspace except the Weil plane. Two integral vectors in that positive plane then span a sublattice of finite index. Clearing that index, rather than estimating an arbitrary inverse matrix, keeps the cost quadratic.

\begin{lemma}[Norm and degree]\label{lem:weil-norm}
The Weil plane is primitive for $\theta$, and
\[
 B(u,v)=(-1)^n\int_Au\smile v
\]
is a positive definite integral form on its integral lattice. With
$v_\theta=\int_A\theta^{2n}/(2n)!\in\Z_{>0}$ and
$\|u\|^2=B(u,u)$, every effective $n$-cycle $Y$ satisfies
\[
 \|e_W\cl(Y)\|\leq\frac{\deg_\theta Y}{n!\sqrt{v_\theta}}.
\]
For any chosen signed presentation $C=C^+-C^-$ whose class lies in $W$, its norm is at most
$(\deg C^++\deg C^-)/(n!\sqrt{v_\theta})$.
\end{lemma}

\begin{proof}
Over $\C$ write $V=H^1(A,\C)=V_\sigma\oplus V_{\bar\sigma}$ for the two eigenspaces of $K$, each of dimension $2n$. The signature hypothesis makes the extreme determinant lines
\[
 W_\C=\bigwedge^{2n}V_\sigma\oplus\bigwedge^{2n}V_{\bar\sigma}
\]
of Hodge type $(n,n)$. Rosati compatibility gives $a^*\theta=q\theta$, so $\theta$ lies in $V_\sigma\otimes V_{\bar\sigma}$. Its product with either extreme line vanishes for dimension reasons. This proves primitivity.

The primitive Hodge-star formula gives $*w=(-1)^nw$ for $w\in W_\R$, and consequently $B(w,w)=\int_A w\wedge*w>0$. This is the middle-degree Hodge--Riemann form, with the stated sign \cite[Theorem~5.42 and (5.43)]{Ballmann}. The cup product is integral on integral cohomology.

Use the translation-invariant K\"ahler metric determined by $\theta$. Harmonic forms are translation-invariant, so a vector $w$ of cohomological norm one has pointwise norm $v_\theta^{-1/2}$. Wirtinger's formula gives $\operatorname{vol}(Y)=\deg_\theta(Y)/n!$, including for singular $Y$ by integration on its regular locus. Therefore
\[
 \left|\int_Yw\right|\leq\frac{\deg_\theta Y}{n!\sqrt{v_\theta}}.
\]
The cup-product projection onto $W$ is the metric projection: for $w\in W$, its metric pairing with any middle-degree class $u$ is $\int_Au\wedge*w=B(u,w)$. Taking the supremum over unit $w$ proves the assertion for an effective cycle, and the triangle inequality proves it for a signed cycle.
\end{proof}

\begin{proof}[Proof of Theorem~\ref{thm:extraction}]
The trace of $a$ is integral, and hence $\bar a=\operatorname{Tr}_{K/\Q}(a)-a$ is also an endomorphism. On $H^{2n}$ set
\[
 S=a^*+\bar a^*,\qquad
 \beta_j=q^j\bigl(a^{2n-2j}+\bar a^{2n-2j}\bigr)
 \quad(0\leq j\leq n).
\]
Here $S$ denotes a sum of pullback operators. On the summand with $j$ factors from $V_{\bar\sigma}$ and $2n-j$ from $V_\sigma$, its eigenvalue is $\beta_j=\beta_{2n-j}$. These numbers are integers, and $|\beta_j|\leq2q^n$. If $z=a/\bar a$, equality of two distinct $\beta_j$, for $0\leq j\leq n$, would give $z^{j-k}=1$ or $z^{2n-j-k}=1$. This is excluded by the hypothesis. Thus
\[
 P(X)=\prod_{j=1}^n(X-\beta_j)\in\Z[X],\qquad
 P(S)=P(\beta_0)e_W,
\]
and $P(\beta_0)\ne0$. The middle eigenspaces are orthogonal to the extreme lines because a nonzero top wedge must have exactly $2n$ factors from each $K$-eigenspace. This verifies that this operator is the projector specified above, multiplied by a nonzero integer.

For a chosen signed cycle presentation put
\[
 \mathcal D(C)=\deg_\theta C^++\deg_\theta C^-.
\]
Pullback by $a$ preserves effectivity and multiplies the degree of an $n$-cycle by $q^n$: indeed $\deg(a)=q^{2n}$ and $a^*\theta=q\theta$. It follows that
\[
 \mathcal D((S-\beta_j)C)
 \leq(2q^n+|\beta_j|)\mathcal D(C)
 \leq4q^n\mathcal D(C).
\]
Choose an irreducible detector $Y$ of degree $d=\delta_W$. The integral signed cycle $C=P(S)Y$ has nonzero class $u\in W$ and has a presentation with
\[
 \mathcal D(C)\leq L:=4^nq^{n^2}d,
 \qquad \|u\|\leq L/n!.
\]
Its pullback $a^*C$ has mass at most $q^nL$ and class $v=a^*u$ with $\|v\|=q^n\|u\|$. The vectors $u,v$ are independent over $\Q$: the eigenvalues $a^{2n},\bar a^{2n}$ of $a^*|_W$ are nonreal conjugates, so this rational operator has no rational eigenline.

Let $\mathcal W=W\cap H^{2n}(A,\Z)$. It is a saturated rank-two lattice with integral positive Gram matrix. Its Euclidean covolume $c$ satisfies $c^2\in\Z_{>0}$, so $c\geq1$. Give $W_\R$ the orientation for which $\det(u,v)>0$, and let
\[
 m=[\mathcal W:\Z u+\Z v].
\]
Then $mc=\det(u,v)\leq\|u\|\|v\|$, and $m\mathcal W\subset\Z u+\Z v$. Therefore there exist integers $r,s$ such that
\[
 m\delta t=ru+sv.
\]
Taking determinants with $v$ and with $u$ gives
\[
 r=\frac{\delta\det(t,v)}c,\qquad
 s=\frac{\delta\det(u,t)}c,
\]
and hence
\[
 |r|\leq\delta\sqrt T\,\|v\|,
 \qquad |s|\leq\delta\sqrt T\,\|u\|.
\]
The signed cycle $rC+s a^*C$ represents $M t$, where $M=\delta m$ is a positive integer. Grouping positive and negative coefficients gives effective integral cycles satisfying
\[
 M\leq\frac{\delta q^nL^2}{(n!)^2},\qquad
 \deg Z^\pm\leq|r|L+|s|q^nL
 \leq\frac{2\delta\sqrt T\,q^nL^2}{n!}.
\]
Substitution of $L$ proves \eqref{eq:sharper-extraction}. Since $q\geq1$, it also proves \eqref{eq:original-extraction}.

Finally, if $Mt=\cl(Z^+)-\cl(Z^-)$, at least one irreducible component of one of these cycles has nonzero Weil projection. Its degree is no greater than that of the corresponding effective cycle. Thus $\delta_W\leq c_t$. If no detector exists, this observation excludes every finite representative.
\end{proof}

\begin{remark}[An endomorphism available in every order]\label{rem:order-endomorphism}
If $K=\Q(\sqrt{-m})$, with $m>0$ squarefree, and the endomorphism order contains $\Z+f\mathcal O_K$, one may take $a=f(2+\sqrt{-m})$. Then $q=f^2(m+4)$. Outside $\Q(i)$ and $\Q(\sqrt{-3})$, the only roots of unity in $K$ are $\pm1$, whereas $a/\bar a$ is neither. In the two exceptional fields its real part is respectively $3/5$ and $1/7$, so it is again not a root of unity.
\end{remark}

\subsection{Bounded detector loci}

\begin{proposition}\label{prop:detector-loci}
In a connected polarized family with a fixed integral quadratic action of signature $(n,n)$, the locus $\{s:\delta_W(s)\leq D\}$ is Zariski closed. The following are equivalent:
\begin{enumerate}[label=\textnormal{(\roman*)}]
\item every Weil class on every fibre is algebraic;
\item $\delta_W$ is bounded on every fibre;
\item $\delta_W$ is bounded on one Zariski-dense set.
\end{enumerate}
For a special Weil family these conditions are also equivalent to one uniform detector-degree bound at infinitely many good reductions of every CM fibre, allowing the auxiliary prime and geometric place to vary.
\end{proposition}

\begin{proof}
On the projective parameter space of effective cycles of degree at most $D$, the projected universal class $e_W\gamma$ is a section of the pulled-back Weil local system. Its nonvanishing locus is a union of connected components, hence is also closed and proper over the base. Its image is exactly the stated locus: if an effective sum detects, some irreducible component detects. This proves closedness.

If every Weil class is algebraic, every fibre has a detector. The closed loci therefore cover the complex base; Baire's theorem gives one bound. Conversely, a bound on a dense set gives a bound everywhere by closedness, and Theorem~\ref{thm:extraction} makes every Weil class algebraic. Alternatively, the integral Weil local system has finite monodromy, because it preserves the positive integral form $B$; after a finite \'etale cover it has a basis of global rational sections to which Theorem~\ref{thm:family} applies.

For the arithmetic assertion, extend the quadratic endomorphisms over an arithmetic open of each CM fibre. The projector is a rational polynomial in their pullbacks and commutes with specialization. The proof of Proposition~\ref{prop:lift} therefore lifts a bounded-degree special cycle with nonzero projected class to a bounded-degree characteristic-zero detector, at one place outside the finite exceptional set. Nonvanishing is preserved by the specialization isomorphism. In the reverse direction, descend and spread a characteristic-zero detector as in Proposition~\ref{prop:spread}. This gives detectors at all but finitely many good places.
\end{proof}

\subsection{An explicit Gaussian extraction constant}

\begin{proposition}\label{prop:gaussian-extraction}
On a Gaussian Weil fourfold with integral $\Z[i]$-action, a detector of degree $D$ gives
\[
 c_t\leq61\,440\,\delta\lceil\sqrt T\rceil D^2.
\]
\end{proposition}

\begin{proof}
Put $I=[i]^*$ and $J=[1+i]^*$. On the five summands of $H^4$ indexed by the number of factors from one $K$-eigenspace, direct substitution gives
\[
 (1+I)(4-J)=16e_W.
\]
For a detector $Y$, set $y=16e_W[Y]$. It is represented by
\[
 4Y+4[i]^*Y-[1+i]^*Y-[i(1+i)]^*Y,
\]
whose two effective parts each have degree $8D$. The norm lemma gives $d:=B(y,y)\leq64D^2$. The operator $A=[2+i]^*$ has eigenvalues $-7\pm24i$ on $W$. Hence $z=(A+7)y$ satisfies
\[
 B(y,z)=0,\qquad B(z,z)=576d,
\]
and has signed effective parts of degree at most $256D$. For the integers $b_1=B(\delta t,y)$ and $b_2=B(\delta t,z)$, orthogonal expansion gives
\[
 (576\delta d)t=576b_1y+b_2z.
\]
Write $t=x y/\sqrt d+x' z/(24\sqrt d)$, so $x^2+(x')^2=T$. The degree of each signed part is at most
\[
 \delta D\sqrt d\,(4608|x|+6144|x'|)
 \leq7680\delta\sqrt T D\sqrt d
 \leq61\,440\delta\sqrt T D^2.
\]
The multiplier is at most $36\,864\delta D^2$, which is bounded by the asserted expression.
\end{proof}

\section{Stabilization after Schoen}
\label{sec:stabilization}

This section explains how bounded detectors in the next split dimension control every determinant in the preceding dimension. Schoen's product argument is stated for fourfolds and sixfolds in \cite[Proposition~10, pp.~332--333]{Schoen98}; the general product and contraction statement is also recorded by Voisin \cite[Lemma~2.9]{VoisinB}. We include the degree estimate, for which the polarization and the chosen effective divisors matter.

\begin{theorem}[Qualitative stabilization]\label{thm:stabilization}
Fix an imaginary quadratic field $K$ and $n\geq1$. Suppose that the detector degree is bounded in every connected split Weil family of dimension $2n+2$. Then it is bounded in every connected Weil family of dimension $2n$, of every determinant. The bound may depend on the family and its polarization. No class-number hypothesis is required.
\end{theorem}

\begin{proof}
The determinant class of a fixed connected Weil family is constant. Write it as $(-1)^n d$ modulo norms, with $d=r/s>0$ and $r,s$ positive integers. On $B=E_K\times E_K$, where $E_K=\C/\mathcal O_K$, use the action $a\cdot(x,y)=(ax,\bar a y)$ and the polarization $\lambda=r\vartheta_1+s\vartheta_2$. Its Hermitian determinant is $-rs$ modulo norms, hence $-d$ because $rs=d s^2$. The product $A\times B$, with polarization $\Theta=\theta+\lambda$, has determinant $(-1)^{n+1}d^2$, so its Hermitian space is split by the rational classification \cite[Proposition~4.1 and Corollary~4.2]{Deligne82}. All this data is fixed in the family.

The extreme $K$-eigenlines give an inclusion
\[
 W_K(A\times B)\subset W_K(A)\otimes W_K(B).
\]
Choose an integral basis $\eta_1,\eta_2$ of $W_K(B)\cap H^2(B,\Z)$. The Lefschetz $(1,1)$-theorem makes them divisor classes. For a sufficiently large fixed $N$, the classes
\[
 d_0=N\lambda,\qquad d_1=N\lambda+\eta_1,
 \qquad d_2=N\lambda+\eta_2
\]
are very ample. Choose $b_B>0$ such that $b_B\lambda-d_j$ is nef for each $j$.

Let $Y$ be an effective detector on the split product, of degree at most $D$. Choose effective divisors $D_j$ in the classes $d_j$, translating them to meet $Y$ properly after pullback, and set
\[
 C_j=p_{A*}(Y\cdot p_B^*D_j).
\]
These are effective $n$-cycles. If all three had zero Weil projection, the projection formula and $d_i-d_0=\eta_i$ would make $[Y]$ orthogonal to every $w\otimes\eta_i$, with $w\in W_K(A)$. It would then be orthogonal to $W_K(A\times B)$, contradicting detection. Thus some $C_j$ detects. Mixed nef intersections give
\[
 \deg_\theta C_j
 \leq b_B\int_Yp_A^*\theta^n p_B^*\lambda
 \leq\frac{b_B}{n+1}\int_Y\Theta^{n+1}
 \leq\frac{b_BD}{n+1}.
\]
The factor $n+1$ is the coefficient of $\theta^n\lambda$ in the binomial expansion. Taking an irreducible detecting component proves a uniform bound on $A$.

If the hypothesis is formulated only for principal polarizations, pass to a fixed isogenous principally polarized split family. The choice of a maximal isotropic subgroup of the polarization kernel becomes constant after a finite \'etale cover. The rational $K$-action transports across the isogeny, and an order in $K$ acts integrally. For an isogeny $f:C\to C'$ with $f^*\theta'=k^2\Theta$, a middle-dimensional cycle satisfies
\[
 \deg_\Theta f^*Y'
 =\frac{\deg f}{k^{2n+2}}\deg_{\theta'}Y'.
\]
This is a constant factor in the family. Pullback preserves a nonzero Weil projection, so the preceding contraction applies. The finite cover is surjective, and therefore the bound descends as a fibrewise existence assertion.
\end{proof}

\begin{theorem}[An explicit contraction bound]\label{thm:explicit-stabilization}
Assume $\mathcal O_K=\Z[\tau]$ has class number one and acts integrally on a polarized Weil $2n$-fold $(A,L)$. Put $\Delta=\chi(L)$ and
$\kappa_K(\Delta)=\max\{1,\Nm(\tau)/\Delta\}$. If every split Weil $(2n+2)$-fold of polarization product $x$ has a detector of degree at most $D_{n+1}(x)$, then
\[
 \delta_W(A,L)\leq\frac{\kappa_K(\Delta)}{n+1}D_{n+1}(\Delta^2).
\]
For $K=\Q(i)$ or $\Q(\omega)$ choose $\tau=i$ or $\omega$, respectively; then $\kappa_K(\Delta)=1$. The premise ranges over all polarization types of the specified product, not only principal polarizations.
\end{theorem}

\begin{proof}
Class number one makes $H_1(A,\Z)$ free over $\mathcal O_K$. Normalize the Hermitian form associated to the alternating Riemann form $E$ by
\[
 h(x,y)=E(x,\tau y)-\bar\tau E(x,y).
\]
Here $h$ is linear in its second variable. The first-variable
convention $H$ in Section~\ref{sec:principal-quotients} satisfies
$H=-\overline h$; the determinants agree in even $K$-rank.
In an integral $\mathcal O_K$-basis, $h$ has determinant $(-1)^n\Delta$. Indeed for one diagonal block $h=(a)$, the alternating form on $(v,\tau v)$ has matrix $\left(\begin{smallmatrix}0&a\\-a&0\end{smallmatrix}\right)$. Both sides of $\det_\Z E=(\det_Kh)^2$ transform by the same norm factor under a change of $K$-basis. The sign is fixed by signature and $|\det_Kh|=\chi(L)$.

Take $B=E_K^2$ with balanced action and $\lambda=\vartheta_1+\Delta\vartheta_2$, where $\vartheta_1$ is the class of $\{0\}\times E_K$ and $\vartheta_2$ that of $E_K\times\{0\}$. Its determinant is $-\Delta$, so $A\times B$ is split and has polarization product $\Delta^2$. Let $Y$ be a detector there, of degree at most $D=D_{n+1}(\Delta^2)$ for $\Theta=\theta+\lambda$, and put
\[
 I=\int_Y\theta^n\lambda.
\]
All terms in the expansion of $\Theta^{n+1}$ have nonnegative intersection with $Y$, since both factors are nef. Hence $(n+1)I\le D$.

Intersect with general translates of the coordinate elliptic divisors and push to $A$, as in the preceding proof. Each contraction has degree at most $I$, since $\vartheta_j\le\lambda$ in the nef order. If either has nonzero Weil projection, an irreducible component gives the result.

Suppose both projected contractions vanish. For $u=1,\tau$ use the two graphs $\Gamma_u=\{x=u y\}$ and $\Gamma_{-u}$. Their classes satisfy
\begin{equation}\label{eq:opposite-graphs}
 [\Gamma_u]+[\Gamma_{-u}]
 =2(\vartheta_1+\Nm(u)\vartheta_2)
 \le2\kappa_K(\Delta)\lambda
\end{equation}
in the nef order. The equality follows either by adding the Hermitian matrices
$\left(\begin{smallmatrix}1&\mp u\\\mp\bar u&\Nm(u)\end{smallmatrix}\right)$ or from the equation $x\mp uy=0$. The axes and the graphs of $1,\tau$ span $\operatorname{NS}(B)_\Q$: the axes give the diagonal entries and $1,\tau$ give a rational basis for the off-diagonal entry. Their off-diagonal parts span $W_K(B)$.

Because $Y$ detects $W_K(A\times B)$, contraction against some element of $W_K(B)$ has nonzero projection to $W_K(A)$. Thus one of the two graph pairs has a detecting member. Its two projected contractions are opposite by \eqref{eq:opposite-graphs} and the vanishing of the axis contractions, so both detect. Choose the member of smaller degree. Its degree is at most the average, which by \eqref{eq:opposite-graphs} is at most $\kappa_K(\Delta)I$. One of its irreducible components detects and has no larger degree. This proves the bound.
\end{proof}

\begin{remark}[The polarization range]\label{rem:polarization-range}
Full integral maximal-order action and principal polarization over a class-number-one field force $\Delta=1$ and determinant $(-1)^n$, hence the split determinant class. Statements covering every determinant must therefore allow nonprincipal polarizations or rational actions transported through isogenies. The preceding theorem does so explicitly.
\end{remark}
\section{Small fourfold detectors after Schoen and van Geemen}
\label{sec:small}
This section computes degrees and integral coefficients for the geometric cycles of Schoen and van Geemen. The construction of the cycles is their input; the estimates keep track of the polarization, the isogeny and the saturated Weil lattice.

\begin{theorem}\label{thm:small-detectors}
Let $(A,\theta)$ be a principally polarized Weil fourfold with full integral maximal-order action. For $K=\Q(i)$ the detector degree is at most $48$ on the even Hermitian lattice $U\oplus U$, and at most $192$ on the odd lattice $I_{2,2}$. For $K=\Q(\sqrt{-3})$ it is at most $36$. In the last case, if $\mathcal W=W_K(A)\cap H^4(A,\Z)$, then $3\mathcal W$ consists of integral algebraic cycle classes. For $0\ne t\in W_K(A)$, $\delta\in\Z_{>0}$ with $\delta t$ integral, and $T=\int_A t^2$, one has
\begin{equation}\label{eq:eisenstein-cost}
 c_t\leq\max\left\{3\delta,
 36\left\lfloor\frac{\delta\sqrt{2T}}3\right\rfloor\right\}
 \leq17\delta\lceil\sqrt T\rceil.
\end{equation}
Here $U=\left(\begin{smallmatrix}0&1\\1&0\end{smallmatrix}\right)$ and $I_{r,s}$ is the diagonal Hermitian form with $r$ entries $1$ and $s$ entries $-1$.
\end{theorem}

\subsection{The Gaussian bound in two lattice types}

Van Geemen \cite[Proposition~3.6 and Theorem~3.7, pp.~623--625]{vG96} constructs a dominant rational map from a general member of the even principal Gaussian family to a smooth quadric fourfold in $\mathbf P^5$. The six defining sections belong to $|2\theta|$ and have finite base locus; the inverse image of a suitable ruling plane has nonzero Weil projection. This yields the degree bound with a short residual-intersection argument.

Choose such a plane $\Pi$ generally in its ruling, and choose a linear $\mathbf P^3$ through it. Its intersection with the quadric is the union of $\Pi$ and a residual ruling plane. The two pulled-back hyperplanes have no common divisor for general choices: away from the finite base locus this follows by avoiding the finitely many divisorial components contracted by the rational map. Their proper intersection is an effective surface cycle of class $4\theta^2$ and degree
\[
 \int_A4\theta^4=96.
\]
It contains a detecting component $Y$ from the inverse image of $\Pi$. If $\deg Y\leq48$, this is the desired detector. Otherwise subtract one copy of $Y$ from the complete-intersection cycle. The residual effective cycle has degree less than $48$ and has nonzero Weil projection, since the full class $4\theta^2$ has zero Weil projection. One of its components detects. Proposition~\ref{prop:detector-loci} extends the bound from the general member to every member of the even component.

For the odd lattice put
\[
 P=\begin{pmatrix}1&1\\1&-1\end{pmatrix},\qquad F=P\oplus P.
\]
Direct multiplication gives $P^*UP=2I_{1,1}$ and $\Nm(\det F)=16$. Transporting the complex structure by $F$ gives a $K$-linear isogeny
\[
 f:A_{\rm odd}\longrightarrow A_{\rm even},\qquad
 \deg f=16,\quad f^*\theta_{\rm even}=2\theta_{\rm odd}.
\]
For an effective surface cycle $Z$ on the target,
\[
 \deg_{\theta_{\rm odd}} f^*Z
 =\frac{16}{2^2}\deg_{\theta_{\rm even}}Z.
\]
Pullback preserves nonzero Weil projection, because it is a $K$-linear cohomology isomorphism. Pulling back an even detector and taking a detecting component proves the bound $192$.

Combining these bounds with Proposition~\ref{prop:gaussian-extraction} gives, respectively,
\[
 c_t\leq141\,557\,760\,\delta\lceil\sqrt T\rceil,
 \qquad
 c_t\leq2\,264\,924\,160\,\delta\lceil\sqrt T\rceil.
\]

\subsection{Degree forty-eight on cyclic-four Jacobians, after Schoen}
\label{sec:cyclic-four-degree}

The following calculation keeps the effective sheet cycles in Schoen's
paired-branch construction.  It gives the same numerical bound as the
quadric construction on the even Gaussian family, but applies here to an
odd principal Gaussian Jacobian.

\begin{proposition}\label{prop:cyclic-four-degree}
Let
\[
 C:\quad y^2=z\prod_{r=1}^4(z^2-a_r^2),\qquad
 a_r\ne0,\quad a_r^2\ne a_s^2\quad(r\ne s),
\]
and let $(A,\theta)=(J(C),\theta_C)$, with Gaussian action induced by
$\rho(z,y)=(-z,iy)$.  There are four effective integral surface cycles
$Z_0,Z_1,Z_2,Z_3$ on $A$, each of ordinary theta degree $48$, such that
\[
 0\ne w=[Z_0]-2\theta^2\in W_{\mathbf Q(i)}(A),
 \qquad \int_A w^2=32,
 \qquad [Z_0]-[Z_2]=2w.
\]
Consequently $\delta_W(A,\theta)\le48$.  The integral Gaussian Hermitian
lattice of $(A,\theta)$ is odd.  Thus this improves the bound $192$
on this Jacobian locus; it does not change the bound for every member of
the odd Gaussian family or determine the minimum $\delta_W$.
\end{proposition}

\begin{proof}
Put $t=z^2$ and $b_r=a_r^2$.  The map $C\to\mathbf P^1_t$ is the
cyclic-four cover
\[
 y^4=t\prod_{r=1}^4(t-b_r)^2.
\]
Its branch exponents at $0,\infty,b_1,b_2,b_3,b_4$ are
$(1,3,2,2,2,2)$.  Pair these points as
$(0,\infty),(b_1,b_2),(b_3,b_4)$.  This is the paired-branch
construction of Schoen \cite[Theorem~2.0, pp.~11--12;
Lemmas~2.3--2.6, pp.~13--17]{Schoen88}.  We give the details needed
for its degree and nonvanishing here.

Write $X=\operatorname{Sym}^4C$, and let
$f:X\to\operatorname{Sym}^4\mathbf P^1=\mathbf P^4$ be the finite
map.  The six evaluation forms for the branch points have a unique linear
relation, with every coefficient nonzero.  After rescaling them, write
this relation as
\[
 u_1+v_1+u_2+v_2+u_3+v_3=0.
\]
The plane $P\subset\mathbf P^4$ given by $u_r+v_r=0$ for all three
pairs is a $\mathbf P^2$.  The cyclic norm equation restricts to a
fourth power on $P$: its numerator is a nonzero constant times
$(u_1u_2u_3)^4$.  The intermediate quotient
\[
 W=C^4/(N\rtimes\mathfrak S_4),\qquad
 N=\{(e_1,e_2,e_3,e_4)\in(\mathbf Z/4)^4:\textstyle\sum e_j=0\},
\]
is a cyclic-four cover of $\mathbf P^4$.  Over the open part of $P$
with four distinct, unramified support points, it has four sheets.
Let $B_j$ be the closures, with their generic multiplicities, of their
inverse images in $X$.  They are effective integral surface cycles.
The finite map $f$ is flat, since its source and target are smooth of
the same dimension.  Hence its pullback of $P$ is a complete intersection
whose fundamental cycle is
\begin{equation}\label{eq:cyclic-four-total}
 \sum_{j=0}^3B_j=f^*[P],\qquad
 \sum_{j=0}^3[B_j]=16x^2,
\end{equation}
where $x$ is the class of $p+\operatorname{Sym}^3C$ for a point $p\in C$.
Indeed $f^*\mathcal O_{\mathbf P^4}(1)$ has class $4x$.

We verify that the four theta degrees are equal.  Let
$\pi:C^4\to X$ and $a:C^4\to A$ be the symmetric quotient and sum
of Abel maps, so $a=u\pi$, where $u:X\to A$ is Abel--Jacobi.
Each cycle $\Gamma_j=\pi^*B_j$ is invariant under $N\rtimes\mathfrak S_4$.
Applying $\rho$ on the first factor permutes the $\Gamma_j$ cyclically.
In the K\"unneth decomposition of $H^4(C^4,\mathbf Q)$, $N$-invariance
forces the cyclic characters of all four factors to be equal.  A factor
of degree zero or two has trivial character.  If some factor has this
character, all four do, so the residual action is trivial.  Otherwise
all four factors have degree one, and their common character is $i$ or
$-i$, because $\rho^2$ is the hyperelliptic involution.  Permutation
invariance, with the signs of odd cohomological degree, gives precisely
\[
 \bigwedge^4H^1(C,\mathbf C)_i\ \oplus\
 \bigwedge^4H^1(C,\mathbf C)_{-i}=a^*W_{\mathbf Q(i)}(A)_{\mathbf C}.
\]
In particular there is no residual character $-1$.  The trivial parts
of $[\Gamma_j]$ are all equal, so any difference
$[\Gamma_j]-[\Gamma_k]$ belongs to $a^*W_{\mathbf Q(i)}(A)$.
The finite map $\pi$ has injective rational cohomology pullback; it
follows that $[B_j]-[B_k]\in u^*W_{\mathbf Q(i)}(A)$.
The class $\theta^2$ is orthogonal to $W_{\mathbf Q(i)}(A)$: with
$H^1(A,\mathbf C)=V_i\oplus V_{-i}$ it has bidegree $(2,2)$,
whereas the two Weil lines have bidegrees $(4,0)$ and $(0,4)$.
Their products have a fifth or sixth factor in a four-dimensional
space.  Thus all four numbers $\int_{B_j}u^*\theta^2$ are equal.

For completeness, the diagonal computation evaluates their sum directly.
Let $\Delta_X$ denote the big diagonal divisor on $X$.  The standard
symmetric-product identity gives $u^*\theta=7x-\Delta_X/2$.
One has
\[
 \int_X x^4=1,\qquad
 \int_X x^3\Delta_X=6,\qquad
 \int_X x^2\Delta_X^2=20.
\]
Here is a direct verification of the last two intersection numbers.
On the curve $p+q+r+C$, each of the three fixed points contributes
multiplicity two to the diagonal.  On $p+q+\operatorname{Sym}^2C$,
the restricted diagonal is $\Delta_2+4x$, where $\Delta_2$ is the big
diagonal of $\operatorname{Sym}^2C$.  Its intersections are
$x\Delta_2=2$ and $\Delta_2^2=2(2-2g)=-12$: under
$C^2\to\operatorname{Sym}^2C$, the diagonal pulls back to twice the
ordinary diagonal, whose self-intersection is $2-2g$.
Therefore $(\Delta_2+4x)^2=-12+16+16=20$.  Equation
\eqref{eq:cyclic-four-total} now gives
\[
 \sum_j\deg_\theta u_*B_j
 =16\left(49-7\cdot6+\frac{20}{4}\right)=192.
\]
Each degree is $48$.

We record nonvanishing and the integral norm without dividing an integral
cycle class.  Blow up successively the strict transforms of the three
planes $K_r=\{u_r=v_r=0\}$ in $\mathbf P^4$.  The strict transform
$\widetilde P$ of $P$ remains a plane, since its intersection with each
centre is a Cartier line.  After the first two blowups the ambient space
is a $\mathbf P^2$-bundle over $\mathbf P^1\times\mathbf P^1$ and
$\widetilde P$ is a fibre.  Near that fibre use base coordinates $r_1,r_2$
and fibre coordinates $[t_1:t_2:u_3]$.  The last blowup has the chart
\[
 r_1t_1+r_2t_2=u_3w,
 \qquad \widetilde P=\{r_1=r_2=w=0\}.
\]
Consequently its normal bundle fits into
\[
 0\longrightarrow N_{\widetilde P}\longrightarrow\mathcal O^3
 \xrightarrow{(t_1,t_2,-u_3)}\mathcal O(1)\longrightarrow0,
\]
so its second Chern number is $1$.  Each paired exponent sum is four;
normalization of the pulled-back cyclic cover therefore introduces no
branching along an exceptional divisor.  The strict branch divisors miss
$\widetilde P$.  Above it are four disjoint planes
$Q_0,Q_1,Q_2,Q_3$, each with this normal bundle.
Thus $c=[Q_0]-[Q_2]$ has square $2$.

The primitive order-four cohomology is unchanged by these blowups.
To see the point also for the exponents $2,2$, the exceptional divisor
is a $\mathbf P^1$-bundle.  On a fibre avoiding the other branch divisors,
the faithful-character eigensheaf is extension by zero of a rank-one
local system on $\mathbf C^*$ with monodromy $e^{2\pi i a/4}\ne1$,
where $a=1$ or $2$.  Its compactly supported cohomology is zero.
On a fibre contained in another branch divisor, its eigensheaf is zero
because the inertia is nontrivial.  Proper base change therefore gives
zero faithful-character cohomology for the whole exceptional divisor.
The inverse image of the original centre has the same vanishing, since
its inertia contains the nontrivial subgroup generated by $a$.
The two localization sequences identify the primitive cohomology before
and after each blowup.  This argument treats order-two inertia directly.

The normal cyclic covers have finite quotient singularities and are
oriented rational homology manifolds.  Poincar\'e duality and the
projection formula therefore apply.  If $\varphi:\widetilde W\to W$
is the resulting birational map and $c_0=\varphi_*c$, then
$\varphi^*c_0=c$ and $\int_Wc_0^2=2$.
The map $r:X\to W$ has degree $4^3=64$.
On a resolution of $X\times_W\widetilde W$, take refined pullbacks
of the regular embeddings $Q_j\subset\widetilde W$ and push them
to $X$.  The result is $B_j$ as defined above: over the open plane it
is the ordinary finite pullback; any extra term has image over a
proper subset of $P$, and finite $f$ forces that image to have dimension
less than two, so its pushforward as a surface cycle vanishes.
It follows that
\[
 [B_0]-[B_2]=r^*c_0,\qquad
 \int_X([B_0]-[B_2])^2=64\cdot2=128.
\]
As shown above, this difference is the pullback of a Weil class.
The map $u$ is birational: a general degree-four divisor on a genus-four
curve has exactly one section.  Thus $u_*u^*=1$, and the difference
$[u_*B_0]-[u_*B_2]$ is an integral Weil class with square $128$.

Put $Z_j=u_*B_j$.  The Poincar\'e formula for the Abel image of
$\operatorname{Sym}^2C$ gives $u_*x^2=\theta^2/2$.  This formula
can also be verified without a geometric assumption on that image:
the Abel curve has cohomology class $\theta^3/3!$, and the addition
map $C^2\to A$ pushes its product class to $\theta^2$; division
by the generic symmetric-quotient degree $2$ gives the formula.
The two identities follow by expanding the symplectic basis of
$H^1(C,\mathbf Z)$ under the addition pullback.  Therefore
\[
 \sum_j[Z_j]=8\theta^2,\qquad
 w_j=[Z_j]-2\theta^2\in W_{\mathbf Q(i)}(A).
\]
These are integral classes.  On the primitive space the residual action
has eigenvalues $i,-i$, so $w_{j+2}=-w_j$.  In particular
$[Z_0]-[Z_2]=2w_0$ and $w_0^2=128/4=32$.  All four $w_j$ are nonzero,
since the residual action permutes them.  Choosing two general translates
$\Theta,\Theta'$ of a theta divisor gives the integral algebraic
representative $Z_j-2(\Theta\cdot\Theta')$ of $w_j$, with both
effective numerator degrees equal to $48$.  One irreducible component
of any $Z_j$ detects with degree at most $48$.

Finally we determine the lattice parity.  Let $E$ be the integral
alternating intersection form on $H_1(C,\mathbf Z)$, and put
$h(v,w)=E(v,\rho w)+iE(v,w)$.  It is Hermitian, antilinear in its
first variable, with $h(v,v)=E(v,\rho v)\in\mathbf Z$.
A change of its overall sign does not affect parity.
For a hyperelliptic curve, the reduction of $H_1(C,\mathbf Z)$ modulo
two is the space of even subsets of the ten branch points, modulo the
full set; its symplectic form is the cardinality of intersection modulo
two.  This follows directly by lifting arcs between branch points:
two such arcs sharing one endpoint give cycles with one transverse
intersection, and disjoint pairs give intersection zero.  The action of
$\rho$ fixes $0,\infty$ and exchanges $a_r,-a_r$.
The subsets $\{0,a_1\}$ and $\{0,-a_1\}$ therefore have intersection
of odd cardinality.  If $v$ lifts the first class, then
$E(v,\rho v)$ is odd.  The Gaussian lattice is consequently odd,
as claimed.
\end{proof}

\begin{remark}\label{rem:cyclic-four-eightfold-seed}
The product of two such degree-$48$ detecting surfaces has nonzero
Weil projection on the product eightfold: on the two complex determinant
lines its projection is the tensor product of the corresponding nonzero
components.  Its degree for the product polarization is
$\binom42 48^2=13824$.  The product has odd Gaussian lattice
$I_{4,4}$.  The Gaussian isogeny with matrix
$\left(\begin{smallmatrix}1&1\\1&-1\end{smallmatrix}\right)^{\oplus4}$
from the odd lattice to $U^4$ pulls back the even principal polarization
to twice the odd one.  Pushforward therefore gives a degree bound
$2^4\cdot13824=221184$ at a principal $U^4$ point.  These points are
isogenous to products.  They lie in proper special subvarieties and
cannot be relative-Weyl points with full symmetric Weyl group.
Neither this construction nor its Hecke translates proves a uniform
bound on the full eightfold family.
\end{remark}

\subsection{The Schoen surface and its principal quotient}

The published geometric input is the deformation of
$\Delta_J\cup(C\times C)\subset J\times J$, for a genus-two Jacobian $J$, to an embedded Schoen surface in a four-dimensional Eisenstein Weil family. We use the form of Schoen's original smoothing theorem \cite[Theorem~1.1(i)--(iii), p.~585]{Schoen07}, together with the ambient-family description of van Geemen \cite[Introduction, p.~1, and \S6.2, pp.~14--15]{vG22}, including the ambient family and its polarization. The embedded flat degeneration is also described in \cite[\S3, pp.~5--6]{Beauville}, and the four-dimensional smooth surface moduli in \cite[Theorem~1]{RRS}.

For clarity, the dimension of the surface moduli really gives dominance onto the Weil component. A smooth Schoen surface $S\subset A$ is Lagrangian for a nondegenerate holomorphic two-form and has $q(S)=4$ \cite[Introduction]{RRS}. Thus $N_{S/A}\simeq\Omega_S^1$ and $h^0(N_{S/A})=4$. Translation supplies four independent infinitesimal embedded deformations: the infinitesimal stabilizer is zero because a smooth surface of general type has no vector fields. Therefore a fixed abelian variety admits no positive-dimensional family of these embedded surfaces modulo translation. The map from the four-dimensional surface moduli to Albanese moduli has zero-dimensional fibres. Its image in the four-dimensional Weil component is dominant. This supplies an analytic open set of the component carrying the required effective surface class.

We compute its degree and Weil part at the degeneration. Let $\vartheta$ be the principal class on $J$, and on $J^2$ set
\[
 x=p_1^*\vartheta,\qquad y=p_2^*\vartheta,\qquad
 z=m^*\vartheta-x-y.
\]
The polarization $a x+b z+c y$ corresponds to the matrix
$\left(\begin{smallmatrix}a&b\\b&c\end{smallmatrix}\right)$ on the two factors. The determinant formula for its Euler characteristic, or expansion in a symplectic basis, gives
\begin{equation}\label{eq:intersection-polynomial}
 \int_{J^2}(a x+b z+c y)^4=24(ac-b^2)^2.
\end{equation}
In particular
\[
 \int x^2y^2=4,\qquad \int xyz^2=-4,\qquad
 \int z^4=24,
\]
and the other degree-four monomial intersections vanish. The central surface has class
\[
 s=[\Delta_J]+[C\times C]
   =\tfrac12(x+y-z)^2+xy.
\]
Take $\omega=(-1+\sqrt{-3})/2$ and the endomorphism and polarization
\[
 R=\begin{pmatrix}0&-1\\1&-1\end{pmatrix},\qquad
 H=2x-z+2y.
\]
The matrix of $H$ is $M=\left(\begin{smallmatrix}2&-1\\-1&2\end{smallmatrix}\right)$ and $R^tMR=M$. Its polarization type is $(1,1,3,3)$ and $H^4=216$. Equation~\eqref{eq:intersection-polynomial} gives
\[
 \int H^2s=36,\qquad \int s^2=8.
\]
For example, the two surface components contribute respectively $8$ and $28$ to the first number. The diagonal has self-intersection zero, the product has self-intersection $4$, and their mutual intersection is $2$, giving the second number.

Put $\ell_+=\omega^2x+\omega z+y$ and
$\ell_-=\omega x+\omega^2z+y$. Their squares span the two extreme $K$-eigenlines. Direct expansion gives
\[
 w_0=s-\tfrac16H^2
   =\tfrac16\bigl(\omega\ell_+^2+\omega^2\ell_-^2\bigr)\in W,
 \qquad \int w_0^2=2.
\]
The same flat class and norm hold on the deformed surface.

The polarization kernel is $\{(u,-u):u\in J[3]\}$, and $R$ acts trivially on it. A maximal isotropic subgroup of order $9$ therefore gives an Eisenstein-linear quotient
\[
 q:A_H\longrightarrow A,\qquad \deg q=9,\qquad q^*\theta=H,
\]
with principal polarization on $A$. The quotient can be chosen over a finite cover of the period family. To identify its principal lattice, specialize $J$ to $E_\omega^2$. The explicit map
\[
 F=\begin{pmatrix}1&\omega\\1&\omega^2\end{pmatrix}
       \otimes I_{E_\omega^2}
\]
has degree $9$, satisfies $F^*\theta_{E_\omega^4}=H$, and intertwines $R$ with $\operatorname{diag}(\omega,\omega^2)$. Its target is the balanced product $E_\omega^4$ with product polarization. Transporting this finite-index lattice along the period domain identifies the whole principal component.

This component covers every principal Eisenstein lattice of signature $(2,2)$. Indeed the maximal order is a principal ideal domain, and the associated integral Hermitian lattice is unimodular with determinant $1$. It is rationally split by the classification in \cite[Proposition~4.1 and Corollary~4.2]{Deligne82}. Choose a primitive integral isotropic vector $v$. Unimodularity gives $w$ with $h(v,w)=1$. If $h(w,w)=b\in\Z$, replacing $w$ by $w+b\omega v$ makes it isotropic because $\operatorname{Tr}(\omega)=-1$. Thus $v,w$ span a unimodular hyperbolic plane, which splits off integrally. Repeat on its signature-$(1,1)$ complement. The lattice is $U\oplus U$. Finally, in $I_{1,1}$ the vectors $(1,1)$ and $(1+\omega,\omega)$ have Gram matrix $U$ and determinant $-1$, so $U\oplus U\simeq I_{2,2}$.

Push the entire effective surface cycle forward, including its multiplicities, and write $Z=q_*S$. Then
\begin{equation}\label{eq:degree36-class}
 \deg_\theta Z=36,\qquad
 [Z]=\tfrac32\theta^2+w,\qquad
 w\in\mathcal W,\qquad \int_Aw^2=18.
\end{equation}
The degree follows by projection formula. For the norm use $q^*q_*s=9s$ in cohomology, so $(q_*s)^2=9s^2=72$, and subtract $(3\theta^2/2)^2=54$. The integrality of $w$ follows from that of $\theta^2/2$, which is immediate by expanding the square of an integral alternating two-form. Proper cycle parameter spaces extend an effective cycle with these same flat relations from the nonempty open set to every fibre: the components on which these relations hold are open and closed, and at least one dominates the base. Their proper image is all the base. This proves the degree-$36$ detector bound.

\subsection{The saturated Eisenstein lattice}

It remains to turn the effective cycle into the integral coefficient bound. Compute at the balanced product $E_\omega^4$. Choose integral one-form bases $x_j,y_j$, adapted to the maximal-order action, and write
\[
 \prod_{j=1}^4(x_j+\omega y_j)=u+\omega v.
\]
The all-$x$ coefficient of $(u,v)$ is $(1,0)$, and any coefficient with exactly one $y$ is $(0,1)$. Consequently a rational combination of $u,v$ is integral only when both coefficients are integers. They are a basis of the saturated Weil lattice. The identities
\[
 (u+\omega v)^2=0,\qquad
 \int(u+\omega v)(u+\omega^2v)=9
\]
give
\begin{equation}\label{eq:eisenstein-gram}
 \operatorname{Gram}(\mathcal W)=
 \begin{pmatrix}6&3\\3&6\end{pmatrix}.
\end{equation}
These are integral local systems, so the result holds throughout the family. If $\rho$ is the automorphism induced by $\omega$, then
\[
 \rho^*u=-v,\qquad \rho^*v=u-v.
\]
The norm-$18$ vectors in this lattice have coordinates
$\pm(1,1),\pm(2,-1),\pm(-1,2)$: this follows by solving
$m^2+mn+n^2=3$. The three unit translates of any such vector have differences giving, up to signs and a lattice isometry, the three directions $3u,3v,3(u-v)$. Applied to the Weil part in \eqref{eq:degree36-class}, the differences of the three degree-$36$ cycles
$Z,\rho^*Z,\rho^{2*}Z$ therefore generate exactly $3\mathcal W$.

Write $\delta t=mu+nv$. It is a sum of
\[
 N=\max\{|m|,|n|,|m+n|\}
\]
signed steps in the directions $u,v,u-v$: for equal signs use $u$ and $v$, and for opposite signs first use $u-v$ and then the remaining direction. Put $Q=m^2+mn+n^2=\delta^2T/6$. The three identities
\[
 4Q-3m^2=(m+2n)^2,\quad
 4Q-3n^2=(2m+n)^2,\quad
 4Q-3(m+n)^2=(m-n)^2
\]
show that $N\leq\delta\sqrt{2T}/3$. Each step, multiplied by $3$, is represented by a difference of two degree-$36$ effective cycles. Hence $3\delta t$ has effective numerator degrees at most $36N$. This proves the first inequality in \eqref{eq:eisenstein-cost}; the second follows from $12\sqrt2<17$ and $\lceil\sqrt T\rceil\geq1$.

\begin{example}\label{ex:eisenstein-integral-class}
For the integral class $t=u$ in \eqref{eq:eisenstein-gram}, one has $T=6$ and may take $\delta=1$. A single unit-orbit difference represents $3u$, with both effective degrees $36$. The floor in \eqref{eq:eisenstein-cost} is then $\lfloor\sqrt{12}/3\rfloor=1$, exactly reproducing this presentation.
\end{example}
\section{Principal hyperbolic quotients over every imaginary quadratic field}
\label{sec:principal-quotients}
This section reduces all split integral polarization types in a fixed dimension to one principal hyperbolic family, with a degree cost linear in the polarization volume.

Let $K$ be an imaginary quadratic field, let $\mathcal O=\mathcal O_K$, and define $t_K\in\{1,2\}$ by
\[
 \operatorname{Tr}_{K/\Q}(\mathcal O)=t_K\Z.
\]
Thus $t_K=1$ when the discriminant is odd and $t_K=2$ when it is even. For a projective $\mathcal O$-module $Q$, write
\[
 Q^\vee=\operatorname{Hom}_{\bar{\mathcal O}}(Q,\mathcal O).
\]
Its elements are additive functions satisfying $p(aq)=\bar a p(q)$. The hyperbolic Hermitian module $H(Q)$ is $Q^\vee\oplus Q$, with pairing
\[
 ((p,q),(p',q'))\longmapsto p(q')+\overline{p'(q)}.
\]
We use Hermitian forms linear in the first variable.

\begin{theorem}\label{thm:principal-quotient}
Let $(A,L)$ be a split Weil $2n$-fold over $K$, with compatible polarization and integral action by $\mathcal O_f=\Z+f\mathcal O$. Put $\Delta=\chi(L)$. There is a $K$-linear isogeny to a principally polarized split Weil variety,
\[
 q:A\longrightarrow(P,\Theta),\qquad
 q^*\Theta=t_K fL,\qquad
 \deg q=(t_K f)^{2n}\Delta,
\]
where equality of polarizations means equality of their first Chern classes. The integral Hermitian lattice of $P$ is $H(Q)$ for a projective rank-$n$ module $Q$.

At most $h_K$ isomorphism types of such hyperbolic modules occur, where $h_K$ is the ideal class number. If $C_I$ bounds detectors throughout the principal family associated with $H(\mathcal O^{n-1}\oplus I)$, then
\[
 \delta_W(A,L)\le t_K^n f^n\Delta\,
       \max_{[I]\in\operatorname{Cl}(\mathcal O)} C_I.
\]
The same conclusion on CM fibres requires the bounds only on the CM points of those principal families. If $\mathcal O$ has class number one, only the lattice $U^n$ is needed.

More generally, put
\[
 m_K=\max_{[I]\in\operatorname{Cl}(\mathcal O)}
       \min_{\substack{J\subset\mathcal O\\{[J]=[I]}}}
                    \operatorname{Nm}(J).
\]
If $C_U$ bounds detectors on the single principal family with lattice $U^n$, then
\[
 \delta_W(A,L)\le(t_Kfm_K)^n\Delta C_U,
 \qquad m_K\le\frac2\pi\sqrt{|\operatorname{disc}K|}.
\]
Again the CM version holds with a supplier only on CM points.
\end{theorem}

\begin{proof}
Write $A=V_{\R}/\Lambda$, with Riemann form $E$, and choose $\tau$ such that $\mathcal O=\Z[\tau]$. Set $S=\mathcal O\Lambda=\Lambda+\tau\Lambda$ and $b=[S:\Lambda]$. The form $E'=fE$ is integral on $S$. Indeed the cross terms use the integral endomorphism $f\tau$, and Rosati compatibility gives
\[
 E(\tau x,\tau y)=\operatorname{Nm}(\tau)E(x,y).
\]
The quotient with lattice $S$ therefore has polarization volume
\[
 \Delta'=\frac{f^{2n}\Delta}{b}.
\]
On $S$ use the integral Hermitian form
\[
 H(x,y)=E'(\tau x,y)-\bar\tau E'(x,y),
 \qquad
 E'(x,y)=\operatorname{Tr}_{K/\Q}
          \frac{H(x,y)}{\tau-\bar\tau}.
\]
The first formula is linear in $x$ and conjugate-linear in $y$. The codifferent of $\mathcal O$ is $(\tau-\bar\tau)^{-1}\mathcal O$. Consequently a perfect integral Hermitian pairing gives a perfect alternating $\Z$-pairing by the second formula. This assertion holds for projective modules and requires no choice of a free basis.

Choose a maximal isotropic $K$-subspace $W$ of dimension $n$. The module $P_0=S\cap W$ is saturated: if a nonzero $a\in\mathcal O$ has $ax\in W$, then $x\in W$. The quotient $S/P_0$ is torsion-free over a Dedekind domain, hence projective. Choose a splitting
\[
 S=P_0\oplus Q.
\]
Nondegeneracy identifies $W$ with the conjugate dual of $Q_K$ by $w\mapsto H(w,\cdot)|_Q$. Let $P'$ be the inverse image of $Q^\vee$. Integrality gives $P_0\subset P'$. On $\Gamma=P'\oplus Q$, the form becomes
\[
 H(p+q,p'+q')=p(q')+\overline{p'(q)}+C(q,q'),
 \qquad C=H|_{Q\times Q}.
\]
The off-diagonal evaluation pairing is perfect, so this Hermitian form is unimodular. Its alternating form is therefore principal. Comparing the polarization volumes on $S$ and $\Gamma$ gives
\[
 [\Gamma:S]=\Delta'.
\]

Choose $\eta\in\mathcal O$ with $\eta+\bar\eta=t_K$, and let $c:Q\to Q^\vee$ be the $\mathcal O$-linear map $c(q)(q')=C(q,q')$. Enlarge the lattice once more, to
\[
 \Gamma_t=\frac1{t_K}P'\oplus Q,
\]
and give it the form $t_KH$. The map
\[
 Q^\vee\oplus Q\longrightarrow\Gamma_t,
 \qquad (p,q)\longmapsto\frac p{t_K}+q-
                                  \frac{\eta c(q)}{t_K}
\]
is an integral module isomorphism. Its off-diagonal pairing is evaluation; its pairing on the second summand is
$t_KC-\eta C-\bar\eta C=0$. It is thus an isometry from $H(Q)$.
In particular this is a principal hyperbolic lattice and
\[
 [\Gamma_t:S]=t_K^{2n}\Delta'.
\]
All these lattices lie in the same real vector space with the same complex structure. The Riemann form has only been multiplied by a positive integer, so positivity is preserved. The inclusion $\Lambda\subset\Gamma_t$ defines a holomorphic $K$-linear isogeny with
\[
 q^*\Theta=t_K fL,
 \qquad \deg q=b\,t_K^{2n}\Delta'=(t_K f)^{2n}\Delta.
\]

For an effective detecting $n$-cycle $Z$ on $P$, finite flat pullback and the projection formula give
\[
 \deg_L(q^*Z)=\frac{\deg q}{(t_K f)^n}\deg_\Theta Z
             =t_K^n f^n\Delta\deg_\Theta Z.
\]
The pullback has nonzero Weil projection because a $K$-linear isogeny identifies the rational Weil spaces. Some irreducible component is a detector, and its degree is at most the total degree.

Finally every projective rank-$n$ module over $\mathcal O$ has Steinitz form $\mathcal O^{n-1}\oplus I$. Isomorphic modules give isometric hyperbolic modules, so at most $h_K$ types are required. For each type the compatible period domain is connected; a neat arithmetic quotient gives its principal family. Passing to a finite level cover does not affect the bound.

Choose the representative $I$ integral, of norm $m$. In coordinates,
\[
 Q=\mathcal O^{n-1}\oplus I,
 \qquad Q^\vee=\mathcal O^{n-1}\oplus\bar I^{-1},
 \qquad m\bar I^{-1}=I.
\]
Thus $(p,q)\mapsto(mp,q)$ embeds $H(Q)$ into the standard lattice $\mathcal O^{2n}$ and multiplies its Hermitian form by $m$. Transporting the complex structure through this map gives an isogeny $P_I\to P_U$ with pullback polarization $m\Theta_I$. Since both polarizations are principal, its degree is $m^{2n}$. Pullback of a middle-dimensional detector therefore costs $m^n$. Taking the smallest norm in each ideal class proves the bound involving $m_K$. All isogenies here preserve CM.

For the numerical estimate, an integral ideal $J$ embedded in $\C$ has covolume $\operatorname{Nm}(J)\sqrt{|\operatorname{disc}K|}/2$. The planar convex-body theorem applied to a disk gives $0\ne\alpha\in J$ with
\[
 |\operatorname{Nm}(\alpha)|\le
 \frac2\pi\sqrt{|\operatorname{disc}K|}\operatorname{Nm}(J).
\]
Indeed a centrally symmetric convex set of area exceeding four times the covolume contains a nonzero lattice point: two points of its half-sized copy have the same image in a lattice fundamental domain, and their difference lies in the original set. Apply this to disks with radii decreasing to the stated bound. The ideal $(\alpha)J^{-1}$ is integral, lies in the inverse class of $J$, and has norm at most $2\sqrt{|\operatorname{disc}K|}/\pi$. This proves the estimate for every class.
\end{proof}

\begin{lemma}[Symmetrization on a projective Hermitian module]
\label{principal-lem:even-symmetrization}
Let $\mathcal O$ be the ring of integers of an imaginary quadratic
field with $\operatorname{Tr}(\mathcal O)=2\mathbb Z$.
Let $Q$ be a projective $\mathcal O$-module of finite rank, and let
$C:Q\times Q\to\mathcal O$ be an integral Hermitian form, linear
in the first variable, such that $C(q,q)\in2\mathbb Z$ for all $q\in Q$.
Then there is an $\mathcal O$-linear map $R:Q\to Q^\vee$ with
\[
 C(q,q')=R(q)(q')+\overline{R(q')(q)}.
\]
\end{lemma}

\begin{proof}
Let $\mathcal S=\operatorname{Hom}_{\mathcal O}(Q,Q^\vee)$,
regarded as a finite free $\mathbb Z$-module, with involution
$R^*(q)(q')=\overline{R(q')(q)}$.
The cokernel of
\[
 \mathcal S\longrightarrow\mathcal S^{*=1},
 \qquad R\longmapsto R+R^*,
\]
is killed by two.
It suffices to show that the class of $C$ vanishes after tensoring
with $\mathbb Z_2$.
The projective module $Q\otimes\mathbb Z_2$ is free over
the semilocal ring $\mathcal O\otimes\mathbb Z_2$.
By continuity, the norm-even hypothesis holds on this completed
module, so a matrix of $C$ has diagonal entries in $2\mathbb Z_2$.
Choose half of each diagonal entry and one triangular half of the
off-diagonal entries to obtain a matrix $R_2$ with
$R_2+R_2^*=C$.
Flat base change commutes with the kernel defining
$\mathcal S^{*=1}$ and with the cokernel of symmetrization.
Thus the cokernel class of $C$ becomes zero over $\mathbb Z_2$.
Since this cokernel is killed by two, the class was already zero.
\end{proof}

\begin{corollary}[Even saturated lattices]
\label{principal-cor:even-saturation}
In Theorem~\ref{thm:principal-quotient}, suppose $t_K=2$ and
the Hermitian form associated with $fE$ on the saturated
maximal-order lattice $S=\mathcal O\Lambda$ has even norms.
Then the factor $t_K$ can be omitted:
there is a principal hyperbolic quotient with
\[
 q^*\Theta=fL,\qquad \deg q=f^{2n}\Delta.
\]
The detector costs are consequently
$f^n\Delta\max_I C_I$ and $(fm_K)^n\Delta C_U$, respectively.
Freeness of $S$ is not required.
\end{corollary}

\begin{proof}
Use $S=P_0\oplus Q$, $P'$ and $\Gamma=P'\oplus Q$ from the proof
of Theorem~\ref{thm:principal-quotient}.
The restriction $C=H|_{Q\times Q}$ has even norms.
By Lemma~\ref{principal-lem:even-symmetrization}, write $C=R+R^*$.
Via $P'\cong Q^\vee$, replace the second summand by the graph
$\{q-R(q):q\in Q\}$.
Its restriction of $H$ is zero, and the cross pairing is unchanged.
Thus $(\Gamma,H)$ is already $H(Q)$; the enlargement to
$t_K^{-1}P'\oplus Q$ and the scaling by $t_K$ are unnecessary.
The degree and detector formulas follow from the same volume and
projection-formula calculation.
\end{proof}

\section{Sixfold cycles after Schoen and Koike, and bounds for every fourfold determinant}
\label{sec:sixfolds}
This section keeps track of the degrees in Schoen's cyclic-cover construction and improves the passage between integral polarization lattices.

Write $U=\left(\begin{smallmatrix}0&1\\1&0\end{smallmatrix}\right)$ for a Hermitian hyperbolic plane. Degrees are unnormalized intersection numbers. All polarizations below are compatible with the indicated quadratic action.

\begin{theorem}\label{thm:sixfold-suppliers}
For full maximal-order action, the following detector bounds hold on the indicated principal split sixfold families:
\[
\begin{array}{lll}
K&\text{Hermitian lattice}&\delta_W\text{ at most}\\
\hline
\Q(i)&U^3&746496\\
\Q(i)&I_{3,3}&5971968\\
\Q(\omega)&U^3\cong I_{3,3}&30375.
\end{array}
\]
For an arbitrary full maximal-order split sixfold of polarization product $\Delta$, one has
\[
 \delta_W\le5971968\Delta\quad(K=\Q(i)),
 \qquad \delta_W\le30375\Delta\quad(K=\Q(\omega)).
\]
For an even Gaussian Hermitian lattice the first bound improves to $746496\Delta$.
\end{theorem}

\begin{theorem}\label{thm:all-determinant-fourfolds}
Let $(A,L)$ be a Weil fourfold over $K$, of any determinant, and put $\Delta=\chi(L)$. Suppose the order $\mathcal O_f=\Z+f\mathcal O_K$ acts integrally. Then
\[
\begin{aligned}
 \delta_W(A,L)&\le1990656\,f^6\Delta^2
 &&\text{if }K=\Q(i),\\
 \delta_W(A,L)&\le10125\,f^6\Delta^2
 &&\text{if }K=\Q(\omega).
\end{aligned}
\]
Each right-hand side may be divided by $d$, where
$d=[\mathcal O_KH_1(A,\Z):H_1(A,\Z)]$.
\end{theorem}

The geometric inputs are published: Schoen constructs the cycles \cite[Theorem~2.0 and \S3]{Schoen88}, and Schoen and Koike prove the required dominance \cite[Lemma~13, pp.~334--335]{Schoen98}, \cite[Theorem~2.1, p.~568]{Koike}. The estimates below follow from their constructions, integral polarization calculations and closedness of bounded detector loci. They do not use the recent CM theorems or equivariant semiregularity.

\subsection{The cyclic sources and their principal polarizations}

Let $C\to X$ be a connected cyclic \'etale cover of degree $m=3$ or $4$, with $g(X)=4$ and generator $\sigma$. Then $g(C)=3m+1$. Choose a topological symplectic basis on $X$ for which all monodromy lies on the first handle. The other three handles lift in $m$ disjoint copies. Their cycles satisfy
\[
 E(a_{j,r},b_{k,s})=\delta_{jk}\delta_{rs},
 \qquad E(a_{j,r},a_{k,s})=E(b_{j,r},b_{k,s})=0,
\]
and $\sigma$ shifts the sheet index. The cover of the first punctured handle is a genus-one surface with $m$ boundary circles. After the other handles are glued, its boundary classes vanish and it contributes an invariant rank-two symplectic summand. This gives an integral orthogonal decomposition into that summand and three pairs of regular representations.

For $m=4$, the saturated primitive lattice on each handle is generated by
\[
 a_0-a_2,\quad a_1-a_3,\quad b_0-b_2,\quad b_1-b_3.
\]
The two nonzero alternating pairings are both two. Thus on the primitive sixfold $B=\ker(1+\sigma^2)^0$ the restriction of the Jacobian polarization is $2\Xi$, where $\Xi$ is principal. The Gaussian Hermitian matrix for $\Xi$ on a handle has off-diagonal entries $i,-i$; changing a basis vector by a unit gives $U$. The principal lattice is $U^3$.

For $m=3$, the norm Prym $P$ has, on each handle, augmentation generators $A=a_0-a_1$ and $B=b_0-b_1$ over $\Z[\omega]$. Using the second-variable convention of
Theorem~\ref{thm:explicit-stabilization}, the form
$h(x,y)=E(x,\omega y)-\bar\omega E(x,y)$ has
\[
 E(A,B)=2,\quad E(A,\omega B)=-1,
 \qquad h(A,B)=1+2\omega=\sqrt{-3}.
\]
Its matrix is $\left(\begin{smallmatrix}0&\sqrt{-3}\\-\sqrt{-3}&0\end{smallmatrix}\right)$ and its alternating type is $(1,3)$. Enlarge the lattice on each handle from $\langle A,B\rangle$ to $\langle A,B/\sqrt{-3}\rangle$. This has index three and gives matrix $U$. On three handles it defines a degree-$27$ quotient $P\to P'$ with a principal polarization $\Xi'$ satisfying
\[
 \varphi^*\Xi'=\Theta_C|_P.
\]
The positive Riemann form on the real vector space has not changed. A finite level cover chooses the kernels in families and identifies the period domains. The Eisenstein basis with columns $(1,-\omega)$ and $(1,\bar\omega)$ takes $U$ to $\operatorname{diag}(1,-1)$ and has unit determinant. Thus the principal lattice is also $I_{3,3}$.

Koike's Theorem~2.1 states that the quartic Prym map is dominant onto the nine-dimensional unitary domain; its codifferential is multiplication of inverse primitive eigendifferentials. Schoen's Lemma~13 proves the cubic dominance by exhibiting nine independent products. These statements apply to the principalized sources just described, since a fixed isogeny change of lattice does not change the period domain.

\subsection{An effective detecting component and its degree}

Put
\[
 N=\ker\bigl((\Z/m)^6\xrightarrow{\sum}\Z/m\bigr),
 \qquad G_0=N\rtimes S_6,
 \qquad W_0=C^6/G_0.
\]
The finite maps
\[
 C^6\xrightarrow{r}C^{(6)}\xrightarrow{p}W_0
 \xrightarrow{\beta}X^{(6)}
\]
have generic degrees $6!,m^5,m$. Schoen's local equation \cite[Lemma~2.2, p.~13]{Schoen88} makes $\beta$ \'etale in the unramified case. The inverse image of $|K_X|\cong\mathbf P^3$ is a disjoint union of $m$ copies of $\mathbf P^3$. Choose one, $Q$.

The nonvanishing is especially simple here. Embed $X^{(6)}$ in $X^{(7)}$ by adding a point. Since $X^{(7)}\to\operatorname{Pic}^7(X)$ is a projective bundle, the canonical $\mathbf P^3$, after this embedding, is a fibre. Its normal bundle there is trivial of rank four; the normal line of $X^{(6)}$ restricts to $\mathcal O_{\mathbf P^3}(1)$. The normal sequence gives
\[
 c(N_{\mathbf P^3/X^{(6)}})=(1+H)^{-1},
 \qquad c_3(N_{\mathbf P^3/X^{(6)}})=-H^3\ne0.
\]
The same holds for $Q\subset W_0$. For a primitive character $\chi$ of the residual deck group, the cycle
$\sum_t\chi(-t)\,\sigma^tQ$ restricts on $Q$ to this nonzero top Chern class, since the other sheets are disjoint. Thus the primitive-character projection is nonzero. This is the unramified calculation of \cite[Lemma~2.6, p.~16]{Schoen88}.

Let $Z$ be an irreducible component of the inverse image of $Q$ in $C^6$ and let $W=r(Z)$ with reduced structure. A general canonical divisor of $X$ has six distinct zeros, so the relevant action is free at the generic point above it. If $H_Z\subset G_0$ stabilizes $Z$, then
\[
 \deg(Z/Q)=|H_Z|,
 \quad \deg(Z/W)=|H_Z\cap S_6|,
 \quad e:=\deg(W/Q)\le m^5.
\]
The primitive determinant plane $U_C\subset H^1(C)^{\otimes6}$ is $G_0$-invariant: the sum-zero sheet translations act trivially, and the alternating determinant and the cohomological permutation signs cancel. Its primitive-character components are those just detected on $W_0$. For $q:C^6\to W_0$, the generic cycle identity is
\[
 \sum_{g\in G_0}g_*[Z]=\deg(Z/Q)q^*[Q].
\]
Pair with an element of $U_C$ for which the right side is nonzero. All summands on the left have the same pairing, so $Z$ itself detects. This produces an effective integral component, without interpreting cyclotomic coefficients as effective cycles.

Let $\xi$ be the incidence class on $C^{(6)}$, and let $\alpha:C^{(6)}\to J(C)$ be an Abel map. Pullback of an incidence hyperplane of $|K_X|$ has class $m\xi$. Consequently
\begin{equation}\label{eq:cyclic-incidence}
 m^3\int_W\xi^3=e,\qquad \int_W\xi^3\le m^2.
\end{equation}
Write $\vartheta=\alpha^*\Theta_C$ and let $\mathcal D$ be the reduced big diagonal of $C^{(6)}$. Then
\begin{equation}\label{eq:symmetric-theta}
 \vartheta=(g(C)+5)\xi-\tfrac12[\mathcal D].
\end{equation}
To verify this, pull back to $C^6$. The incidence class becomes the sum of the point classes. The pair diagonal is the sum of two point classes minus the degree-one cross pairing; the Abel sum pulls theta back to $g(C)$ times the sum of point classes plus these cross pairings. Sum over the fifteen pairs and use $r^*\mathcal D=2\sum_{i<j}\Delta_{ij}$. Rational pullback is injective for the finite quotient.

The class $\xi$ is ample because its finite pullback is ample, and $\vartheta$ is nef. Since $W$ is not contained in $\mathcal D$, the restricted diagonal is effective. With $b=g(C)+5$,
\[
 (b\xi)^3-\vartheta^3
 =\tfrac12[\mathcal D](b^2\xi^2+b\xi\vartheta+\vartheta^2)
\]
has nonnegative intersection with $W$, computed on a resolution if necessary. This uses effectivity of the restricted diagonal, not its nefness. Hence
\begin{equation}\label{eq:cyclic-theta-degree}
 \int_W\vartheta^3\le b^3m^2.
\end{equation}

For $m=4$ use $\phi=(1-\sigma^2)\alpha$ with image in $B$. Its pullback of the Weil plane to $C^6$ is $U_C$, up to nonzero scalars. Indeed expand the exterior product of six primitive one-forms under the Abel sum. Two forms of the same primitive character on one factor multiply to zero, since their product has nontrivial character whereas $H^2(C)$ has trivial character. The surviving alternating tensor is nonzero; $1-\sigma^2$ is invertible on the primitive eigenspaces. Projection formula shows that $\phi(W)$ has dimension three and detects.

The theorem of the square and $\Theta_C|_B=2\Xi$ give
\[
 (1-\sigma^2)^*\Xi
 =2\Theta_C-\tfrac12(1+\sigma^2)^*\Theta_C
 \le2\Theta_C.
\]
Thus its effective pushforward has degree at most
\[
 2^3\cdot18^3\cdot4^2=746496.
\]
For $m=3$, $\alpha(W)$ lies in a translate of the norm Prym. The same argument using the invertible endomorphism $1-\sigma$ proves detection, and its restricted-theta degree is at most
\[
 15^3\cdot3^2=30375.
\]
The principal quotient has pullback polarization equal to that restricted theta, so adds no degree factor. Selecting a reduced image can only decrease either bound.

The bounds are uniform on the dense cyclic-source loci. Bounded-detector closedness extends them to the whole principal components. For the odd Gaussian bound, the three-block matrix with blocks $M=\left(\begin{smallmatrix}1&1\\1&-1\end{smallmatrix}\right)$ satisfies $M^*I_{1,1}M=2U$. It gives an isogeny from the even principal family to the odd one with polarization multiplier two. Threefold degrees increase by at most eight, giving $5971968$.

These are all Gaussian principal parity types. Split off a primitive isotropic vector and a vector pairing to one in a unimodular lattice. The partner norm reduces modulo two. An even norm gives $U$ after subtracting an integral multiple of the isotropic vector; an odd norm gives $I_{1,1}$. Repeat. A mixture containing an odd block becomes $I_{3,3}$: in $\langle1\rangle\oplus U$, with basis $e,a,b$, the vectors $e+a,e-b,e+a-b$ are a unimodular orthogonal basis of norms $1,1,-1$. This converts each remaining hyperbolic block. The rational isotropic spaces exist because the Hermitian space is split.

Theorem~\ref{thm:principal-quotient}, with $n=3$, transfers the principal bounds to every split polarization: $t_K^3\Delta$ is $8\Delta$ for the Gaussian ring and $\Delta$ for the Eisenstein ring. In the even Gaussian case the final paragraph of Section~\ref{sec:principal-quotients} gives the factor $\Delta$. This finishes Theorem~\ref{thm:sixfold-suppliers}.

\subsection{All determinants and nonmaximal orders}

\begin{proof}[Proof of Theorem~\ref{thm:all-determinant-fourfolds}]
First suppose the full maximal order acts. The balanced surface of Theorem~\ref{thm:explicit-stabilization} makes $A\times E_K^2$ a split sixfold of polarization product $\Delta^2$. In both fields $\kappa_K(\Delta)=1$, so the sixfold bound and the factor $1/3$ give
\[
 \delta_W(A,L)\le\frac{5971968\Delta^2}{3}
 =1990656\Delta^2
\]
in the Gaussian case, and $30375\Delta^2/3=10125\Delta^2$ in the Eisenstein case. The Gaussian product contains a norm-one vector from the surface, so no even-lattice improvement is used.

For conductor $f$, put $\Lambda'=\mathcal O_K\Lambda=\Lambda+\tau\Lambda$ and $d=[\Lambda':\Lambda]$. The form $fE$ is integral on $\Lambda'$: cross terms use the integral endomorphism $f\tau$, while
$E(\tau x,\tau y)=\Nm(\tau)E(x,y)$ follows from Rosati compatibility. The quotient isogeny $q:A\to A'$ has a compatible polarization with
\[
 q^*\theta'=f\theta,
 \qquad \Delta'=f^4\Delta/d.
\]
Apply the maximal-order bound $C(\Delta')^2$, where $C$ is the appropriate constant. Finite flat pullback preserves integral effectivity and a nonzero Weil projection, and multiplies the surface degree by $d/f^2$. Selecting a detecting irreducible component gives
\[
 \delta_W(A,L)\le\frac d{f^2}C\left(\frac{f^4\Delta}{d}\right)^2
 =C\frac{f^6}{d}\Delta^2.
\]
\end{proof}

The principal quotient improves transfer once a principal detector is known; it supplies no principal detector in a new dimension. Cyclic \'etale covers of genus-five curves have twelve moduli, whereas the signature-$(4,4)$ unitary domain has dimension sixteen. The genus-four dominance used here therefore does not extend to eightfolds by the same unramified cyclic construction.
\section{An explicit sixfold bound over every imaginary quadratic field, after Markman and Perry}
\label{sec:effective-sixfolds}
This section makes the ideal-curve and Fourier-transform construction quantitative over every imaginary quadratic field. The geometric design is that of Markman's detailed preprint \cite[\S\S8--9]{MarkmanSecant}; its qualitative consequence is also stated in his ICM account \cite[Theorem~1.2]{MarkmanICM}. The deformation input is Perry's Theorem~1.2 \cite{Perry}, an unrefereed preprint: a perfect complex with no negative self-Ext, a coherent finite translation--tensor action, injective semiregularity on invariant obstructions, and an algebraic rational $B$-field whose entire twisted character stays Hodge deforms locally as a twisted perfect complex. We use that theorem with its hypotheses as stated. The estimates below do not use either recent CM Hodge theorem.

\begin{theorem}[With Perry's semiregularity theorem]\label{thm:all-field-sixfolds}
Let $d$ be a positive squarefree integer, let $K=\Q(\sqrt{-d})$, and let $(A,L)$ be a split Weil sixfold with compatible polarization, integral $\mathcal O_f$-action and $\Delta=\chi(L)$. With $m_K$ as in Theorem~\ref{thm:principal-quotient},
\[
 \delta_W(A,L)\le8m_K^3 C_{4d} f^3\Delta
             <10^{179}d^{53}f^3\Delta.
\]
The positive integer $C_D$ is defined explicitly in \eqref{eq:effective-degree} below. When $K$ has class number one, the sharper bound is
\[
 \delta_W(A,L)\le8C_{4d}f^3\Delta
                  <10^{178}d^{51}f^3\Delta.
\]
The detecting cycle may be chosen reduced and irreducible.
\end{theorem}

We prove the theorem in four steps. First a genus-three curve supplies quotient ideals whose obstruction trace is injective. Two distinct elliptic directions make their Fourier transform, after duality, a coherent sheaf of projective dimension at most one. An explicit regularity estimate then turns its local deformation into a globally generated vector bundle, whose third Chern class is an effective detector. Finally a fixed isogeny places this local construction in the principal hyperbolic family, and Theorem~\ref{thm:principal-quotient} transports it to every target lattice.

\subsection{The Fermat curve and its quotient ideals}
Let $C\subset\mathbf P^2$ be $x^4+y^4+z^4=0$. It is a smooth nonhyperelliptic genus-three curve. The three sign involutions give maps $f_j:C\to E_i$ of degree two. Each has four fixed points, so its quotient has genus one. One quotient is $u^2=-(t^4+1)$; the automorphism $(t,u)\mapsto(it,u)$ fixes a point and acts by $i$ on its differential, identifying its Jacobian with $E_i$. The Klein-four quotient of $C$ is the conic $X^2+Y^2+Z^2=0$, so the three nontrivial characters exhaust $H^1(C,\Q)$.

Put $X=E_i^3$ and let $H$ be its product principal polarization. The map
\[
 q:X\longrightarrow J(C),\qquad q(z_1,z_2,z_3)=\sum_j f_j^*z_j
\]
satisfies $q^*\Theta=2H$, $\deg q=8$ and $\ker q\subset X[2]$. Indeed the cross norms $f_{j,*}f_k^*$ vanish for distinct Klein characters and the diagonal norms are $[2]$. Thus the polarization matrix is diagonal $[2]$; Riemann--Roch gives the degree, and the norm identity gives the kernel assertion.

Each $f_j^*$ is injective. A nonzero kernel point would be two-torsion and would factor $C\to E_i$ through a connected \'etale double cover, forcing the remaining map to have degree one and contradicting $g(C)=3$. For an Abel image of $C$, one has the stronger identity
\begin{equation}\label{eq:bielliptic-disjointness}
 f_j^*\operatorname{Pic}^0(E_i)\cap(C-C)=\{0\}.
\end{equation}
To prove it, write $p-q=f_j^*M$ and let $\sigma$ be the involution. Then $p+\sigma q\sim q+\sigma p$. On a nonhyperelliptic curve equal degree-two divisor classes have equal effective representatives. If $p\ne q$, equality of these unordered pairs forces both points to be fixed by $\sigma$. The natural linearization on $\mathcal O_C(p-q)$ has fibre sign $-1$ at these two fixed points and $+1$ at the other two: a local parameter and its inverse both change sign. Any other linearization differs by one global character. Neither choice has trivial inertia at all four points, whereas a pullback from the quotient does. This proves \eqref{eq:bielliptic-disjointness}.

Set $D=4d$ and $k=D+1$. For $j=1,2$, choose a point
$\tau_j\in E_i[k]$ of exact order $k$, and let
$u_j:J(C)\to Y_j=J(C)/\langle f_j^*\tau_j\rangle$ be the quotient.
It embeds $C$ as a smooth curve $\bar C_j\subset Y_j$, because its
translates are disjoint by \eqref{eq:bielliptic-disjointness}.
Let $h_j\in\operatorname{NS}(Y_j)_{\Q}$ be the rational polarization
class with $u_j^*h_j=\Theta$. Then
\[
 \operatorname{ch}(I_{\bar C_j})=1-\frac{k}{2}h_j^2+\frac{k}{3}h_j^3.
\]
This follows by pulling back to the disjoint union of $k$ Abel curves, using $[C]=\Theta^2/2$ and $\chi(\mathcal O_C)=-2$.

For completeness, we check the whole obstruction trace. The restriction $H^1(\mathcal O_{Y_j})\to H^1(\mathcal O_C)$ is an isomorphism. In the normal sequence, the dual of $H^1(T_C)\to H^1(T_{Y_j}|_C)$ contains
$\operatorname{Sym}^2H^0(K_C)\to H^0(K_C^2)$. This map is an isomorphism: both dimensions are six and a plane quartic has no quadratic equation. Hence $h^0(N_{C/Y_j})=3$.
The ideal sequence and Serre duality bound $\operatorname{ext}^1(I,I)$ above by $3+3$. Translations give three independent determinant-preserving directions, while tensoring by $\operatorname{Pic}^0$ gives three independent determinant directions. Thus
\[
 \operatorname{ext}^1(I,I)=\operatorname{ext}^2(I,I)=6.
\]

The categorical trace followed by Hochschild--Kostant--Rosenberg identifies evaluation on $I$ with contraction against $\operatorname{ch}(I)$; the precise compatibility is \cite[Proposition~6.2.1, Corollary~6.3.2 and Theorem~6.4.2, author pp.~29--31]{BFHH}. The tangent bundle is trivial, so the Todd correction is one. The degree-two target is
\[
 HH_{-2}(Y_j)=H^2(\mathcal O_{Y_j})\oplus H^3(\Omega^1_{Y_j}),
\]
of dimension six. The first block is reached isomorphically through the rank term. Contraction from $H^1(T_{Y_j})$ with the nonzero multiple $-kh_j^2/2$ surjects onto the second block: contraction with $h_j$ reaches $H^2(\mathcal O_{Y_j})$, and wedging with $h_j$ identifies that space with $H^3(\Omega^1_{Y_j})$. In coordinates $h_j=\sum x_iy_i$, the latter sends $y_1y_2$ to a nonzero multiple of $x_3y_1y_2y_3$, and likewise for the other basis vectors. The trace therefore has rank six and is injective. In degree one the same calculation separates the Picard directions through the rank term and the translation directions through contraction with $h_j^2$. The degree-zero trace is the nonzero Chern character.

External product preserves this injectivity in degree two: the three K\"unneth terms of self-Ext map to the distinct Hochschild bidegrees $(2,0)$, $(1,1)$ and $(0,2)$, by tensor products of the three injective trace maps just checked. Pullback from a finite quotient identifies invariant self-Ext with the self-Ext downstairs. Trace is natural under that pullback. Thus the pullbacks of these two quotient ideals have a coherent finite descent action and injective invariant obstruction trace.

\subsection{A Fourier transform that is a sheaf after duality}
Put $A=X\times\widehat X$, and write $H_x,H_\alpha$ for the product principal classes on its factors. Write $r_j=u_jq$, $K_j=\ker r_j$, and $F_j=r_j^*I_{\bar C_j}=I_{Z_j}$. Then $|K_j|=8k$; concretely $K_1=\ker q+\langle(\tau_1,0,0)\rangle$ and $K_2=\ker q+\langle(0,\tau_2,0)\rangle$. These summands intersect trivially by injectivity of $f_j^*$. Fourier transform with the normalized Poincar\'e bundle is an equivalence by \cite[Theorem~2.2, p.~156]{Mukai}; its translation, tensor and isogeny identities are in \cite[\S3, especially (3.1) and Proposition~3.11]{Mukai}.
Let $L=\mathcal O_X(2H)$, $\mu(x,y)=(x+y,y)$, and use the inverse
Fourier kernel $\mathcal P^{-1}[3]$ in the second variable. Define
\begin{equation}\label{eq:effective-transform}
 E_{\mathrm{ref}}=
 (\operatorname{id}\times\mathrm{FM}_{\mathcal P^{-1}[3]})
 \mu^*\bigl(([-1]^*F_2\otimes L)\boxtimes(F_1\otimes L)\bigr).
\end{equation}
The order of the factors gives the following incidence description. On $X_y\times A$, put
\[
 a(y,x)=y+x,\quad
 \mathcal L=L_y\otimes a^*L\otimes\mathcal P_{y,\alpha}^{-1},
 \quad D_1=\{y\in Z_1\},\quad D_2=\{y+x\in-Z_2\}.
\]
For $\pi:X_y\times A\to A$ and $T=D_1\cap D_2$, set
\begin{equation}\label{eq:effective-duality}
 G=R\pi_*(\mathcal L I_{D_1}I_{D_2}),
 \qquad E_{\mathrm{ref}}=G[3],\qquad F=G^\vee[-1]
               =E_{\mathrm{ref}}^\vee[2].                 
\end{equation}
Every equality here is a derived-functor identification.

Here is the additional geometric fact that makes $F$ a sheaf.
The curve $C'=q^{-1}C$ is connected: the Abel map induces the
surjection $\pi_1(C)\to\pi_1(J(C))$, hence has transitive monodromy
on the fibre of the connected etale cover $q$. It is an etale
degree-eight cover of $C$, and therefore has genus $17$. Projection
of $C'$ along the first elliptic factor is injective. Indeed, two
points in one such fibre differ by a first-factor point; their
images in $C$ differ by an element of $f_1^*\operatorname{Pic}^0(E_i)$, so \eqref{eq:bielliptic-disjointness} makes those
images equal. The difference then lies in
$\ker q\cap(E_i\times0\times0)=0$, since $f_1^*$ is injective.
The same argument applies to the second factor.

Write $\eta_j$ for the pullback of the point class from the $j$th
elliptic factor of $X$, so $H=\eta_1+\eta_2+\eta_3$.
The two projected images are irreducible divisors $B_1,B_2$ on
the corresponding elliptic surfaces. Their projection maps have
degree one, so the inverse-image divisors $S_1,S_2$ on $X$
have classes
\begin{equation}\label{eq:effective-divisor-classes}
 [S_1]=4(\eta_2+\eta_3),\qquad
 [S_2]=4(\eta_1+\eta_3).                                  
\end{equation}
This follows by pushing forward
$[C']=q^*(\Theta^2/2)=2H^2$ along each elliptic projection.
There is no unaccounted projection degree. One may also check
smoothness: the projected divisor has self-intersection $32$ and
arithmetic genus $17$, equal to the genus of its degree-one
normalization $C'$, so the normalization has no singularity defect.

The $k$ translates defining $Z_1$ lie on $S_1$. On every generic
first-factor fibre of $S_1$ they supply $k$ distinct points.
The $k$ reflected translates in $-Z_2-x$ have the analogous
property on the second-factor fibres of $-S_2-x$. A section of
$\mathcal O_X(4H)\otimes P$ vanishing on both unions, where
$P\in\operatorname{Pic}^0(X)$ is arbitrary, must vanish on each
of these entire elliptic fibres: its restriction has degree four,
whereas $k=4d+1\ge5$. Hence its divisor contains both $S_1$ and
$-S_2-x$. They are distinct for every $x$, because their classes
in \eqref{eq:effective-divisor-classes} differ. The residual effective divisor would have class
\[
 4H-[S_1]-[S_2]=-4\eta_3.
\]
Its intersection with the ample class $H^2$ would be $-8$, which
is impossible for an effective divisor. Thus that space of
sections is zero, for every $x$ and every flat factor $P$.

\paragraph{An explicit derived-to-coherent calculation.}
The coordinates $(y,z=y+x)$ identify the two curve-union ideals
with pullbacks from independent factors. They are Tor independent,
and their product equals their intersection. Therefore
\begin{equation}\label{eq:effective-ideal-sequence}
 0\to\mathcal L I_{D_1}I_{D_2}\to\mathcal L
 \to\mathcal L|_{D_1}\oplus\mathcal L|_{D_2}
 \to\mathcal L|_T\to0                                    
\end{equation}
is exact on the total space.
On a fibre of $\pi$, $\mathcal L$ has class $4H$, with
$h^0=64$ and no higher cohomology. Its restriction to a component
$C'$ has degree $48$. Since $g(C')=17$, its first cohomology is zero
and its space of sections has dimension $32$. Consequently
\[
 W=\pi_*\mathcal L,\qquad
 U=\pi_*(\mathcal L|_{D_1}\oplus\mathcal L|_{D_2})
\]
are vector bundles of ranks $64$ and $64k$, respectively.

The map $f:T\to A$ is finite flat over each of its smooth divisor
images. To see the precise degree, ignore translations and write
one component as $C'\times C'\times\widehat X$, with its map to
$X$ given, up to sign, by addition. The diagonal action
$(p',q')\mapsto(p'+h,q'-h)$ of $\ker q$ is free. Its quotient is
the base change of $C\times C\to C^{(2)}\simeq\Theta$ by $q$.
The quotient map is etale of degree eight; the latter map is
finite flat of degree two onto the smooth divisor $q^{-1}\Theta$.
Thus the composite has degree $16$. Smoothness of $\Theta$ follows
from the absence of a degree-two pencil on the nonhyperelliptic
genus-three curve. Each summand of
\[
 \mathcal T=f_*(\mathcal L|_T)
\]
is a vector bundle on a smooth divisor, pushed forward to $A$.
It has projective dimension one.

Applying direct image to \eqref{eq:effective-ideal-sequence}, including the preceding higher
cohomology vanishings, represents $G$ by
\begin{equation}\label{eq:effective-three-terms}
 [W\longrightarrow U\longrightarrow\mathcal T]
 \quad\hbox{in degrees }0,1,2.                            
\end{equation}
The section argument proves that $W\to U$ is fibrewise injective.
Hence $Q=U/W$ is a vector bundle of rank $64(k-1)=r$, and
$G=[Q\to\mathcal T]$ in degrees $1,2$. No assertion that fibrewise
$H^2(G)$ vanishes has been made or used.

Finite duality gives
\[
 \mathcal T^D=\mathcal E xt^1_A(\mathcal T,\mathcal O_A)
 =\bigoplus f_{ab,*}
       ((\mathcal L|_{T_{ab}})^{-1}\otimes\omega_{T_{ab}}),
\]
with all other sheaf Ext groups zero. Dualizing the last two-term
complex and shifting by $-1$ gives an ordinary coherent sheaf and
the exact extension
\begin{equation}\label{eq:effective-extension}
 0\longrightarrow Q^*\longrightarrow F
   \longrightarrow\mathcal T^D\longrightarrow0.          
\end{equation}
Both the local extension and its finite-duality description prove
$\operatorname{pd}F\le1$; local freeness outside a prescribed
codimension is unnecessary.

\paragraph{The finite action and its trace.}
In \eqref{eq:effective-transform}, the first translation parameter belongs to $K_2$ and
the second to $K_1$. With the convention $t_x^*$ for translation,
the transported action on $E_{\mathrm{ref}}$ has point map
\begin{equation}\label{eq:effective-action}
 (x,y)\longmapsto
 \bigl((x-y,2(x+y)),\,(-2x,y)\bigr)
 \quad\hbox{in }A\times\widehat A.                       
\end{equation}
Tensoring by $L$ sends a translation parameter $x$ to translation $x$ with tensor parameter $-2x$. Pullback by $\mu$ sends the two translation coordinates to $(x-y,y)$ and the tensor coordinates to $(-2x,-2(x+y))$. The inverse Fourier kernel exchanges translation $y$ with tensor $P_y$, and tensor $P_b$ with translation $-b$. This proves the displayed map. Its last coordinate recovers $y$ and its first then recovers $x$, so it is injective, with group order $64k^2$. Its coherence is transported from finite descent. Duality replaces a translation/tensor
point $(a,b)$ by $(a,-b)$; invert the dual linearizations to obtain
the coherent action on $F$. The trace is preserved up to invertible
signs under this operation. Explicitly, derived transpose identifies
the degree-two self-Ext spaces, the Atiyah class of the dual is
the negative transpose, and graded cyclicity of the trace gives
the sign $(-1)^p$ on
$\operatorname{tr}(\xi\operatorname{At}^p/p!)$.
The reversal signs for the cohomological and differential-form
factors cancel. Thus the complete trace kernels identify, including
the invariant subspaces. The even shift $[2]$ contributes sign $+1$.
This proves the claimed strong equivariant semiregularity of $F$.
Negative self-Ext groups vanish because $F$ is a sheaf; they also
vanish through the displayed equivalence with the original ideals.

\subsection{The entire Chern character and the split Weil action}
Put $s^2=-D$, $\theta_D=4DH_x+H_\alpha$, $B=-P/2$ and $\lambda_D=(D+1)/(2D)$, where $P$ is the first Chern class of the normalized Poincar\'e bundle on $X\times\widehat X$. The factor used before Fourier transform has character
\[
 \operatorname{ch}(F_j\otimes L)=1+2H-2DH^2-\frac{4D}{3}H^3
     =\alpha_+e^{2sH}+\alpha_-e^{-2sH},
 \quad\alpha_\pm=\frac{1\pm1/s}{2}.
\]
Reflection preserves this even character.

Here is the Fourier calculation, including the signs. On one elliptic pair let $a,b$ be the first coordinates, $c,e$ the integrated coordinates, and $u,v$ the dual coordinates; normalize $\int ce=1$. For $s_1,s_2\in\{s,-s\}$ the kernel exponent is
\[
 2s_1(a+c)(b+e)+2s_2ce-(cv-eu).
\]
Its integral is
\[
 e^{2s_1ab}\{2(s_1+s_2)-(2s_1b-v)(u-2s_1a)\}.
\]
This follows by collecting the coefficient of $ce$ in the exponential. Multiplying by $e^{-(av-bu)/2}$ and expanding the remaining nilpotent terms gives
\[
 \begin{cases}
 4s_1\exp(s_1(4Dab+uv)/(4D)),&s_1=s_2,\\
 \Omega_{s_1}:=4Dab+2s_1(av-bu)-uv,&s_1=-s_2.
 \end{cases}
\]
Take the product over the three elliptic pairs. The shift $[3]$ in the inverse Fourier transform contributes a minus sign. The equal-sign terms give the polynomial
\begin{equation}\label{eq:effective-polynomial}
 \begin{split}
 P_D(z)={}&64D-8(D-1)z-2z^2+\frac{D-1}{12D}z^3\\
 &+\frac{z^4}{96D}-\frac{D-1}{3840D^2}z^5
       -\frac{z^6}{46080D^2}.
 \end{split}
\end{equation}
The opposite-sign terms give $\lambda_Dt_D$, where
\[
 t_D=-\frac12\left(\prod_{j=1}^3\Omega_{s,j}
                      +\prod_{j=1}^3\Omega_{-s,j}\right).
\]
Duality changes the sign of the codimension-$j$ character by $(-1)^j$; the shift $[2]$ has sign $+1$. Consequently the sheaf $F$ satisfies the full identity
\begin{equation}\label{eq:effective-character}
 e^{-B}\operatorname{ch}(F)=-\lambda_Dt_D+P_D(-\theta_D).
\end{equation}

The rational action of $s$ on each elliptic pair is
$J_D=\left(\begin{smallmatrix}0&1/2\\-2D&0\end{smallmatrix}\right)$.
It squares to $-D$, is skew-adjoint for $\theta_D$, and has signature $(3,3)$. Since
$\Omega_s=-(u-2sa)(v-2sb)$, its conjugate products are the two Weil determinant lines. The class $t_D$ is primitive integral: the odd powers of $s$ cancel and the coefficient of $\prod_j u_jv_j$ is one. Direct multiplication gives
\[
 \theta_Dt_D=0,\qquad \int t_D^2=-2048D^3,
 \qquad \chi(\theta_D)=64D^3.
\]
The Hermitian space is split. For example, the two real cycles in the first elliptic factor form a $K$-basis of each pair, with zero diagonal Hermitian entries and nonzero conjugate off-diagonal entries; rescaling one vector gives $U$.

Set $r=64D$ and $\mathcal H=3(H_x+H_\alpha)$. We prove the following stronger form of $560$-regularity: for every $m\ge1$ and every line bundle $M$ such that $c_1(M)-560m^2\mathcal H$ is nef,
\begin{equation}\label{eq:effective-vanishing}
 H^j(A,[m]^*F\otimes M)=0\qquad(j>0).
\end{equation}
\subsection{A regularity estimate valid after every multiplication isogeny}
All the following vanishings allow arbitrary flat factors.
Relative duality and Serre duality show, for every integer $q\ge1$,
\begin{equation}\label{eq:effective-duality-vanishing}
 H^j(A,W^*(q\mathcal H))=H^j(A,U^*(q\mathcal H))=0
 \qquad(j>0).                                             
\end{equation}
Here is the numerical positivity test, so that \eqref{eq:effective-duality-vanishing} does not
depend on a qualitative regularity assertion. Project the duality
calculation onto the fixed $y$ coordinate, or onto the fixed curve
coordinate for a summand of $U$. The line bundle on the remaining
six-dimensional fibre has matrix
\[
 \begin{pmatrix}2-3q&0\\0&-3q\end{pmatrix}
 \quad\hbox{or}\quad
 \begin{pmatrix}2-3q&1\\1&-3q\end{pmatrix}.
\]
Both are negative definite: the largest eigenvalue before
subtracting $3qI$ is at most $1+\sqrt2<3$. Only the top direct
image of this antiample line bundle can be nonzero, which gives
the required cohomology vanishing below degree six on the dual
side. This proves \eqref{eq:effective-duality-vanishing}. In particular $W^*$ and $U^*$ are
$7$-regular.

Take the evaluation sequence
\[
 0\to K_U\to\mathcal O_A(-7\mathcal H)^N\to U^*\to0.
\]
Its kernel is locally free and $14$-regular. For degrees at least
two this follows directly from \eqref{eq:effective-duality-vanishing} and positivity of
$(7-j)\mathcal H$. For degree one, the multiplication map from
the degree-seven sections, followed by six hyperplanes in the
embedding defined by $\mathcal H$, is surjective by regularity.
Compose the evaluation map with $U^*\to W^*$ and call the kernel
$J$. Then
\begin{equation}\label{eq:effective-kernel-sequence}
 0\to K_U\to J\to Q^*\to0.                              
\end{equation}
The Buchsbaum--Rim resolution of $J$ has terms
\begin{equation}\label{eq:effective-br}
 C_j=\bigwedge^{64+j}\mathcal O_A(-7\mathcal H)^N
          \otimes\operatorname{Sym}^{j-1}W\otimes\det W,
 \qquad j\ge1.                                           
\end{equation}
The first term surjects onto $J$. Local exactness follows by
splitting the surjection as $K\oplus V\to V$ and decomposing
according to the number of factors from $K$. The resulting
strands are the polynomial de Rham complex in positive degree,
contracted by the Euler vector field in characteristic zero;
the degree-one strand maps isomorphically to $K$.

For an explicit bound on \eqref{eq:effective-br}, the theorem of the square gives
\[
 W=\mathcal O_A(2H_x)\otimes b^*\mathrm{FM}(\mathcal O_X(4H)),
 \quad b=2x-\alpha,
 \quad [4]^*W=Q_0^{\oplus64},
\]
with
\begin{equation}\label{eq:effective-four-pullback}
 c_1(Q_0)=16H_x-4H_\alpha+8P,\qquad
 \operatorname{Mat}(Q_0)=\begin{pmatrix}16&8\\8&-4\end{pmatrix}.
\end{equation}
Indeed the pullback of the Fourier bundle by the polarization
isogeny of $\mathcal O(4H)$ is
$H^0(\mathcal O(4H))\otimes\mathcal O(-4H)$; the displayed
formula is its pullback by $b$ with the first line-bundle factor.
The eigenvalues of the matrix in \eqref{eq:effective-four-pullback} are $6\pm2\sqrt{41}$,
whose minimum is greater than $-7$, while
$[4]^*\mathcal H$ has matrix $48I$. The symmetric and determinant
factor in \eqref{eq:effective-br} pulls back to copies of $Q_0^{j+63}$.
Kodaira vanishing and the trace splitting of the isogeny therefore
give
\begin{equation}\label{eq:effective-term-regularity}
 \operatorname{reg}_{\mathcal H}
   (\operatorname{Sym}^{j-1}W\otimes\det W)\le j+69,
 \qquad \operatorname{reg}C_j\le517+8j.                  
\end{equation}
Only six terms are needed to bound $J$ on a sixfold: successive
cohomology sequences leave a final kernel in cohomological
degree seven. Thus
\begin{equation}\label{eq:effective-reg560}
 \operatorname{reg}J\le
 \max_{1\le j\le6}(\operatorname{reg}C_j-j+1)
 \le560,\qquad\operatorname{reg}Q^*\le560.               
\end{equation}

For the torsion dual, use $\rho:C'\times C'\to C\times C$.
On $C\times C$ write $F_1,F_2$ for the point fibres,
$S=F_1+F_2$, and $\Delta_C$ for the diagonal. Addition pulls
$\Theta$ back to $A_C=4S-\Delta_C$. The divisor
$R_C=3S-\Delta_C$ is nef. A degree-three map $C\to\mathbb P^1$
from plane projection supplies its effective residual fibre
correspondence. If its monodromy is $S_3$, that residual divisor
is irreducible and has square $2$, hence is nef. If its monodromy
is cyclic of order three, it is the sum of the two automorphism
graphs: each has square $-4$, and they intersect transversely
in the five fixed points given by Riemann--Hurwitz. Their sum
therefore has intersection $1$ with each component and nonnegative
intersection with every other curve, again proving nefness.

Put $A'=\operatorname{add}_X^*H=\rho^*A_C/2$ and
$H'_i=\rho^*(3F_i/2)$. Up to flat factors on each translated
component,
$c_1(\mathcal L_T)=2H'_1+2H'_2-P_{1,\alpha}$ and
$f^*\mathcal H=3A'+3H_\alpha$. For $q\ge4$ decompose
\begin{equation}\label{eq:effective-torsion-positive}
 \begin{split}
 qf^*\mathcal H-c_1(\mathcal L_T)
 ={}&[3qA'-3H'_1-2H'_2]+(3q-1)H_\alpha\\
    &+[H'_1+H_\alpha+P_{1,\alpha}].
 \end{split}                                               
\end{equation}
The last bracket is a nef addition pullback. The first equals
\[
 \rho^*\left[\frac{3q}{2}R_C+
       \left(\frac{3q}{2}-\frac92\right)F_1+
       \left(\frac{3q}{2}-3\right)F_2\right]
\]
and is ample. Consequently \eqref{eq:effective-torsion-positive} is ample on the smooth
fivefold $C'\times C'\times\widehat X$. Kodaira vanishing and
finite duality give $9$-regularity of $\mathcal T^D$.
Together with \eqref{eq:effective-extension} and \eqref{eq:effective-reg560}, this proves the stated
$560$-regularity of $F$.

For the stronger assertion \eqref{eq:effective-vanishing}, pull the exact constructions
back by $[m]$. The matrices proving \eqref{eq:effective-duality-vanishing} are multiplied by
$m^2$, so the assumed nef difference makes the duality
vanishings hold after tensoring by $M$. The pulled evaluation
sequence has ample middle terms and gives vanishing of
$H^j([m]^*K_U\otimes M)$ for $j\ge2$.
After one further pullback by $[4]$, the first six terms of
\eqref{eq:effective-br} are sums of line bundles whose matrices, after removal
of a nef class, are
\[
 m^2\{48(112-7j)I+(j+63)\operatorname{Mat}(Q_0)\},
 \qquad1\le j\le6.
\]
Their least eigenvalue is at least
$m^2(3360-7\cdot69)=2877m^2>0$. The same six-term truncation
and \eqref{eq:effective-kernel-sequence} give the desired vanishing for $[m]^*Q^*\otimes M$.
For the torsion term, finite flat base change gives a smooth
etale cover $T'$ of $T$; the required line bundle is the sum
of the pullback of \eqref{eq:effective-torsion-positive} with $q=4$ and a nef class.
Kodaira vanishing on $T'$ proves the last vanishing.
The extension \eqref{eq:effective-extension} proves \eqref{eq:effective-vanishing}.

\paragraph{Normalization with an explicit regularity integer.}
Keep $D=4d$, $r=64D$ and the sheaf $F$ just constructed. Set
\begin{equation}\label{eq:effective-normalization}
\begin{gathered}
 h=3\theta_D,\quad \mathcal D=\det F,\quad
 a=\frac{D-1}{8D},\\
 R=560r^2+8Dr+6=2294272D^2+6,\qquad
 B_0=[r]^*F\otimes\mathcal D^{-r}.
\end{gathered}                                             
\end{equation}
The notation $B_0$ here denotes the ordinary normalized sheaf;
the rational B-field used below is written $\mathcal B_0$.
The class $h-\mathcal H$ is nef. Formula \eqref{eq:effective-character} gives
$c_1(F)=-32DP+8(D-1)\theta_D$, and the class
$8Dh-c_1(F)$ has matrix
\[
 \begin{pmatrix}32D(2D+1)&32D\\32D&8(2D+1)\end{pmatrix},
 \qquad\det=256D(4D^2+1)>0.
\]
It is ample. For $1\le i\le6$, the line bundle
$\mathcal D^{-r}\otimes h^{R-i}$ therefore satisfies
\[
 c_1(\mathcal D^{-r}\otimes h^{R-i})-560r^2\mathcal H
 =560r^2(h-\mathcal H)+r(8Dh-c_1(F))+(6-i)h,
\]
which is ample. Assertion \eqref{eq:effective-vanishing} proves that $B_0$ is
$R$-regular. It still has projective dimension at most one.
Its complete character is
\begin{equation}\label{eq:effective-normal-character}
 \operatorname{ch}(B_0)=e^{-r^2a\theta_D}
       \{-\lambda_Dr^6t_D+P_D(-r^2\theta_D)\}.            
\end{equation}
In particular $c_1(B_0)=0$, and, writing its first two relevant
scalar coefficients relative to $h$,
\begin{equation}\label{eq:effective-a2a3}
 \begin{aligned}
 \operatorname{ch}_2(B_0)&=a_2h^2,
 &a_2&=-\frac{r^4(D+1)^2}{18D},\\
 \operatorname{ch}_3(B_0)&=-\lambda_Dr^6t_D+a_3h^3,
 &a_3&=\frac{r^6(D-1)(D+1)^2}{648D^2}.
 \end{aligned}                                             
\end{equation}

\paragraph{Deformation and the ordinary sheaf family.}
Apply Perry's Theorem~1.2 \cite{Perry} to the sheaf $F$, the
coherent finite action proved above, and the algebraic
rational B-field $\mathcal B_0=-c_1(F)/r$. The whole twisted
character is the normalized expression in \eqref{eq:effective-normal-character} before the
$[r]$ scaling: it is a polynomial in the fixed polarization
plus a nonzero Weil class. It is flat and remains of Hodge type
through the balanced Weil family. The determinant monodromy is an integral unit of $K$ of norm one and hence a root of unity; a finite cover makes the Weil class flat.

For precision about untwisting, Perry's \S6.2, proof of Theorem~1.2, pp.~26--28, begins with the $\mu_n$-gerbe representing
the B-field denominator. Here take $n=r$. Its subsequent pushout
to $\mu_m$ does not change the associated multiplicative-group gerbe. Pullback along the map from the original $\mu_r$-gerbe to this pushout is an equivalence on the weight-one categories: locally both are modules on which the band acts through its standard character, and the descent cocycle is unchanged. We therefore pull the deformed weight-one object back to the original finite-band gerbe before untwisting. On a
contractible analytic base the original class remains in the
local system $H^2(A,\mu_r)$. Multiplication by $r$ acts as
$r^2$ on this group and therefore kills it. Pullback by $[r]$
and an analytic Cech trivialization give an ordinary perfect
family. Here is its exact seed comparison, including the sign.
The seed gerbe is the $r$th-root gerbe of $\mathcal D^{-1}$;
its tautological weight-one line $\mathcal T_0$ satisfies
$\mathcal T_0^r=\mathcal D^{-1}$, and the seed object is
$F\otimes\mathcal T_0$. Write its class without an exponential
convention as the Kummer boundary
$\vartheta_0=\delta(\mathcal D^{-1})\in H^2(A,\mu_r)$.
Its rational B-field is
$c_1(\mathcal T_0)=c_1(\mathcal T_0^r)/r
=-c_1(\mathcal D)/r=\mathcal B_0$.
After pullback by $[r]$ and neutralization, the line becomes
an ordinary line $N_0$ with
$N_0^r=[r]^*\mathcal D^{-1}$. Hence
$c_1(N_0)=-r c_1(\mathcal D)$ and
$N_0=\mathcal D^{-r}\otimes P_0$ for some
$P_0\in\operatorname{Pic}^0(A)$. The ordinary seed is therefore
exactly $B_0\otimes P_0$, as asserted.
For these conventions see \cite[Definition~4.5, Lemma~4.6, Remark~4.13, Definition~4.16 and Theorem~4.17, author pp.~15--18]{Hotchkiss}. The displayed root and tensor identities give the required
comparison directly. All preceding estimates allow the flat
factor $P_0$.

The locus where the derived fibres of a relative perfect
complex have cohomology only in degree zero is open. At the
seed this property holds by \eqref{eq:effective-extension}; after shrinking the base
the ordinary family is therefore a coherent sheaf flat over
the base. This can equally be seen by the relative Tor-amplitude
criterion applied to a local finite vector-bundle resolution.
Properness permits a common neighbourhood for the whole fibre.
Analytic upper semicontinuity preserves the six regularity
vanishings of $B_0$ on a smaller neighbourhood. No stability
of all sheaves with this character is used.

\paragraph{A vector-bundle kernel supplies the effective cycle.}
At the seed and throughout that neighbourhood, evaluation gives
\begin{equation}\label{eq:effective-evaluation}
 0\longrightarrow V\longrightarrow
    \mathcal O_A(-Rh)^N\longrightarrow B_0\longrightarrow0,
 \qquad N=h^0(B_0(Rh))=\chi(B_0(Rh)).                     
\end{equation}
At the seed $V$ is locally free, because $\operatorname{pd}B_0\le1$.
This follows locally from the exact Ext sequence, or by comparing
the evaluation presentation with a projective resolution of
length one. In the family the direct image of $B_0(Rh)$ is a
rank-$N$ vector bundle, by the regularity vanishings and base
change; evaluation is surjective. Its kernel is flat, and its
locally free locus is open. The closed complement misses the
proper seed fibre, so it misses every fibre after one more
shrinking. Thus \eqref{eq:effective-evaluation} is a relative vector-bundle
presentation on the neighbourhood.

The kernel $V$ is $(R+7)$-regular. For cohomological degrees
at least two, use the sequence \eqref{eq:effective-evaluation}, the regularity of $B_0$,
and the positive line bundles $(7-i)h$ for $2\le i\le6$.
For degree one, use the embedding $A\hookrightarrow\mathbb P^M$ defined by $h$.
Then $R$-regularity gives the surjection
\[
 H^0(B_0(Rh))\otimes\operatorname{Sym}^6H^0(\mathbb P^M,\mathcal O(1))
 \longrightarrow H^0(B_0((R+6)h)),
\]
which factors through the degree-six sections in \eqref{eq:effective-evaluation}.
It gives the remaining $H^1$ vanishing. Therefore
$V'=V((R+7)h)$ is a globally generated vector bundle.

Let $q_0=R+7$ and define the explicit rational numbers
\begin{equation}\label{eq:effective-uvw}
 \begin{split}
 u&=7N-rq_0,\\
 v&=\frac{49N}{2}-\frac{rq_0^2}{2}-a_2,\\
 w&=\frac{343N}{6}-\frac{rq_0^3}{6}-q_0a_2-a_3.
 \end{split}                                               
\end{equation}
The identity in $K_0(A)$ is
\[
 [V']=N[\mathcal O_A(7h)]-[B_0(q_0h)].
\]
Consequently $c_1(V')=uh$, $\operatorname{ch}_2(V')=vh^2$,
and $\operatorname{ch}_3(V')=wh^3+\lambda_Dr^6t_D$.
The exact Newton identity
$c_3=2\operatorname{ch}_3-c_1\operatorname{ch}_2+c_1^3/6$
therefore gives
\begin{equation}\label{eq:effective-c3}
 c_3(V')=2\lambda_Dr^6t_D+
                (2w-uv+u^3/6)h^3.                       
\end{equation}
Its Weil projection is nonzero. In particular the rank of $V'$
is at least three, since lower-rank bundles have $c_3=0$.

General sections of a globally generated vector bundle give
an effective maximal-rank-drop locus of expected codimension
three, or the empty locus, representing $c_3(V')$. This is the
elementary Thom--Porteous construction: take
$\operatorname{rank}(V')-2$ sections, use the universal surjective
evaluation to get transversality to the rank strata, and then
take a general choice. Equation \eqref{eq:effective-c3} rules out the empty case.
It is an effective cycle with integral coefficients with nonzero Weil projection,
so at least one reduced irreducible component detects the Weil
plane. No condition on the codimension of the nonfree locus of
$B_0$ is required.

Write $P_D(z)=\sum_{j=0}^6p_jz^j$, with coefficients given in
\eqref{eq:effective-polynomial}. Since the Todd class of an abelian variety is $1$ and
$\theta_Dt_D=0$, the integer $N$ is explicitly
\[
 N=(6!\,64D^3)\sum_{j=0}^6
       p_j(-r^2)^j\frac{(3R-r^2a)^{6-j}}{(6-j)!}.
\]
The exact degree of the effective cycle above is the nonnegative
integer
\begin{equation}\label{eq:effective-degree}
 C_D=(3^6\,6!\,64D^3)(2w-uv+u^3/6).                     
\end{equation}
The formula defines every quantity by rational operations on
the specified positive integer $D$. Integrality and nonnegativity
also follow from the constructed vector bundle and effective
cycle. These same Chern classes and degrees persist on the
analytic neighbourhood. GAGA makes each fibre and cycle algebraic.

\subsection{The degree and its transport to every field and lattice}
The formulas above specify $C_D$ by rational arithmetic in $D$. We first give a convenient upper bound. Put $s_0=3R-r^2a$. From \eqref{eq:effective-polynomial}, $j!|p_j|\le64D$ for $0\le j\le6$, while $0<s_0$ and $s_0+r^2<7\cdot10^6D^2$. The Euler-characteristic formula therefore yields
\[
 0<N\le4096D^4(s_0+r^2)^6<10^{45}D^{16}.
\]
Also $q_0<3\cdot10^6D^2$, $|a_2|<4\cdot10^6D^5$ and $|a_3|<5\cdot10^8D^7$. Substitution in \eqref{eq:effective-uvw} gives
\[
 |u|\le8\cdot10^{45}D^{16},\quad
 |v|\le26\cdot10^{45}D^{16},\quad
 |w|\le59\cdot10^{45}D^{16}.
\]
As $10^{45}D^{16}\ge1$, the absolute value of $2w-uv+u^3/6$ is at most $512\cdot10^{135}D^{48}$. Hence
\[
 0<C_D<10^{146}D^{51},\qquad C_{4d}<10^{177}d^{51}.
\]
Positivity follows also from the nonzero effective detecting cycle, and integrality follows from its degree.

We next give a fixed map from the source to a principal hyperbolic lattice. Put $s=\sqrt{-d}$ and let $e_d=1$ for $d\equiv1,2\pmod4$, and $e_d=2$ for $d\equiv3\pmod4$. Thus $\tau=s$ or $(1+s)/2$ generates $\mathcal O_K$, and $\tau-\bar\tau=2s/e_d$. In one elliptic pair write the real homology basis as $u=e_x$, $v=i e_x$, $w=e_y$, $z=i e_y$. The $s$-action is
\[
 su=-4dw,\quad sv=-4dz,\quad sw=u/4,\quad sz=v/4.
\]
The source polarization $h=3(16dH_x+H_y)$ has $E_h(u,v)=48d$ and $E_h(w,z)=3$. In the $K$-basis $u,v$ its normalized Hermitian matrix, linear in the first variable, is
\[
 \begin{pmatrix}0&48ds/e_d\\-48ds/e_d&0\end{pmatrix}.
\]
In the principal hyperbolic lattice $\mathcal O_K e\oplus\mathcal O_K f$, define
\[
 T(u)=4s e,\quad T(v)=12d f,\quad
 T(w)=e,\quad T(z)=-3s f.
\]
These values are integral and compatible with the $K$-action, and $H_U(Tu,Tv)=48ds$. Thus $T^*E_U=e_dE_h$. Three copies give an isogeny $\varphi:A_0\to P_U$ with $\varphi^*\Theta_U=e_dh$. In particular the pushforward of a three-cycle costs exactly $e_d^3$ in total cycle degree. Its Weil projection stays nonzero. This calculation uses no class-number assumption.

The fixed rational similitude identifies the period domains. After choosing a common finite level it sends the analytic neighbourhood where the sheaf construction works onto an analytic open subset of the principal $U^3$ family. Bounded-detector closedness therefore gives a detector of degree at most $e_d^3C_{4d}$ on every member of that family. Theorem~\ref{thm:principal-quotient}, and $e_dt_K=2$, now give
\[
 \delta_W(A,L)\le(t_Kfm_K)^3\Delta\,e_d^3C_{4d}
                 =8m_K^3C_{4d}f^3\Delta.
\]
Since $m_K\le4\sqrt d/\pi$, one has $8m_K^3<17d^{3/2}$. The asserted coarse bound follows from $C_{4d}<10^{177}d^{51}$ and $d^{105/2}\le d^{53}$. For class number one, $m_K=1$. A detecting irreducible component never has greater degree than the effective cycle containing it. This proves Theorem~\ref{thm:all-field-sixfolds}.

The exterior-algebra and integer-arithmetic calculations are independently reproducible with the programs \emph{Sixfold character verification} and \emph{Sixfold degree verification} in the \href{\workbenchurl}{workbench}. The proofs above work for every positive squarefree $d$. The programs check the formulas and signs, and do not replace the geometric deformation argument.
\section{A dominant Hurwitz family of split Eisenstein Weil eightfolds}
\label{sec:eisenstein-hurwitz-dominance}

Put $K=\mathbf Q(\omega)$, with $\omega^2+\omega+1=0$.
The following geometric construction is the source needed for the
split Eisenstein eightfold argument.

\begin{theorem}\label{thm:eisenstein-hurwitz-dominance}
The product-one tuple of twenty involutions defined below generates
$G_{34}$ and determines a seventeen-dimensional Hurwitz family of
connected covers of $\mathbf P^1$.  A primitive factor of their Jacobians
has dimension eight, a $K$-action of signature $(4,4)$, and split
Hermitian form.  Its period map has rank sixteen at a smooth point.
After finite level and marking, its algebraic moduli map is dominant
onto its connected split Weil component.
\end{theorem}

\paragraph{Proof overview.}
We compute the eight-dimensional parabolic multiplicity space directly
from the branch matrices.  Seventeen compatible boundary twists act on
it by square-zero operators with a common four-dimensional kernel and
image.  Their upper-right blocks span all four-by-four matrices.  A
positive-cone certificate ensures that a simultaneous degeneration has
maximal rank.  Schmid's limiting mixed Hodge structure then places the
limiting Hodge plane transverse to that common image, converting the
matrix span into sixteen independent period derivatives.  One actual
monodromy logarithm also proves that the geometric Hermitian form is
split.

The ten roots, cocycle representatives, and boundary intervals are the
numerical data of Severn's Appendix~A \cite{SevernData}.  They are all
printed here, together with the reconstruction rules and output
matrices.  The proof uses their exact verification, the published
$G_{34}$ root description, the Chevalley--Weil formula, geometric
parabolic cohomology, and Schmid's theorems.  It invokes no theorem or
geometric criterion from Severn's manuscript.

\subsection{The finite group and the branch datum}

Equip \(K^6\) with \(\langle x,y\rangle=\overline{x}^{\,t}y\).  The roots
\(v_i\) are the rows of the following array, regarded as column vectors when
forming matrices:
\[
\begin{array}{c|rrrrrr}
i&\multicolumn{6}{c}{v_i^t}\\\hline
1&1&-1-\omega&1&\omega&-1-\omega&\omega\\
2&1&0&0&1+\omega&0&0\\
3&0&1&1+\omega&0&0&0\\
4&1&\omega&-1-\omega&\omega&-1-\omega&1\\
5&1&-1-\omega&-1-\omega&-1-\omega&-1-\omega&\omega\\
6&1&-1-\omega&\omega&\omega&\omega&\omega\\
7&0&0&1&0&0&-1\\
8&1&\omega&\omega&-1-\omega&1&-1-\omega\\
9&1&-\omega&0&0&0&0\\
10&0&1&0&0&0&-1
\end{array}
\]
Define
\[
 s_i=1-\frac{2v_i\overline v_i^{\,t}}{\langle v_i,v_i\rangle},\qquad
 (t_1,\ldots,t_{20})=(s_1,s_1,\ldots,s_{10},s_{10}).
\]
Each root has norm \(2\) or \(6\), and each \(s_i\) is an involutive
reflection.  Consequently \(t_1\cdots t_{20}=1\).

\textbf{Claim.} The \(s_i\) generate \(G_{34}\), of order \(39\,191\,040\).
Their natural representation is absolutely irreducible over \(K\), has
character field \(K\), and its rational simple block is \(M_6(K)\).

\begin{proof}
Consider the following set \(R\) of norm-two vectors:
\[
\begin{split}
 R_{\rm sp}&=\{\omega^ae_i-\omega^be_j:
          i\ne j,\ a,b\in\{0,1,2\}\},\\
 R_{\rm de}&=\left\{\frac{\pm1}{1+2\omega}
       (\omega^{a_1},\ldots,\omega^{a_6}):
       a_i\in\{0,1,2\},\ \sum a_i\equiv0\pmod3\right\}.
\end{split}
\]
They have cardinalities \(270\) and \(486\), respectively.  The exact
reconstruction checks that every \(s_i\) permutes \(R=R_{\rm sp}\cup R_{\rm de}\),
that their generated action is transitive on its \(756\) elements, and that
the permutation group has order \(39\,191\,040\).  Its stabilizer-chain
orbit lengths are
\[
                   756,\quad80,\quad27,\quad12,\quad2,
\]
whose product is the stated order.  A base for this stabilizer chain is
$e_1-e_2,e_1-e_3,e_1-e_4,e_1-e_5,e_1-\omega e_2$.
The permutation action is faithful because
\(R\) spans \(K^6\).

Davis \cite[Theorem 5.1 and its proof, pp.~312--313]{Davis08} describes
\(G_{34}\) by the reflections in the hyperplanes
\(x_i-x_j=0\), \(x_1-\omega x_2=0\), and
\(x_1+\cdots+x_6=0\).  He gives the invariant degrees
\(6,12,18,24,30,42\), whose product is \(39\,191\,040\), and the same
\(270+486\) root set.  One root reflection in our group, together with
transitivity on \(R\), gives the reflections in every root direction.  In
particular it gives Davis's standard generators.  Thus the published
\(G_{34}\) is a subgroup of the computed group; equality of the exact orders
proves equality of the groups.  This is an identification by containment and
order, not just by a remembered order and dimension.

The determinant of the six root rows with indices $1,2,3,4,5,7$ is
$-3-6\omega$, so the roots have rank six.  The first root has nonzero
inner product with each of the other nine, so their nonorthogonality
graph is connected.  If a complex invariant subspace \(W\) is given, then
\(W^\perp\) is invariant.  Since \(1-s_i\) has rank one, its root line
lies in \(W\) or \(W^\perp\).  Nonorthogonal roots cannot belong to opposite
summands.  Connectedness puts all roots in the same summand, and their span
is all of \(\mathbb C^6\).  This proves absolute irreducibility.

All character values lie in $K$, while direct multiplication gives
$\operatorname{tr}(s_1s_2s_4)=2-\omega$.  The character field is therefore
exactly $K$.  The displayed realization
over that field has Schur index one, giving the block \(M_6(K)\).
\end{proof}

Riemann existence now supplies a connected \(G_{34}\)-cover
\(C\to\mathbb P^1\) with this product-one generating tuple.  The corresponding
Hurwitz component is a finite cover of the ordered configuration space
\(M_{0,20}\), hence has dimension \(17\).  Its genus is
\[
 2g(C)-2=|G_{34}|\left(-2+20\left(1-\frac12\right)\right)
        =8|G_{34}|,
 \qquad g(C)=156\,764\,161.
\]
For the natural character \(\sigma\) and its conjugate, every inertia
eigenvalue list is \((1,1,1,1,1,-1)\).  Chevalley--Weil in the precise form
of \cite[Theorem~2.10, arXiv version p.~8]{FredianiGhigiPenegini15} gives holomorphic multiplicity
\[
                         -6+20\cdot\frac12=4
\]
for each embedding.  A rational primitive idempotent of \(M_6(K)\) cuts out
an abelian variety \(A\) of dimension eight and Weil signature \((4,4)\).
The parabolic multiplicity space has \(K\)-dimension eight.  The following
calculation verifies the latter dimension independently at the cocycle level.

\subsection{Parabolic cocycles and boundary monodromy}

Use left cocycles \(u(gh)=u(g)+g u(h)\).  Write
\[
 u_i=u(\gamma_i)=\lambda_i v_{\lceil i/2\rceil},\qquad
 Q_i=t_1\cdots t_{i-1}.
\]
The parabolic relation and the coboundaries are
\[
 \sum_{i=1}^{20}Q_i\lambda_i v_{\lceil i/2\rceil}=0,
 \qquad u_i=(1-t_i)a.
\]
In \(\lambda\)-coordinates the coboundary matrix has row
\(2\overline v_{\lceil i/2\rceil}^{\,t}/
\langle v_{\lceil i/2\rceil},v_{\lceil i/2\rceil}\rangle\).
The relation matrix and coboundary matrix have ranks six and six, and their
product is zero.  Thus the quotient dimension is \(20-6-6=8\).

For \(j=1,\ldots,4\), let \(b_j\) have
\(\lambda_{2j-1}=\lambda_{2j}=1\) and all other entries zero.  The last four
vectors are supported at the ordered indices
\[
                       (1,3,5,7,10,12,14,16,18,20)
\]
and have coefficient rows
\[
\begin{aligned}
 b_5:&(1,0,0,0,1+\omega,1,0,-1-\omega,0,-1+\omega),\\
 b_6:&(0,1,0,0,1+\omega,\tfrac23-\tfrac23\omega,-1,
                     -\tfrac23-\tfrac13\omega,0,1+2\omega),\\
 b_7:&(0,0,1,0,-\tfrac13-\tfrac23\omega,-\tfrac23-\tfrac13\omega,
                          1+\omega,0,1+\omega,-\omega),\\
 b_8:&(0,0,0,1,1+\omega,1,-2-\omega,-1-\omega,0,3\omega).
\end{aligned}
\]
All eight satisfy the relation.  Adjoining the six coboundary columns gives
rank fourteen, so these are a quotient basis, rather than just candidate
representatives.  Put \(I=\langle b_1,b_2,b_3,b_4\rangle_K\).

The intervals and return powers are as follows:
\[
\begin{array}{c|rrr|r@{\qquad}c|rrr|r}
 j&a&b&k&\operatorname{rank}N_j&j&a&b&k&\operatorname{rank}N_j\\\hline
1&1&8&1&4&10&9&12&1&2\\
2&1&18&1&1&11&9&16&1&4\\
3&2&8&2&3&12&9&17&2&4\\
4&3&8&1&3&13&9&18&1&4\\
5&4&8&2&2&14&11&12&1&1\\
6&5&6&1&1&15&13&14&1&1\\
7&5&8&1&2&16&13&16&1&2\\
8&7&8&1&1&17&15&16&1&1\\
9&9&10&1&1&&&&&
\end{array}
\]
They are pairwise nested or disjoint and each gives a stable partition.
They therefore determine seventeen disjoint nonisotopic essential circles
on the twenty-punctured sphere, a maximal degeneration in
\(\overline M_{0,20}\).  For an interval \([a,b]\), put
\(c=t_a\cdots t_b\).  The listed \(k\) satisfies \(c^k=1\).

The \(k\)-th boundary twist fixes the tuple exactly.  Its pullback is computed
from
\[
\begin{split}
 u(c^k)&=(1+c+\cdots+c^{k-1})
             \sum_{i=a}^b(t_a\cdots t_{i-1})u_i,\\
 u'_i&=u_i+(1-t_i)u(c^k)\quad(a\le i\le b),\qquad
 u'_i=u_i\quad(i\notin[a,b]).
\end{split}
\]
Indeed the twist conjugates each inside loop by the boundary word
\(\gamma_a\cdots\gamma_b\); the cocycle identity gives the formula, and
\(c^k=1\) gives the returning action.  Reducing the resulting cocycles
modulo the verified coboundary space gives matrices
\[
 T_j=1+N_j,\qquad N_j=
 \begin{pmatrix}0&X_j\\0&0\end{pmatrix},\qquad
 N_iN_j=0.
\]
The ranks are in the table.  The cocycle formulas determine these matrices from the roots.

Let $D=\operatorname{diag}(-3,-1,-1,-3)$.  The matrices $H_j=-DX_j$
are Hermitian positive semidefinite.  Here is a short exact certificate
that can be checked from the matrices printed below.  Start with
$R=H_j$; for $p=1,\ldots,4$, set $d_p=R_{pp}$.  If $d_p=0$, its
whole row and column are zero.  Otherwise replace
$R$ by $R-d_p^{-1}R_{\cdot p}R_{p\cdot}$.  All final remainders are
zero, and the successive pivots are the following nonnegative rationals:
\[
\begin{array}{c|rrrr@{\qquad}c|rrrr}
j&d_1&d_2&d_3&d_4&j&d_1&d_2&d_3&d_4\\\hline
1&6&3/2&4/3&9/4&10&9&1&0&0\\
2&6&0&0&0&11&15&9/5&4/3&5/4\\
3&0&3&8/3&9/2&12&21&2&16/7&9/4\\
4&0&2&2&3&13&15&9/5&13/9&24/13\\
5&0&0&4&6&14&6&0&0&0\\
6&0&0&2&0&15&0&2&0&0\\
7&0&0&2&9/2&16&6&3/2&0&0\\
8&0&0&0&6&17&6&0&0&0\\
9&6&0&0&0&&&&&
\end{array}
\]
This expresses each $H_j$ as a sum of positive multiples of Hermitian
rank-one matrices.  The leading principal minors of their sum are
\[
                         96,\quad2379,\quad58722,\quad2344041.
\]
They are positive.  Every positive weighted sum of the \(N_j\) consequently
has rank four, with kernel and image both \(I\).

Flatten the sixteen matrices with indices
\[
                    1,2,3,4,5,6,7,9,10,11,12,13,14,15,16,17
\]
row by row into vectors of length sixteen over \(K\).  Their determinant is
\[
                                  1728\ne0.
\]
Thus the \(K\)-span of all \(X_j\) is \(M_4(K)\), and their complex span
under either embedding is \(M_4(\mathbb C)\).

\paragraph{Explicit output matrices.}
For compactness write $Y_j=3X_j$.  Their entries are in
$\mathbf Z[\omega]$ and are as follows.  Thus the cocycle recipe above,
ordinary exact row reduction over $K$, and the following matrices
provide independent ways to check every finite-dimensional assertion.
\[
\begin{gathered}
Y_{1}=\begin{pmatrix}6&1-\omega&1+2\omega&-3\omega\\
6+3\omega&6&0&6+3\omega\\
-3-6\omega&0&6&-3+3\omega\\
3+3\omega&1-\omega&-2-\omega&6\end{pmatrix},\\[1ex]
Y_{2}=\begin{pmatrix}6&-6\omega&-2+2\omega&6-6\omega\\
18+18\omega&18&-12-6\omega&36+18\omega\\
-12-6\omega&-6+6\omega&6&-18\\
12+6\omega&6-6\omega&-6&18\end{pmatrix} .
\end{gathered}
\]
\[
\begin{gathered}
Y_{3}=\begin{pmatrix}0&0&0&0\\
0&9&-3\omega&9+9\omega\\
0&3+3\omega&9&9\omega\\
0&-3\omega&-3-3\omega&9\end{pmatrix},\\[1ex]
Y_{4}=\begin{pmatrix}0&0&0&0\\
0&6&0&6+3\omega\\
0&0&6&-3+3\omega\\
0&1-\omega&-2-\omega&6\end{pmatrix} .
\end{gathered}
\]
\[
\begin{gathered}
Y_{5}=\begin{pmatrix}0&0&0&0\\
0&0&0&0\\
0&0&12&-6+6\omega\\
0&0&-4-2\omega&9\end{pmatrix},\\[1ex]
Y_{6}=\begin{pmatrix}0&0&0&0\\
0&0&0&0\\
0&0&6&0\\
0&0&0&0\end{pmatrix} .
\end{gathered}
\]
\[
\begin{gathered}
Y_{7}=\begin{pmatrix}0&0&0&0\\
0&0&0&0\\
0&0&6&-3+3\omega\\
0&0&-2-\omega&6\end{pmatrix},\\[1ex]
Y_{8}=\begin{pmatrix}0&0&0&0\\
0&0&0&0\\
0&0&0&0\\
0&0&0&6\end{pmatrix} .
\end{gathered}
\]
\[
\begin{gathered}
Y_{9}=\begin{pmatrix}6&6&-4-2\omega&6\\
18&18&-12-6\omega&18\\
-6+6\omega&-6+6\omega&6&-6+6\omega\\
6&6&-4-2\omega&6\end{pmatrix},\\[1ex]
Y_{10}=\begin{pmatrix}9&6-3\omega&-6-3\omega&9\\
27+9\omega&24&-15-12\omega&27+9\omega\\
-9+9\omega&-3+12\omega&9&-9+9\omega\\
9&6-3\omega&-6-3\omega&9\end{pmatrix} .
\end{gathered}
\]
\[
\begin{gathered}
Y_{11}=\begin{pmatrix}15&4-7\omega&-4+\omega&9-6\omega\\
33+21\omega&24&-15-9\omega&33+12\omega\\
-15-3\omega&-6+9\omega&12&-15+6\omega\\
15+6\omega&7-4\omega&-7-2\omega&15\end{pmatrix},\\[1ex]
Y_{12}=\begin{pmatrix}21&2-11\omega&-2+5\omega&9-12\omega\\
39+33\omega&27&-15-6\omega&39+15\omega\\
-21-15\omega&-9+6\omega&15&-21+3\omega\\
21+12\omega&8-5\omega&-8-\omega&21\end{pmatrix} .
\end{gathered}
\]
\[
\begin{gathered}
Y_{13}=\begin{pmatrix}15&4-7\omega&-1+4\omega&9-6\omega\\
33+21\omega&24&-9&33+12\omega\\
-15-12\omega&-9&9&-15-3\omega\\
15+6\omega&7-4\omega&-4+\omega&15\end{pmatrix},\\[1ex]
Y_{14}=\begin{pmatrix}6&4-4\omega&-4-2\omega&6\\
24+12\omega&24&-12-12\omega&24+12\omega\\
-6+6\omega&12\omega&6&-6+6\omega\\
6&4-4\omega&-4-2\omega&6\end{pmatrix} .
\end{gathered}
\]
\[
\begin{gathered}
Y_{15}=\begin{pmatrix}0&0&0&0\\
0&6&-6-6\omega&12+6\omega\\
0&6\omega&6&-6+6\omega\\
0&2-2\omega&-4-2\omega&6\end{pmatrix},\\[1ex]
Y_{16}=\begin{pmatrix}6&-3\omega&1+2\omega&3-3\omega\\
9+9\omega&9&-6-3\omega&18+9\omega\\
-3-6\omega&-3+3\omega&6&-9\\
6+3\omega&3-3\omega&-3&9\end{pmatrix} .
\end{gathered}
\]
\[
\begin{gathered}
Y_{17}=\begin{pmatrix}6&2-2\omega&0&6\\
12+6\omega&6&0&12+6\omega\\
0&0&0&0\\
6&2-2\omega&0&6\end{pmatrix} .
\end{gathered}
\]

The permutation calculation uses the root set in the order given by
$i,j,a,b$, followed by the dense vectors in lexicographic exponent order.
For the latter the two signs are taken in the order $-,+$.
Each reflection is evaluated on this finite set, and the displayed
stabilizer chain computes the exact group order.  On the parabolic
space, solve the relation first and then reduce the twist pullbacks
modulo the six coboundary columns.  The selected determinant and the
Hermitian elimination above require only field arithmetic.  The program \emph{Exact reconstruction of the Eisenstein Hurwitz family}, in the \href{\workbenchurl}{workbench}, performs these operations from the printed roots, representatives, and intervals.

\subsection{The period differential}

\begin{proof}[Proof of Theorem~\ref{thm:eisenstein-hurwitz-dominance}]
Pass to a finite marking cover so that the \(G\)-cover and its multiplicity
factor are families.  Near the maximal degeneration above, the stable pointed
sphere has seventeen independent plumbing parameters.  On the punctured
polydisc the Hurwitz cover has finite degree.  Since its fundamental group
is \(\mathbb Z^{17}\), a finite Kummer base change kills its finite
permutation monodromy.  Choose the exponents to be multiples of the listed
return powers.  A further common multiple kills any finite central marking
action.  A base point outside the interval circles and a marking of its lifted
fibre give the cocycle pullback calculated above; equivalently this is the
geometric-monodromy identification in
\cite[Lemma~2.2, p.~7, and Theorem~2.5, p.~12]{DettweilerWewers04}.

Every positive plumbing loop is the positive twist about its vanishing
circle.  The passage from a twist to cohomological pullback may reverse all
signs, but does not independently reverse individual signs.  Therefore the
logarithms of the resulting unipotent monodromies are
\[
                        M_j=\varepsilon c_jN_j,
             \qquad c_j\in\mathbb Z_{>0},\quad\varepsilon\in\{1,-1\},
\]
with one common \(\varepsilon\).  Powers only multiply these logarithms,
since \(N_j^2=0\).

Write the punctured coordinates as \(z_j=e^{2\pi i w_j}\).
Schmid's multivariable nilpotent-orbit theorem
\cite[Theorem (4.12), p.~233]{Schmid73} gives
\[
              F(w)=\exp\!\left(\sum_j w_jM_j\right)\Psi(z),
\]
where \(\Psi\) extends holomorphically to the polydisc with values in the
compact dual.  Let \(F_0=\Psi(0)\).  Restrict to the positive diagonal
\(z_1=\cdots=z_{17}=z\).  Its logarithm is \(M=\sum_jM_j\).
The semidefinite certificate, which is purely a statement about the computed
matrices and does not require identifying the auxiliary form with the
polarization, gives
\[
              \operatorname{rank}_K M=4,\qquad
                         \operatorname{im}M=\ker M=I.
\]
On the underlying rational \(H^1\), the rank is eight.  The monodromy weight
filtration centered at one is
\[
                 W_0=W_1=I_{\mathbb Q},\qquad W_2=V_{\mathbb Q}.
\]
The diagonal restriction has the same limiting filtration \(F_0\).
Schmid's limiting mixed-Hodge-structure theorem
\cite[Theorem (6.16), p.~255]{Schmid73} applies.  Since
\(F_0^0=V_{\mathbb C}\) and \(F_0^2=0\), the weight-zero graded piece is
of type \((0,0)\).  More explicitly, both the filtration and its conjugate
have index-zero term equal to the whole space, so Hodge indices on a pure
graded piece are nonnegative; at weight zero they must both be zero.
Consequently
\[
                     F_0^1\cap I_{\mathbb C}=0.
\]
On either \(K\)-embedding block \(V_\sigma\), dimensions are four and four,
and hence
\[
                         V_\sigma=I_\sigma\oplus F_{0,\sigma}^1.
\]

The tangent space of the compact dual \(\operatorname{Gr}(4,V_\sigma)\)
at \(F_{0,\sigma}^1\) is
\(\operatorname{Hom}(F_{0,\sigma}^1,V_\sigma/F_{0,\sigma}^1)\).
Under the displayed decomposition, the infinitesimal action of \(M_j\)
is its induced map from \(F_{0,\sigma}^1\) to \(I_\sigma\).
Changing from the fixed quotient complement in the finite calculation to
\(F_{0,\sigma}^1\) is a fixed invertible map.  Thus these tangent maps
span the same dimension as the \(c_jX_j\), namely all sixteen dimensions
by the determinant \(1728\).

This limiting rank persists at an interior point.  Differentiate the period
map with respect to \(w_j\), and translate its tangent vectors back by
\(\exp(-\sum w_iM_i)\).  The result is
\[
                   M_j\Psi(z)+2\pi i z_j\frac{\partial\Psi}{\partial z_j}.
\]
The second term tends to zero, and the first to \(M_jF_0\).  In a holomorphic
coordinate chart near \(F_0\), take the determinant of the sixteen columns
selected in the finite certificate.  Its limit is nonzero.  It is therefore
nonzero for some sufficiently small point with all coordinates nonzero.
Translation by the group preserves rank.  The smooth period map consequently
has rank at least sixteen; the unitary period domain has dimension
\(4\cdot4=16\), so the rank is exactly sixteen.

It remains to verify that the component is split using the actual geometric
polarization.  The logarithm \(M_1\) is skew-adjoint for that polarization,
its square is zero, and its image is \(I\).  For any \(u,v\), the pairing
of \(M_1u,M_1v\) is minus the pairing of \(u,M_1^2v\), hence zero.  Thus
\(I\) is an actual \(K\)-rational totally isotropic four-space in the
rank-eight Hermitian multiplicity space.  The form is split and its
determinant class is \(1\).  This reasoning is independent of the auxiliary
cup-form convention.  The period image lies in the required component;
rank sixteen makes it analytically open locally, and its algebraic moduli
map is dominant.  The source-dimension inequality was never used as a
sufficient dominance test.
\end{proof}


\section{An explicit Eisenstein eightfold bound after Schoen and Severn}
\label{sec:eisenstein-eightfolds}
This section bounds the degrees in the cyclic-cover transfer after
Schoen and Severn.  The dominant $G_{34}$ source is proved in
Theorem~\ref{thm:eisenstein-hurwitz-dominance} from its explicit finite
data and published inputs.  The exceptional-divisor repair and the
determinant-character transfer are written out below.  Severn's
manuscript \cite{Severn} is credited for the construction and numerical
data; no geometric theorem from it is assumed here.  This argument
does not use Perry's theorem or a CM Hodge-conjecture manuscript.

Put
\begin{align*}
 G&=39191040,& g&=4G+1,\\
 R&=(2G+1)G!G^G,& L_0&=2gR,\\
 S_0&=G^{2g}(g!)^2,& S_B&=(L_0/2)!S_0^R,\\
 k&=L_0/2+4,&
 B_0&=(40G)^k\binom{k}{4}S_B\,8!G^8.
\end{align*}
These formulas specify integers; they do not require decimal expansion.
The constant is deliberately coarse, so that no estimate for coefficients
in an invariant-theory decomposition enters the proof.

\begin{theorem}\label{thm:eisenstein-eightfolds}
Let $K=\Q(\sqrt{-3})$, let $\mathcal O_f=\Z+f\mathcal O_K$ act
integrally on a polarized abelian variety of Weil type $(A,L_A)$, and
assume that the polarization is compatible with the action. Write
$\Delta=\chi(L_A)$. Then
every split Weil eightfold satisfies
\[
 \delta_W(A,L_A)\le B_0 f^4\Delta.
\]
Every Weil sixfold over $K$, of arbitrary determinant, satisfies
\[
 \delta_W(A,L_A)\le \frac{B_0}{4}
          \frac{f^9\Delta^2}{b},
 \qquad b=[\mathcal O_K H_1(A,\Z):H_1(A,\Z)].
\]
The proof uses the dominant split family of group-algebra factors of
Jacobians established above, and gives effective integral detecting
cycles and all the degree losses from that family to $A$.
\end{theorem}

\subsection{The geometric input and the factor lattice}

Theorem~\ref{thm:eisenstein-hurwitz-dominance} gives the required
Hurwitz family of $G_{34}$-covers of genus $g=4G+1$, with a rank-eight
reflection multiplicity and dominant split period map.  The reflection
group contains the scalar $\omega I$: with the ten reflections
defined above, exact multiplication gives
\[
                 (s_1s_2s_3s_4s_7s_5)^{14}=\omega I_6.
\]
This identity also follows by their faithful permutation action
on the displayed root set.

Let $\chi$ be the character of the reflection representation, and let
$s$ be a reflection. In the rational group algebra put
\[
 \varepsilon=\frac6G\sum_{a\in\mathcal G}
       \operatorname{Tr}_{K/\Q}(\chi(a^{-1}))a,
 \qquad e=\varepsilon(1-s)/2.
\]
The first operator is the central projector for the conjugate pair of
reflection representations. The second has rank one in its matrix
algebra $M_6(K)$. Its image on $J(C)$ is an abelian eightfold, denoted
$A_C$, with the restricted polarization $\theta_C$. Character values
are algebraic integers, so their rational traces are integers and
$Ge\in\Z[\mathcal G]$. The Rosati involution sends $a$ to $a^{-1}$;
both commuting factors of $e$ are self-adjoint. The scalar $\omega I$
preserves $A_C$, giving integral $\mathcal O_K$-action.

Let $\Lambda$ be the unimodular lattice of $J(C)$ and
$\Lambda_C=\Lambda\cap e\Lambda_\Q$. This sublattice is primitive.
Every integral functional on it extends to $\Lambda$, and
self-adjointness identifies its dual lattice with $e\Lambda$.
Consequently
\[
 \Lambda_C^{\#}=e\Lambda,\qquad
 G\Lambda_C^{\#}\subset\Lambda_C.
\]
Every elementary divisor of the restricted polarization divides $G$.
In particular
\begin{equation}\label{eq:eis-source-volume}
 \chi(\theta_C)\mid G^8,
 \qquad \int_{A_C}\theta_C^8\le8!G^8.
\end{equation}

\subsection{A bounded integral relation between determinant characters}

Let $U$ be the full centralizer of $\mathcal G$ in the symplectic group
of $H^1(C,\Q)$. Over $\C$, a nonreal conjugate pair of irreducible
representations gives a general linear multiplicity group, an orthogonal
irreducible representation gives a symplectic multiplicity group, and
a quaternionic one gives an orthogonal multiplicity group. Hence the
character group of $U$ has a free determinant coordinate for each
nonreal pair and a coordinate of order two for each quaternionic
multiplicity group. The full orthogonal groups must be used here.

For a subgroup $H'\le\mathcal G$ and a linear character $\chi'$ of
$H'$, let $a(\chi')$ be the determinant character of its multiplicity
space. Frobenius reciprocity identifies this with the image of
$\operatorname{Ind}_{H'}^{\mathcal G}\chi'$ under the homomorphism
which takes differences of conjugate multiplicities and quaternionic
multiplicities modulo two. Brauer induction \cite[\S10.5, Theorem~20, p.~78]{SerreRepresentations}, followed by this character homomorphism, gives an integral relation
\begin{equation}\label{eq:eis-supplier-relation}
 e_\sigma=\sum_i n_i a(\chi_i),
\end{equation}
where $e_\sigma$ is the character of one target Weil line.

Characters of order one or two contribute zero. Their multiplicities
in conjugate nonreal representations agree; their multiplicities in
a quaternionic representation are even, since their isotypic subspace
inherits a nondegenerate alternating form. A character with zero
cohomological multiplicity also contributes zero. All remaining terms
come from cyclic intermediate covers
\[
 C_i'=C/\ker\chi_i\longrightarrow X_i=C/H_i'
\]
of order $m_i\ge3$. Since the original inertia groups have order two,
every nonzero branch exponent in such a cyclic cover is $m_i/2$.
The sum of the exponents is zero modulo $m_i$, so their number $s_i$ is
even and they can be paired with their negatives. If $q_i=g(X_i)$,
the primitive eigenspaces in $H^1(C_i')$ have common dimension
$h_i=2q_i-2+s_i$. A nonzero retained multiplicity has $h_i>0$, and
$h_i$ is even. Unramified intermediate covers are included.

There is a uniform bound
\begin{equation}\label{eq:eis-relation-mass}
 \sum_i|n_i|\le R.
\end{equation}
To prove it, append columns $2e_j$ for the order-two coordinates and
lift the relation to an integer column lattice. There are at most $G$
rows, and the column entries have absolute value at most $G$. Restrict
to independent rows; this does not change relations between the
columns. If their rank is $r$, choose an invertible integer matrix $B$
of $r$ columns and put $d_0=|\det B|\le r!G^r$. In the subgroup of
$\Z^r/B\Z^r$ generated by all columns, a shortest expression for the
target uses at most $d_0-1$ signed columns: a shortest path has no
repeated vertex. The residual vector has coordinates at most $d_0G$.
Cramer's rule, with cofactors at most $(r-1)!G^{r-1}$, expresses it in
the columns of $B$ with each coefficient at most $r!G^r$. The total
coefficient sum is therefore at most $(r+1)r!G^r\le R$. Removing the
auxiliary parity columns leaves \eqref{eq:eis-supplier-relation} and
cannot increase that sum.

\subsection{Integral cycles from the cyclic intermediate covers}

Fix one term, and write $C'$, $X$, $m$, $h$ for its data. We use
Schoen's construction \cite[Theorem~2.0 and Lemmas~2.3--2.6,
pp.~11--17]{Schoen88}. The primitive rational part has dimension
$\varphi(m)h$, so
\[
 h\le g(C')\le g,
 \qquad m\le G.
\]
Let $G_0$ be the kernel of the sum map on $(\Z/m)^h$, extended by
the permutation group on the $h$ factors. It has order
$m^{h-1}h!$. There are finite quotient maps
\[
 q:C'^h\longrightarrow W_0=C'^h/G_0
       \longrightarrow\operatorname{Sym}^hX,
\]
the second with residual cyclic group $\Z/m$. Schoen chooses a linear
$P_0\simeq\mathbf P^{h/2}$ in a complete linear system on $X$. The
strict transform of $P_0$ is unchanged by a sequence of blowups and
avoids the branch divisors. The cyclic cover therefore splits over it.
A component descends to $Q_0\subset W_0$ birationally over $P_0$,
and its primitive-character components in cohomology are nonzero.

We specify the exceptional-divisor issue in this last assertion.
When $4\mid m$ and all exponents are $m/2$, the prime chosen in
\cite[Sublemma~2.9, p.~19]{Schoen88} does not exist. The replacement
is the faithful-character argument of \cite[Lemma~4.4, pp.~9--10]{Severn}.
At a blowup of the intersection of a paired pair of branch divisors,
their total exponent is zero modulo $m$. The exceptional divisor is a
$\mathbf P^1$-bundle and is not itself a branch divisor. The two
strict transforms meet each fibre in two points. For a faithful
character, the eigensheaf is extension by zero of a rank-one local
system on their complement $\C^\times$, with nontrivial monodromy.
Its compactly supported cohomology vanishes: degree zero is zero,
degree two is dual to zero invariants, and the Euler characteristic is
zero. If a whole fibre lies in another branch divisor, its eigensheaf
is zero because its inertia is nontrivial. Proper base change thus
kills the faithful-character cohomology of the entire exceptional
divisor. The original centre has nontrivial inertia as well and has
zero faithful-character eigensheaf. The relative-cohomology sequence
then identifies the primitive parts before and after the blowup.
Schoen's nonzero normal-bundle intersection on the final component
therefore remains nonzero on $Q_0$. This proves the needed descent for
these exponents, including those omitted by the prime-choice argument.

Let $H$ denote the sum of point classes on $C'^h$. An incidence divisor
$D_x=x+\operatorname{Sym}^{h-1}X$ restricts to the hyperplane class on
$P_0$, and pulls back to $mH$ on $C'^h$. If $Z_0$ is an irreducible
component above $Q_0$ and has degree $e_0$ over $Q_0$, then
\[
 e_0\le |G_0|,
 \qquad \deg_H Z_0=e_0/m^{h/2}.
\]
The integral orbit sum
$\mathcal R_0=\sum_{a\in G_0}a_*Z_0$, with repetitions, has effective
degree at most $|G_0|^2/m^{h/2}$. Rational transfer for a finite
quotient gives $[\mathcal R_0]=e_0q^*[Q_0]$. This cohomological
identity does not require the quotient to be flat along $Q_0$.

For the residual action $\delta$ of $\Z/m$, form the integral signed
cycle
\[
 Y=\sum_{t=0}^{m-1}c_m(t)\delta(t)_*\mathcal R_0,
\]
where $c_m(t)$ is the Ramanujan sum. This operator is $m$ times the
rational primitive projector, and $|c_m(t)|\le\varphi(m)$. It keeps
the nonzero primitive components. In $G_0$-invariants a primitive
residual character forces all coordinate characters to be the same
primitive character. Consequently only
$H^1(C')^{\otimes h}$ occurs, and permutation invariance identifies
its component with the corresponding determinant line; the signs in
the permutation action are the cohomological signs in degree one.
Writing the mass of a signed cycle for the sum of the degrees of its
positive and negative parts, we obtain
\begin{equation}\label{eq:eis-cyclic-mass}
 \operatorname{mass}_H(Y)
 \le m\varphi(m)\frac{(m^{h-1}h!)^2}{m^{h/2}}
 \le m^{3h/2}(h!)^2\le S_0.
\end{equation}
The orbit sum in this estimate includes its integral multiplicities.

\subsection{Transfer by a finite collection of divisors}

Take $|n_i|$ copies of the cycle for each term of
\eqref{eq:eis-supplier-relation}. Their exterior product $Y_B$ lies on
a product $B$ of $\ell=\sum_i|n_i|h_i$ curves, with
$\ell\le gR\le L_0$. For the sum $H_B$ of point classes,
\begin{equation}\label{eq:eis-product-mass}
 \operatorname{mass}_{H_B}(Y_B)
 \le \frac{(\ell/2)!}{\prod_i((h_i/2)!)^{|n_i|}}S_0^R
 \le S_B.
\end{equation}
Every selected determinant component is nonzero. The perfect pairing
between inverse primitive eigenspaces on each curve gives its dual
determinant tensor. Choose the component according to the sign of
$n_i$. The product of the dual tensors with one target Weil eigenline
has degree $\ell+8$ on $B\times A_C$, is $U$-invariant by
\eqref{eq:eis-supplier-relation}, and sends $[Y_B]$ to a nonzero vector
of that line. Denote this tensor by $\Gamma$. Both the source cycles
and their dual tensors have middle degree on the curve products.

We need control of the divisors representing $\Gamma$. Let
$D$ be the image of $\Q[\mathcal G]$ in $\operatorname{End}H^1(C)$.
The semisimple-algebra lemma and the double-centralizer argument in
\cite[Lemma~3.5]{MilneLefschetz} give
$\operatorname{End}_U H^1(C)=D$. The invariant theorem for full
symplectic and orthogonal groups and paired general linear groups
\cite[Proposition~3.6 and Remark~3.7]{MilneLefschetz} says that their
exterior invariants on powers are generated in degree two.
The degree-two tensors correspond to self-adjoint matrices over $D$.
Off-diagonal entries and their adjoints are spanned by group elements;
the diagonal self-adjoint entries are spanned by $a+a^{-1}$.
It follows that these tensors are spanned by the divisor classes of
the globally generated line bundles
\begin{equation}\label{eq:eis-divisor-list}
 \operatorname{pr}_i^*(2\Theta),\qquad
 (\operatorname{pr}_i+a\operatorname{pr}_j)^*(2\Theta),
 \qquad a\in\mathcal G,
\end{equation}
on powers of $J(C)$. Repeated indices are allowed: the identity
$(1+a)^\dagger(1+a)=2+a+a^{-1}$ supplies the diagonal entries.

For a quotient $\pi:C\to C'$, the map
$C'\to J(C')\xrightarrow{\pi^*}J(C)$ identifies its first cohomology
with a rational direct summand, equivariantly for $U$. The idempotent
$e$ gives the corresponding summand for $A_C$. Thus $\Gamma$ lifts
to an exterior invariant on a power of $J(C)$. Its degree is $\ell+8$,
so it is a complex linear combination of monomials of
$c=\ell/2+4$ divisors from \eqref{eq:eis-divisor-list}. Since its
action on $Y_B$ is nonzero, some monomial has nonzero projection to
the target Weil line. No bound for the coefficients of the linear
combination is required.

We now bound every such monomial. If $a_0=\deg\pi$ and
$\eta:C'\to J(C)$ is the preceding map, the norm identity and
Riemann--Hurwitz give
\[
 \deg\eta^*\Theta=a_0g(C')
 \le g(C)-1+a_0\le5G.
\]
On $B\times A_C$, the restricted product polarization is therefore
at most $5G(H_B+\theta_C)$ in the nef order. The theorem of the
square bounds each line bundle in \eqref{eq:eis-divisor-list} by eight
times that product polarization, also for repeated indices. Hence
\[
 D_j|_{B\times A_C}\le40G(H_B+\theta_C)
\]
in the nef order. Successive general members of these globally
generated systems meet every component of $|Y_B|\times A_C$
properly. They give effective integral four-cycles. Their signed
pushforward to $A_C$ has the prescribed correspondence class, and
its degree mass is at most
\begin{align*}
 &(40G)^c\int_{|Y_B|\times A_C}
                    (H_B+\theta_C)^c\theta_C^4\\
 &\qquad=(40G)^c\binom c4
              \operatorname{mass}_{H_B}(Y_B)
              \int_{A_C}\theta_C^8
 \le B_0.
\end{align*}
Only the term with exponents $(\ell/2,4)$ contributes. Components
whose images have smaller dimension push to zero. Since the total
Weil projection is nonzero, one irreducible component of the remaining
effective cycles detects. This proves $\delta_W(A_C,\theta_C)\le B_0$.

\subsection{From the source to every polarization and determinant}

The source has maximal-order Eisenstein action and split Hermitian
space. Since $t_K=1$ and the class number is one,
Theorem~\ref{thm:principal-quotient} gives an isogeny to the principal
$U^4$ family pulling the target polarization back to $\theta_C$.
Pushing forward a detecting cycle costs at most its original degree.
The induced map between rational period domains is an isomorphism,
so the principal family still has an open image with detector bound
$B_0$. The closed bounded-detector locus of
Proposition~\ref{prop:detector-loci} extends this bound to its entire
connected component. Pullback through the principal quotient of an
arbitrary target costs $f^4\Delta$, proving the eightfold estimate.

For a maximal-order sixfold of arbitrary determinant, the balanced CM
surface in Theorem~\ref{thm:explicit-stabilization} has polarization
product $\Delta$. Its product with the sixfold is a split eightfold
of polarization product $\Delta^2$. The opposite-graph contraction
in that theorem has cost $1/4$ in the Eisenstein case. This gives
$B_0\Delta^2/4$. For an order of conductor $f$, saturation has degree
$b$, polarization multiplier $f$, and polarization product
$\Delta'=f^6\Delta/b$. Pulling back a three-cycle costs $b/f^3$.
Thus the bound is
\[
 \frac b{f^3}\frac{B_0}{4}
       \left(\frac{f^6\Delta}{b}\right)^2
 =\frac{B_0}{4}\frac{f^9\Delta^2}{b},
\]
as asserted. The factor $8!$ in $B_0$ makes $B_0/4$ an integer.
\section{A numerical detector bound for Gaussian sixfolds of discriminant \texorpdfstring{$[-3]$}{[-3]}}
\label{sec:gaussian-discriminant-verified}

This section gives a uniform numerical detector bound for the
nonsplit Gaussian sixfold family.  The geometric source and determinant
construction are due to Severn \cite[Theorems~A--C, Sections~3--7
and Appendices~A--B]{SevernGaussian}.  That unrefereed preprint proves
algebraicity using an $E_8$ cover family.  We write out a quantitative
form, including the integral polarization and the degree of the
transferred cycles.  The group identity and the period rank used below
are checked by exact and interval certificates.

Put $K=\mathbf Q(i)$ and $\mathcal O=\mathbf Z[i]$.
For a polarized sixfold $(A,\theta)$ of Weil type define
$\delta_W(A,\theta)$ to be the least degree $\int_Z\theta^3$ of an
irreducible reduced threefold $Z\subset A$ whose rational class has
nonzero projection to $W_K=\bigwedge_K^6H^1(A,\mathbf Q)$.

\subsection{The explicit bound}
Set
\[
 \begin{gathered}
 G_0=1393459200,\qquad g=3077222401,\qquad D=174182400,\\
 T=10^{755},\qquad k=gT+3,\qquad \lambda=8(g+G_0)=35765452808.
 \end{gathered}
\]
For the Fermat construction set
\[
 \begin{gathered}
 M_0=16(4\cdot12^3)^2=764411904,\qquad B_0=12(M_0+10)=9172942968,\\
 E_{2p}=B_0^p2^{p^2},\qquad E_0=1.
 \end{gathered}
\]
For positive even $S$ and $R\ge1$, put $n=S+2(R-1)$, $a=n+S$ and
\begin{equation}\label{eq:gaussian-discriminant-fermat}
 C_{12}(S,R)=
 2^{S+R+(n/2)^2}(a+1)(a!)^{a+2}S!\,
 12^{(R+1)(a+1)+2R}B_0^{n/2}.
\end{equation}
The required source bound is the integer
\begin{equation}\label{eq:gaussian-discriminant-bound}
 \begin{split}
 S_{\rm ord}&=G_0^{2g}(g!)^2,\\
 S_{\rm prod}&=(156341088!)^2\,12^{312682176}C_{12}(609948,45),\\
 S_0&=S_{\rm ord}+S_{\rm prod},\\
 \mathcal B&=\lambda^k\binom{k}{3}(gT)!S_0^T\,720D^6.
 \end{split}
\end{equation}
There is no minimization over unspecified cycles in this expression.
For example its single Fermat-function value is
\[
 \begin{split}
 C_{12}(609948,45)={}&
 2^{93036590317}\,1219985\,(1219984!)^{1219986}\,609948!\\
 &\hspace{1cm}\cdot12^{56119400}\,9172942968^{305018}.
 \end{split}
\]

\begin{theorem}\label{thm:gaussian-volume-bound}
Let $(A,\theta)$ be a complex abelian sixfold with an integral
$\mathcal O=\mathbb Z[i]$-action compatible with its polarization,
of signature $(3,3)$ and Hermitian determinant class $[-3]$.
Put $\Delta=\int_A\theta^6/6!$.  Then $3\mid\Delta$ and
\begin{equation}\label{eq:gaussian-discriminant-target}
              \delta_W(A,\theta)\le\frac{\Delta}{3}\mathcal B.
\end{equation}
In particular every compatible complex sixfold on
$\mathcal O^6$ with form
$H_{\rm ref}=U\perp U\perp\langle1,-3\rangle$,
$U=\left(\begin{smallmatrix}0&1\\1&0\end{smallmatrix}\right)$,
satisfies $\delta_W\le\mathcal B$.
Every sixfold obtained from the fixed norm-two-root projector in
Section~\ref{subsec:gaussian-discriminant-source}, with its restricted
Jacobian polarization, also satisfies $\delta_W\le\mathcal B$.
\end{theorem}

The factor $\Delta/3$ is the degree of an isogeny from the target
to a member of the reference family preserving the polarization by
pullback.  The statement uses the integral Gaussian action and the
complete Hermitian pairing; rational determinant class alone does
not specify an integral polarization volume.

\subsection{The source properties, including conventions}
\label{subsec:gaussian-discriminant-source}

Let $C_8$ be the Cartan matrix with diagonal two and edges
\[
 (1,3),(3,4),(4,5),(5,6),(6,7),(7,8),(2,4),
\]
and set $H_0=C_8^{-1}$.
Use Severn's fundamental-weight matrices $a,b,c$, with
$d=(abc)^{-1}$, and the positively oriented pullback tuple
\begin{equation}\label{eq:gaussian-discriminant-tuple}
 (ia^5,a^{-4}ba^4,a^{-4}ca^4,a^{-3}ba^3,a^{-3}ca^3,
 a^{-2}ba^2,a^{-2}ca^2,a^{-1}ba,a^{-1}ca,b,c,-id^5).
\end{equation}
The matrices are printed in \cite[Appendix~A]{SevernGaussian}.
Starting from the Cartan reflections, the finite certificate constructs the
240 roots and the action on their 480 real and imaginary copies.
This verifies product one, membership and generation of
$G=W(E_8)\cup iW(E_8)$, of order $G_0$.
The matrix/permutation convention is explicitly checked:
column-root action $p(M)$ is antihomomorphic for GAP right
multiplication, so a weight matrix $g$ is represented by $p(g^t)$.
The matrices are recovered from their eight basis-root images.

The local classes are ten reflections, one element with natural
fixed-space codimension four, and one with codimension eight.
Their orders are eleven twos and one twelve.  Thus
$g(C)=1+53G_0/24=g$, the natural parabolic multiplicity is six,
and the two unitary Hodge multiplicities are three.
The oriented tuple is checked directly; no unproved Hurwitz equivalence
to a differently ordered displayed tuple is used.

The Weyl group has 112 irreducible characters.  A
Dixon--Schneider calculation on this permutation group verifies the
entire candidate character table, class fusion, sizes and power maps.
All 112 Frobenius--Schur indicators are one.
For 95 ordinary suppliers and the 12 distinct slots of three product
suppliers, the finite certificate verifies generator membership, subgroup order
and the existence and exact image order of the proposed character,
using a checked homomorphism from generator images.  Actual subgroup
classes then give fusion, fixed-coset and exact return counts, branch
exponents, genus, and all 45 determinant coordinates.
The equality $h=2q-2+s$ is also checked against the full cohomological
character inner product.  In particular, exponent assignments are not
assumed to extend to characters.

The selected rows satisfy
\begin{equation}\label{eq:gaussian-discriminant-identity}
             \sum_{j=1}^{98}a_jv_j=e_2,\qquad
             \sum_j|a_j|<10^{755}=T.
\end{equation}
The second coordinate is the natural representation.
The 95 ordinary rows have paired nonzero branch exponents.
The three product rows have the following exact data, counting repeated
slots:
\[
\begin{array}{c|r|r|r|r}
 \text{row}&R&H&S=\sum(s_j-2)&q\\ \hline
 3231&2&19784&6748&3259\\
 3244&45&156341088&609858&1730347\\
 3269&4&42739160&4816&5341793
\end{array}
\]
In the middle row the six distinct slots have multiplicities
$(1,11,1,20,11,1)$ and branch counts
$(201712,112,201712,112,124,201688)$, respectively.
Thus $\sum s_j=609948$ and $S=609948-2\cdot45=609858$.
This lift differs from the one printed in
\cite[Table~4]{SevernGaussian}; the finite group calculation verifies
the same determinant row and all four simultaneous Hodge balances.
All slots in one row have the same base subgroup.
Their orders are four or twelve, their branch counts exceed one,
and simultaneous conjugation by $1,5,7,11$ gives Hodge dimensions
$(H/2,H/2)$.  Every supplier uses at most $g$ curve factors.
The public witness hash and the exact 755-digit coefficient mass are
recorded in the independent certificate.

For the oriented tuple let $P_j=g_1\cdots g_{j-1}$.
Write a parabolic cocycle as $c_j=(1-g_j)u_j$ with
$\sum_jP_jc_j=0$, and set $F_j=\sum_{l<j}P_lc_l$.
The geometric cup form, in column convention, is
\[
 S(u,v)=\sum_jh_0(P_j^{-1}F_j^u+u_j^u,c_j^v),\qquad H=iS.
\]
It is obtained from Dettweiler--Wewers \cite[Theorem~2.6, p.~158]{DettweilerWewersPoincare}, by reversing
indices and transposing their right-module complex.
The exact calculation gives a 14-dimensional cocycle space,
an eight-dimensional coboundary radical, and, on an integral quotient
complement, the matrix
\begin{equation}\label{eq:gaussian-discriminant-cup}
 H=\begin{pmatrix}
0&i&-1-i&-1&-2&1+i\\
-i&0&-1-2i&-1-i&-2+2i&1\\
-1+i&-1+2i&0&2i&0&i\\
-1&-1+i&-2i&0&-2+2i&1+i\\
-2&-2-2i&0&-2-2i&0&-2-2i\\
1-i&1&-i&1-i&-2+2i&2
 \end{pmatrix}.
\end{equation}
It has determinant $-24$ and signature $(3,3)$.
Since $8=N(2+2i)$ its determinant class is $[-3]$.
A separate rational echelon complement gives determinant $-5/3$,
whose ratio to $-3$ is $N((1+2i)/3)$.

Use the covariant deck action $a_h=(L_h)_*$ on the Jacobian and
$T_h=(L_h^{-1})^*$ on curve cohomology.
For the norm-two root $v=C_8e_1$, matrix-coefficient orthogonality gives
the primitive self-adjoint projector
\[
 e=\frac8{|G|}\sum_{h\in G}
       \operatorname{Re}h_0(v,\rho(h^{-1})v)h,\qquad
 De\in\mathbf Z[G],\qquad e\mathbf Q[G]e\simeq K.
\]
Thus $A=\operatorname{im}(De:J(C)\to J(C))^0$ is an abelian sixfold.
Its integral Gaussian action is $i_A=a_z^{-1}$, $z=iI_8$.
The regular-fibre map $j_x(h)=h_0(v,\rho(h)x)/2$ satisfies
\[
 \sum_h\overline{j_x(h)}j_y(h)=|G|h_0(x,y)/16.
\]
Consequently $\Phi(u)=\operatorname{PD}(j_u+\overline{j_u})$ is
$K$-linear for this action and pulls the restricted Jacobian form back
to $H_s=DH$.  Integral quotient cochains yield the lattice sandwich
\begin{equation}\label{eq:gaussian-discriminant-lattice}
             \mathcal O^6\subseteq\Lambda_A
                              \subseteq H_s^{-1}\mathcal O^6.
\end{equation}
The right containment is obtained by integral pairing with the left
one.  There is also a direct polarization bound:
primitivity of $\Lambda_A$ in the unimodular Jacobian lattice and
self-adjointness of $e$ identify $eH_1(C,\mathbf Z)=\Lambda_A^*$.
Thus $D\Lambda_A^*\subseteq\Lambda_A$, so all six polarization
elementary divisors divide $D$ and
\begin{equation}\label{eq:gaussian-discriminant-volume}
                         \int_A\theta_A^6\le720D^6.
\end{equation}

\subsection{An effective Fermat construction}

For a degree-twelve Fermat surface, Looijenga's integral primitive
group-ring generator \cite[Corollary~2.2]{LooijengaFermat} has pairing $\pm\prod_{j=0}^3(1-u_j)$.
Every pairing of two translates therefore has absolute value at most
16.  The integral character-orbit projector has coefficient mass at
most $4\cdot12^3$.  Its value $D_\alpha$ on the generator is a nonzero
integral primitive $(1,1)$ class, with nonzero projection to every
character in the orbit, and $D_\alpha^2\ge-M_0$.
The integral $(1,1)$ theorem supplies its line bundle.
For $L=D_\alpha+(M_0+10)H$, use $H^2=12$, $K=8H$ and
$\chi(\mathcal O)=166$ to obtain
\[
 \chi(L)=166+\tfrac12\{D_\alpha^2+
              12(M_0+10)^2-96(M_0+10)\}>0.
\]
Since $(K-L)H<0$, Serre duality gives an effective divisor of degree
$B_0$.  One reduced component detects the chosen character.

Shioda's published strong condition $(P_{12})$ decomposes every higher
Hodge character by joins and hash correspondences.  The statement for
$12\le m\le20$ is explicit in \cite[p.~182; Theorems~I--IV]{ShiodaFermat}.
Joins multiply cycle degrees.  The hash map is defined by
\[
 (x,y)\dashrightarrow
 [x_0y_{s+1}:\cdots:x_ry_{s+1}:
   \epsilon x_{r+1}y_0:\cdots:\epsilon x_{r+1}y_s],
 \qquad\epsilon^{12}=-1.
\]
Its codimension-two base locus is resolved by a blow-up; the resolved
hyperplane class is bounded by $H_1+H_2$ on strict transforms.
For a middle cycle of dimension $p$ the degree cost is at most $2^p$.
The exceptional summand has trivial deleted-coordinate character, so
it cannot account for the nontrivial primitive character being
transferred.  Induction, using $2ab\ge a+b$ for $a,b\ge1$, gives
$E_{2p}=B_0^p2^{p^2}$.

For $V=\prod_{j=1}^RF_{12}^{r_j}$, with positive-dimensional factors
and total even dimension $S$, its iterated join lies in
$F_{12}^n$, $n=S+2(R-1)$.
The graph $\Gamma$ of the map from
$\mathbf P_V(\bigoplus_j\mathcal O(-H_j))$ is a closed subvariety of
$X=F_{12}^n\times V$.  With $a=n+S$ and the Segre polarization,
\[
 \deg X\le a!12^{R+1},\qquad
 \deg\Gamma\le2^{S+R-1}S!12^R,\qquad
 \deg(Z_n\times V)\le a!12^R E_n.
\]
The graph estimate follows by the projective-bundle formula and the
coefficientwise inequality between complete symmetric polynomials and
powers of $\sum H_j$.  The transpose join map is nonzero on the
relevant character by the inverse-character Poincare pairing.
Truong's moving-degree lemma \cite[Lemma~2.2]{TruongMoving} bounds each effective part of the
intersection by $a(\deg X)^a$ times the two degrees.
Replacing this by the larger $(a+1)(\deg X)^{a+1}$ and adding both
parts gives
\[
 2(a+1)(a!12^{R+1})^{a+1}
 (2^{S+R-1}S!12^R)(a!12^R E_n)=C_{12}(S,R).
\]
After projection to $V$, a reduced component of one part detects the
character.  It meets the all-coordinate-open locus, since a coordinate
boundary is fixed by a coordinate scalar on which that character is
nontrivial.  Zero-dimensional factors may be supplied by points and
removed; retaining their number in the increasing upper bound is safe.

\subsection{Integral cycles for the supplier rows}

For a connected cyclic cover $D'\to X$ of order $e\ge3$, write
$q=g(X)\ge1$, $s\ge2$ for its branch count, and $h=2q-2+s$.
The finite quotient
\[
 W=D'^h/\bigl(\ker((\mathbf Z/e)^h\to\mathbf Z/e)
                         \rtimes\mathfrak S_h\bigr)
\]
has residual cyclic action.  Its faithful eigencohomology is zero
except in degree $h$, where it is
$\bigwedge^hH^1(D',\mathbf C)_\psi$, a line.
This follows from finite-group invariants and the Koszul sign in
permuting $H^1$ factors.

The eigensheaf on $\operatorname{Sym}^hX$ is extension by zero off
the branch incidence divisors.  Even at a repeated branch coordinate,
inertia fixes a point above it and acts nontrivially through the
faithful character, so its stalk vanishes.
The Abel--Jacobi map is a projective bundle since $h\ge2q$.
Away from $L_0=K_X+\sum b_j$, the complement of its branch
hyperplanes is $(\mathbf C^*)^{s-1}\times\mathbf A^{q-1}$ with
nontrivial rank-one monodromy and zero compactly supported cohomology.
At $L_0$ it is a rank-$q$ vector bundle over
$\{[x_1:\cdots:x_s]:\sum x_j=0,\ x_j\ne0\}$.
Its cyclic cover is the finite quotient of the appropriate Fermat
open determined by its meridian exponents.
For $e\mid12$, it is dominated by $F_{12}^{s-2}$ with character
entries $12\alpha_j/e$ modulo twelve.
Proper base change and Thom identify the determinant line with this
Fermat character, twisted by $(-q)$; the same holds for its inverse.
Coordinate-boundary eigencohomology vanishes.
These assertions also have the elementary affine interpretation for
$s=2$, which is not needed by the selected positive-dimensional slots.

For a balanced product of $R$ slots, put $H=\sum h_j$ and
$S=H-2Rq=\sum(s_j-2)$.
An effective Fermat detector of degree at most $C_{12}(S,R)$ maps to
the finite quotient, pulls through the vector bundles, and closes in
the product of the $W$'s.  Pairing against the inverse character shows
that its closure is still a detector: proper base change, Thom and
the finite projection formula change the pairing by a nonzero scalar.
Lift one component through the curve-power quotient and sum its
distinct translates.  This is an integral effective cycle and its
pushforward is a positive multiple of the closure.

The map to the hyperplane bases is coordinatewise degree twelve.
The image of the Fermat cycle has degree at most $12^{S/2}\deg Z$.
Coning adds $q$ to each multidegree index; its factorial-weight ratio
is at most $(H/2)!/(S/2)!\le H!$.
The multicone thus has degree at most $H!12^H\deg Z$.
The full map from curve powers to symmetric powers is finite flat of
degree $\prod e_j^{h_j}h_j!\le12^HH!$, and its incidence polarization
pulls back to $\sum e_j\eta_j$, which nef-dominates the sum $\eta$ of
point classes.  The distinct-translate cycle is a subcycle of the
flat inverse image.  Its degree is therefore at most
\[
                         (H!)^2\,12^{2H}C_{12}(S,R).
\]
The coordinate-power degree is already charged in this estimate.
The enlarged even argument $609948>609858$ and monotonicity of
\eqref{eq:gaussian-discriminant-fermat} give $S_{\rm prod}$ for all three rows.

For ordinary paired-branch rows, Schoen's construction \cite[Theorem~2.0 and Lemmas~2.3--2.8]{Schoen88} supplies a
splitting $\mathbf P^{h/2}$ in the symmetric power.
The exceptional-divisor descent works for every nonzero paired
exponent, not merely order-two inertia: on an exceptional fibre the
faithful eigensheaf is extension by zero of a nontrivial rank-one
local system on $\mathbf C^*$, whose compactly supported cohomology
vanishes.  Proper base change gives the required global vanishing.
The quotient group has order $m^{h-1}h!$.
An orbit lift with repetitions has degree at most
$m^{2h-2}(h!)^2m^{-h/2}$; applying the integral Ramanujan projector
costs at most $m\varphi(m)\le m^2$.
Thus a signed integral supplier has mass at most
$m^{3h/2}(h!)^2\le S_{\rm ord}$.
Here $h\le g(D')\le g$, since the character and its inverse are
distinct.  No inverse Lefschetz correspondence is used in this
degree estimate.

\subsection{The invariant divisor contraction}

For $V=H^1(C,\mathbf C)$ and $U=\operatorname{Cent}_{\operatorname{Sp}(V)}G$,
the verified indicator calculation gives
\[
 U\simeq\prod_{\tau\ {\rm real}}\operatorname{Sp}(M_\tau)
       \times\prod_{\{\tau,\bar\tau\}\ {\rm nonreal}}
                         \operatorname{GL}(M_\tau),
 \qquad M_\tau=\operatorname{Hom}_G(\tau,V).
\]
There is no omitted orthogonal sign character.
The evaluation isomorphism in Frobenius reciprocity shows that a
subgroup-character determinant has exponent
$\langle\tau|_H,\chi\rangle-\langle\bar\tau|_H,\chi\rangle$
on $\det M_\tau$.  Hence \eqref{eq:gaussian-discriminant-identity} is an identity of
characters of the $U$.

The quotient curves map by
$D'\to J(D')\xrightarrow{p_J^*}J(C)$.
On one-forms this is the trace, whose composite with curve pullback
is multiplication by $\deg p$.
The pullback of a determinant under addition on $h$ curve factors
is its alternating tensor times the nonzero factorial $h!$ and trace
factors; repeated equal-character forms on a curve vanish because
$H^2(D')_{\chi^2}=0$ for character order at least three.
Thus the supplier cycles pair with the required tensors through
these geometric maps.

Use $|a_j|$ copies of each supplier, taking its inverse determinant
line when $a_j>0$ and its determinant line when $a_j<0$.
Rationality gives nonzero evaluation at conjugates.
The external product $Y$ lies on a product of
$\ell=\sum|a_j|h_j<gT$ quotient curves and pairs nontrivially with
a pure tensor $\zeta$ of $U$-character $(\det M_\rho)^{-1}$.
For the corresponding Weil-line generator $x$ of $A$,
$\zeta\boxtimes x$ is invariant.
Pulling through $q=De:J(C)\to A$ and restricting through the inclusion
multiplies this tensor by $D^6$.

The symplectic/general-linear invariant theorem generates these
invariants by degree-two pairings.
Let $B$ be the image of $\mathbf C[G]$ in $\operatorname{End}(V)$.
Its centralizer is semisimple and adjoint-stable.
Milne's Lemma~3.5 \cite{MilneLefschetz} says that its unitary elements generate it, so
double centralizers give $\operatorname{End}_U(V)=B$.
The pairings are consequently spanned by group correspondences,
and therefore by the classes of
\[
 x_i^*(2\Theta_C),\qquad (x_i+g_0x_j)^*(2\Theta_C),\quad g_0\in G.
\]
Repeated indices are included.  These are globally generated divisors.
At least one product of $c=\ell/2+3$ such divisors has nonzero Weil
contraction against $Y$, since their span contracts to the nonzero
class $D^6(\int_Y\zeta)x$.
The coefficients of the spanning combination do not enter the degree
bound: one monomial suffices.

For a quotient curve of degree $t$, its theta-degree is
$tg(D')\le g-1+G_0$.
The positive-form inequality
$\theta(v+w)\le2\theta(v)+2\theta(w)$ gives the common nef bound
$\lambda(\eta+\theta_A)$ for every displayed divisor, including
repeated indices.
Also
$\operatorname{mass}_\eta Y\le(\ell/2)!S_0^{\sum|a_j|}
\le(gT)!S_0^T$.
Intersect $Y\times A$ with general members of the chosen systems and
push the resulting signed integral three-cycle to $A$.
Its mass is at most
\[
 \lambda^c\binom c3\,\operatorname{mass}_\eta(Y)
                   \int_A\theta_A^6\le\mathcal B.
\]
The binomial is the surviving term of
$(\eta+\theta_A)^c\theta_A^3$, with curve degree $\ell/2$ and
abelian degree six.
Some reduced irreducible component of one signed part has nonzero
Weil projection and degree no greater than this mass.
This proves the source bound.

\subsection{The verified rank-nine point}

The interval existence certificate encloses a four-pole system
with monodromy exactly $(a,b,c,d)$.
It verifies the contraction inequality, the minors restoring the
dropped equations, and the infinity exponent labels
\[
 -1,-1,-11/12,-2/3,-7/12,-5/12,-1/3,-1/12.
\]
Its reported bounds are $\eta<1.054\cdot10^{-35}$,
$\kappa<0.00560$, and $\rho<1.054\cdot10^{-34}$, with
$\eta+\kappa\rho<\rho$.
Pullback by $z=s^5$ and the scalar gauge $s^{1/4-3}$ produces
\eqref{eq:gaussian-discriminant-tuple}.

The holomorphic dictionary requires an exact resonance argument.
At every pole of the original system, eigenvalue differences have
absolute value less than one, so the Frobenius operators
$m-\operatorname{ad}(A_0)$ are invertible for every positive integer
$m$.  Their scalar-majorant recurrence converges.
The repeated infinity block is semisimple because its nonresonant
gauge identifies local monodromy with $\exp(2\pi i A_0)$, and finite
monodromy has no nontrivial unipotent part.
After fifth-power pullback, zero is the only resonant pole:
positive integral exponent differences there are at most four,
whereas the off-diagonal inherited gauge starts in order five.
After the Mobius coordinate change used by the certificate, the extra
coordinate-unit power is diagonal.
Thus the resonant off-diagonal coefficients of orders one through four
are exactly zero, as chosen by the program.
At infinity the original differences have denominator twelve, so
multiplication by five introduces no nonzero integral difference.
This proves the normalization; merely having zero in an obstruction
ball would not prove it.

If $Y$ is the resulting fundamental matrix, the vectors of forms
$Y^{-1}r\,ds$ and $Y^tq\,ds$ identify the two multiplicity spaces with
rational vectors satisfying holomorphy.
The converse follows by equivariant descent and the local
power-times-holomorphic gauges.
At a branch point, with $s-p=y^e$ and exponent $\lambda$, a coefficient
of order $b$ upstairs has order $eb-e\lambda+e-1$.
It is holomorphic precisely when $b\ge\lfloor\lambda\rfloor$;
for the dual vector use $\lfloor-\lambda\rfloor$.
At the regular point infinity in the final coordinate, the rational
vector is $O(s^{-2})$.
The finite principal-parts conditions are therefore complete.

The interval calculation gives matrices of sizes $45$ by $48$ and
$117$ by $120$, with certified invertible pivot blocks.
Their Neumann bounds are less than $1.650\cdot10^{-24}$ and
$3.385\cdot10^{-26}$.
The exact kernels have dimension three: this also follows independently
from the unitary Hodge multiplicities, so selected-row solutions
satisfy every exact holomorphy condition.
The products are exactly $q^tr\,ds^2$.
An invariant holomorphic quadratic differential upstairs has at most a
simple pole downstairs, so its residues satisfy the three exact
relations defining $H^0(\mathbf P^1,K^2(B))$, of dimension nine.
Consistency residual balls are checks, not proofs of these identities.

The certified minor in residue columns four through twelve is
\[
 [\,0.10557441259506563\ \pm6.51\cdot10^{-18}\,]
 +i[\,-0.04409991368479836\ \pm3.88\cdot10^{-18}\,],
\]
which excludes zero.
Kodaira--Spencer and Serre duality identify this multiplication map
with the transpose period differential.
On the natural isotypical factor the coevaluation trace introduces
only the nonzero scalar eight; finite-cover trace introduces only the
nonzero degree $|G|$.
The primitive factor is polarized-isogenous to the multiplicity
Hodge structure, so its period rank is nine.
This equals the dimension of the unitary domain, and gives an analytic
open period image.  One such point suffices.

\subsection{The maximal Gaussian lattice of determinant class \texorpdfstring{$[-3]$}{[-3]}}

The integral transport can be made sharper by using a maximal lattice.
Here integrality means integrality of every Hermitian pairing, not only
of the norms.  We use the convention that Hermitian forms are conjugate
linear in the first variable.

\begin{lemma}\label{lem:gaussian-maximal-lattice}
Let $(V,h)$ be a six-dimensional Hermitian space over $K=\mathbb Q(i)$,
with signature $(3,3)$ and determinant class $[-3]$.  Every full
$\mathcal O=\mathbb Z[i]$-lattice $L\subset V$ such that
$h(L,L)\subset\mathcal O$ admits a $K$-linear isometry
\[
 F:(V,h)\longrightarrow(K^6,H_{\rm ref}),\qquad
 F(L)\subset\mathcal O^6,
 \quad H_{\rm ref}=U\perp U\perp\langle1,-3\rangle.
\]
No scalar multiplication of the form or lattice is required.
\end{lemma}

\begin{proof}
Landherr's classification by dimension, determinant norm class and real
signature identifies the rational Hermitian space with the reference
space \cite[pp.~246--247]{GaussianLandherr}.
Enlarge $L$ to a maximal lattice $M$ among those with integral pairings.
Such a lattice exists: every integral overlattice of $L$ lies in its dual
$L^\#$, and $L^\#/L$ is finite.  Each localization $M_p$ is maximal
with the same property.  Otherwise the corresponding submodule of the
$p$-primary part of $M^\#/M$ would give a larger global integral lattice.
We first determine these local lattices and then prove that their genus
contains only one global isometry class.

At a split prime, the two idempotent parts of a lattice are paired
perfectly.  Write them as $A\oplus B$.  Integrality says
$B\subset A^\vee$, and maximality says $B=A^\vee$.  A basis change
on $A$, and its inverse transpose on the other part, identifies every
such pair.  This gives the reference local lattice.

At an odd inert prime $p$, the extension $K_p/\mathbb Q_p$ is
unramified.  An integral Hermitian lattice diagonalizes as
$\langle p^{a_1}u_1,\ldots,p^{a_6}u_6\rangle$, with $a_j\ge0$ and
unit coefficients $u_j\in\mathbb Z_p^\times$.
Indeed a vector whose norm has the least pairing valuation splits off
an orthogonal line; such a vector exists because the residual trace
pairing is nondegenerate.  Induction proves diagonalization.
Every unit of $\mathbb Z_p$ is a norm of a unit of $\mathcal O_p$:
the residue-field norm is onto, and successive lifts follow from
\[
 N(1+p^m x)\equiv1+p^m\operatorname{Tr}(x)\pmod{p^{m+1}}
 \qquad(m\ge1).
\]
Thus the $u_j$ may all be removed by unit changes of basis.

If $a_j\ge2$, dividing that basis vector by $p$ enlarges the lattice
integrally.  If at least two exponents equal one, their residual
Hermitian space has a nonzero isotropic vector: the norm map of the
quadratic finite field is onto.  A primitive lift $v$ then has
$h(v,M_p)\subset p\mathcal O_p$ and $h(v,v)\in p^2\mathbb Z_p$.
Adjoining $v/p$ again enlarges the lattice integrally.
Consequently a maximal lattice has all exponents zero or exactly one
exponent one.  Their determinant valuations have different parity;
the determinant class $[-3]$ selects the second case precisely at
$p=3$.  These are the local classes of the reference lattice.  The
one $3$-modular line cannot be enlarged: a nonzero class in its dual
quotient has norm of valuation $-1$, which cannot be cancelled by
integral coordinates.  This verifies maximality as well as uniqueness.

At $2$, put $R=\mathbb Z_2[i]$ and $\pi=1+i$.
If an integral lattice $N$ is not unimodular, its Gram matrix modulo
$\pi$ has a nonzero radical.  Lift a radical vector to
$v\in N\setminus\pi N$.  Then
\[
 h(v,N)\subset\pi R,\qquad
 h(v,v)\in\pi R\cap\mathbb Q_2=2\mathbb Z_2.
\]
It follows that $N+R(v/\pi)$ is a proper integral overlattice, since
$h(v/\pi,v/\pi)=h(v,v)/2$ is integral.
Every maximal lattice at $2$ is therefore unimodular.

The unit norm group is
\[
 N(R^\times)=1+4\mathbb Z_2.
\]
One inclusion follows by writing a unit as $a+ib$, with one of $a,b$
odd and the other even.  For the converse, a unit congruent to one
modulo eight is a square in $\mathbb Z_2$.  If $u\equiv5\pmod8$,
then $u-4$ is such a square and
$u=N(\sqrt{u-4}+2i)$.
In particular $-3$ is a unit norm, whereas $-1$ and $3$ are not.

Call a lattice even if all its norms lie in $2\mathbb Z_2$.
An even unimodular lattice has a nondegenerate alternating reduction
modulo $\pi$, so it splits off a binary unimodular summand with matrix
\[
 \begin{pmatrix}2a&b\\\bar b&2c\end{pmatrix},
 \qquad a,c\in\mathbb Z_2,\quad b\in R^\times.
\]
Its determinant $\delta=4ac-N(b)$ satisfies $-\delta\equiv1\pmod4$,
so the binary Hermitian space is isotropic.  This follows directly by
completing the square:
\[
 2a\,h((X,Y),(X,Y))
   =N(2aX+bY)+\delta N(Y)
\]
when $a\ne0$; if $a=0$, a basis vector is already isotropic.
Clear denominators of an isotropic vector and make it primitive in the
binary lattice.  Unimodularity gives a partner pairing to one.
Its even norm can be removed by subtracting one half of that norm
times the isotropic vector.  This produces $U$, and the orthogonal
complement is again even and unimodular.  Thus an even unimodular
rank-six lattice is $U^3$, with determinant class $[-1]$.
Since $3$ is not a local norm, it cannot occur in the space of our lemma.

An odd unimodular lattice splits off a line of unit norm.  Continue
until the remaining summand is even, and hence a sum of hyperbolic
planes.  A unit-norm line absorbs such a plane through the integral
identity
\[
 \langle a\rangle\perp U\simeq\langle a,a,-a\rangle
 \qquad(a\in\mathbb Z_2^\times).
\]
For a basis $e,x,y$ with $h(e,e)=a$ and $h(x,y)=1$, use
$e+x,e-ay,e+x-ay$; these are orthogonal and their change of basis has
unit determinant $a$.  All unit coefficients may therefore be
normalized to $1$ or $-1$.  The determinant class $[-3]=[1]$ requires
an even number of negative signs.

These signs can be removed in pairs.  Choose $c=\sqrt{-7}+2i\in R$,
so that $N(c)=-3$, and put
\[
 P=\begin{pmatrix}\pi&c\\-\bar c&\bar\pi\end{pmatrix}.
 \qquad P^*P=-I_2,\qquad\det P=-1.
\]
The square root exists because $-7\equiv1\pmod8$.
Thus $\langle1,1\rangle\simeq\langle-1,-1\rangle$ over $R$.
Every odd unimodular rank-six lattice in this space has one local
isometry class, containing the reference lattice.

We have proved that $M$ lies in the genus of
$L_0=(\mathcal O^6,H_{\rm ref})$.  We give the global step explicitly.
Put $G=U(H_{\rm ref})$, $G^1=SU(H_{\rm ref})$ and
$K_0=\prod_p U((L_0)_p)$.  Choose local isometries
$g_p(L_0)_p=M_p$, equal to the identity at almost all primes, and let
$g=(g_p)\in G(\mathbb A_f)$.
At every prime the determinant group of $U((L_0)_p)$ is the entire
group of norm-one units: multiplication by any such unit on the
orthogonal norm-one line, fixing its complement, realizes that determinant.

The norm-one idele $t=\det g$ defines a fractional ideal $I$ of
$\mathcal O$ with $I\bar I=\mathcal O$.
Since $\mathbb Z[i]$ is a principal ideal domain, $I=a\mathcal O$.
Then $a\bar a$ is a positive rational unit, and so is one.
The rational isometry $\gamma$ multiplying the norm-one line by $a$
has determinant $a$.  The idele $t/a$ is a norm-one unit everywhere;
choose $k\in K_0$ with that determinant.  Hence
\[
 \gamma^{-1}gk^{-1}\in G^1(\mathbb A_f).
\]
Strong approximation for the simply connected group
$G^1$, whose real group is $SU(3,3)$, gives
$G^1(\mathbb A_f)=G^1(\mathbb Q)(K_0\cap G^1(\mathbb A_f))$.
The precise indefinite special-genus consequence is
\cite[Remark~4.8(2), author p.~9]{GaussianKirschmer}; the determinant
and ideal-class steps above explain why there is only one special genus
here.  Thus $g\in G(\mathbb Q)K_0$, and $M$ is globally isometric to
$L_0$.  Restrict this isometry to $L\subset M$ to prove the lemma.
\end{proof}

The pairing hypothesis is essential.  In a hyperbolic plane over
$\mathbb Z_2[i]$, adjoining $e/(1+i)$ preserves integral norms because
$\operatorname{Tr}(R/(1+i))\subset\mathbb Z_2$, but its pairing with
the other basis vector is not integral.  Such a lattice cannot embed
isometrically into an integral-pairing lattice.

\subsection{The reference family and all integral target lattices}

We now combine the integral lattice statement with the bounded source
cycles and the open period image.  This proves the polarization-volume
factor in Theorem~\ref{thm:gaussian-volume-bound}.

\begin{proof}[Completion of the proof of Theorem~\ref{thm:gaussian-volume-bound}]
Let $E$ be the integral alternating polarization form of a Gaussian
sixfold, and let $I$ be the integral action of $i$.  Compatibility gives
$E(Ix,Iy)=E(x,y)$.  The Hermitian form
\[
 h(x,y)=E(x,Iy)+iE(x,y)
\]
is conjugate linear in the first variable and satisfies
$h(L,L)\subset\mathcal O$ on its full homology lattice.
Thus the integrality hypothesis of
Lemma~\ref{lem:gaussian-maximal-lattice} holds for every sixfold in the
theorem, including the source with its restricted Jacobian polarization.

Apply the lemma first to the source lattice $L_s$.  Its isometric
embedding $F_sL_s\subset L_0$ gives, at corresponding complex periods,
a Gaussian-linear isogeny
\[
 f:A_s\longrightarrow A_0,\qquad f^*\theta_0=\theta_s.
\]
Here $A_0$ has lattice $L_0$ and the reference polarization.  If $Z$
is a source detector, the projection formula gives
\[
 \deg_{\theta_0}(f_*Z)=\int_Zf^*\theta_0^3
                     =\deg_{\theta_s}Z\le\mathcal B.
\]
The Gaussian-linear isogeny identifies the two rational Weil planes,
so this pushforward has nonzero Weil projection.  Its reduced image
is a detector of no greater degree.

The rank-nine period differential proved above gives an analytic open
set in the connected reference unitary domain.  We explain why the
same bound holds at every reference period.  Pass to a fine level
moduli space for this fixed lattice and polarization.  For a fixed
multiple of the polarization giving a projective embedding, the
relative Chow spaces of effective three-cycles of $\theta_0$-degree at most
$\mathcal B$ are proper and of finite type.  Only finitely many degrees
and connected components occur.  The class of a family of cycles is
locally constant in the pulled-back rational cohomology local system;
the projector to the Weil plane is flat because it is determined by
the Gaussian action.  Nonzero Weil projection therefore selects whole
connected components of these Chow spaces.  Their images are closed
by properness.  Their finite union contains the analytic open source
image and hence the entire irreducible moduli component.  Equivalently,
a bounded family of detectors specializes with its nonzero Weil class.
A reduced irreducible component of the resulting effective cycle is
still a detector and has no greater degree.  The connected period
domain parametrizes every compatible complex structure on the
reference lattice.  Consequently
\[
                  \delta_W(A_0,\theta_0)\le\mathcal B
\]
for all of them.

Now let $(A,\theta)$ be any target in the theorem, with homology
lattice $L$.  The lemma gives an isometric embedding $FL\subset L_0$.
Choose a Gaussian basis of $L$ and write the embedding as a matrix $B$.
Then
\[
 H_L=B^*H_{\rm ref}B,\qquad
 [L_0:FL]_{\mathbf Z}=N_{K/\mathbf Q}(\det B)
                    =\frac{|\det H_L|}{3}.
\]
For our normalization $E=\operatorname{Im}h$, the determinant of the
underlying alternating integral form is $(\det H_L)^2$.
The polarization elementary divisors therefore give
\[
 \Delta=\frac1{6!}\int_A\theta^6
       =|\det H_L|.
\]
Thus $j=[L_0:FL]_{\mathbf Z}=\Delta/3$ is a positive integer.
The resulting isogeny has the direction
\[
 \psi:A\longrightarrow A_0,\qquad
 \deg\psi=j,\qquad\psi^*\theta_0=\theta.
\]
Take a reference detector $Z_0$ at the transported period.
Finite flat pullback gives an effective cycle with nonzero Weil
projection, since $\psi^*$ is an isomorphism on that plane.  Its degree is
\[
 \int_{\psi^*Z_0}\theta^3
  =\int_A\psi^*([Z_0]\theta_0^3)
  =j\int_{Z_0}\theta_0^3
  \le\frac{\Delta}{3}\mathcal B.
\]
At least one reduced irreducible component has nonzero projection;
its degree is no greater.  This proves the asserted bound for every
target and every period.  The earlier construction proves the source
bound $\mathcal B$ directly, regardless of its polarization volume.
\end{proof}

For the reference lattice, $\Delta=3$.  The bound is therefore
$\mathcal B$.  More generally the integer $\Delta/3$, which is
read directly from the specified integral Gram matrix, gives the
transport cost.  There is no uncomputed isogeny denominator in the
bound.

\paragraph{Reproducible finite calculations.}
The Gaussian sixfold verifier, available through the workbench,
reconstructs the group, subgroup characters, cup pairing, lattice
transport and determinant identity, and records the interval checks.
The geometric uses of these calculations, including the resonance
normalization and the passage from source to target, are proved above.

\section{An exact obstruction-space criterion for a pure Weil character}
\label{sec:pure-weil-complexes}
This section turns strong semiregularity for a pure Weil character into an exact dimension test.

Let $A$ have dimension $g=2n\ge4$, let $K$ act with signature $(n,n)$, and write $W_K\subset H^{2n}(A,\Q)$ for its Weil plane. For a perfect complex $E$, its Hochschild semiregularity map is
\[
 \sigma_E^2:\operatorname{Ext}^2(E,E)\longrightarrow HH_{-2}(A).
\]
If a finite group $G\subset A\times A^\vee$ acts coherently by translations and tensor products with degree-zero line bundles, and $E$ is $G$-equivariant, call $E$ \emph{strongly $G$-semiregular} when this map is injective on the invariant self-Ext space. The coherence and the equivariant structure are part of the hypothesis. A finite group of other automorphisms is not covered by this definition here.

\begin{theorem}[The exact dimension test]\label{thm:pure-weil-ext}
Suppose that $B\in H^{1,1}(A)\cap H^2(A,\Q)$ and
\[
 e^B\operatorname{ch}(E)=w\in W_K\setminus\{0\}.
\]
Then, for the preceding action and equivariant structure,
\[
 E\text{ is strongly }G\text{-semiregular}
 \quad\Longleftrightarrow\quad
 \dim_{\C}\operatorname{Ext}^2(E,E)^G
     =2\binom{2n}{2}=2n(2n-1).
\]
The assertion includes the trivial group. Every such equivariant object has invariant second self-Ext dimension at least $2n(2n-1)$. Neither simplicity nor vanishing of negative self-Ext is needed for this equivalence. The latter vanishing is a separate hypothesis for the deformation theorem.
\end{theorem}

\begin{proof}
Put $T=H^0(A,T_A)$, $U_0=H^1(A,\mathcal O_A)$ and $H=U_0\oplus T$. The tangent bundle is trivial and the Todd class is one. Multiplicative Hochschild--Kostant--Rosenberg identifies
\[
 HH^*(A)=\bigwedge^*H,\qquad \dim H=2g,
\]
where $(\alpha,\xi)$ acts on Hodge cohomology by $\alpha\wedge-+\iota_\xi$. The characteristic maps
\[
 \operatorname{ev}_E^j:HH^j(A)\longrightarrow\operatorname{Ext}^j(E,E)
\]
form a graded algebra homomorphism, and the trace identity gives
\begin{equation}\label{eq:pure-weil-trace}
 \sigma_E^j\operatorname{ev}_E^j(x)=x\cdot\operatorname{ch}(E).
\end{equation}
We use the full Hochschild action, including the bivector and gerbe summands; the relevant Atiyah--Chern identities are \cite[Proposition~6.2.1, Corollary~6.3.2 and Theorem~6.4.2]{BFHH}.

Choose a nonzero holomorphic volume form. Cap product with it identifies $HH^j(A)$ with $HH_{g-j}(A)$, so the character corresponds to an element of $\bigwedge^gH$. Serre duality pairs $\operatorname{Ext}^2(E,E)$ perfectly with $\operatorname{Ext}^{g-2}(E,E)$. Under this pairing and the Hochschild trace pairing, $\sigma_E^2$ is adjoint to $\operatorname{ev}_E^{g-2}$. Indeed pairing the categorical trace of an endomorphism $u$ with a Hochschild operation $x$ is the Serre trace of $u\operatorname{ev}_E(x)$. The categorical adjunction is \cite[Proposition~11, p.~83; \S7.3, pp.~88--89]{CaldararuWillerton}. Equivalently, the same identity follows term by term from the expression $\operatorname{tr}(u\operatorname{At}(E)^j/j!)$.

Translations and degree-zero tensor products act trivially on Hochschild cohomology and preserve the volume and Serre trace. Evaluation therefore lands in invariant self-Ext. Averaging a dual vector shows that Serre duality restricts to a perfect pairing on invariants. If $R^j=\operatorname{Ext}^j(E,E)^G$, the adjunction becomes
\begin{equation}\label{eq:pure-weil-dual}
 \sigma_E^2|_{R^2}\text{ injective}
 \quad\Longleftrightarrow\quad
 HH^{g-2}(A)\longrightarrow R^{g-2}\text{ surjective}.
\end{equation}

The rational $B$-field is handled without choosing a line bundle representing it. Define the invertible shear
\[
 T_B(\alpha,\xi)=(\alpha+\iota_\xi B,\xi).
\]
The type $(1,1)$ of $B$ places the first component in $U_0$. Directly,
\[
 x\cdot(e^Bz)=e^B(T_Bx\cdot z).
\]
Thus $\operatorname{ev}_{E,B}=\operatorname{ev}_E\circ\bigwedge T_B$ is still an algebra homomorphism, $\sigma_{E,B}=e^B\sigma_E$ has the same injectivity, and their composite is cap product with $w$. Surjectivity in \eqref{eq:pure-weil-dual} is unchanged by this shear.

Write $H^{1,0}(A)=X\oplus Y$ for the two $K$-eigenspaces, each of dimension $n$. Then
\[
 H^1_\sigma=X\oplus\overline Y,
 \qquad H^1_{\bar\sigma}=Y\oplus\overline X.
\]
The rational nonzero class $w$ has nonzero coefficients on both conjugate determinant lines. Under the volume identification it becomes
\[
 \eta=a\operatorname{vol}_{H_1}+b\operatorname{vol}_{H_2},
 \quad ab\ne0,\qquad
 H_1=Y^*\oplus\overline Y,\quad H_2=X^*\oplus\overline X.
\]
Here $H=H_1\oplus H_2$ and $\dim H_i=g$. Multiplication by $\eta$ on $\bigwedge^2H$ has kernel exactly $H_1\otimes H_2$ and rank $2\binom g2$. The mixed products kill both top forms. The pure second exterior powers multiply injectively by the opposite top form, with distinct bidegrees since $g\ge4$.

The composite in \eqref{eq:pure-weil-trace}, after the shear, consequently has rank $2\binom g2$. This proves the lower bound on $\dim R^2$ and proves sufficiency when equality holds.

For necessity suppose that the semiregularity map on $R^2$ is injective. Every element of its quadratic annihilator $H_1\otimes H_2$ evaluates to zero, by \eqref{eq:pure-weil-trace}. Evaluation is an algebra homomorphism, so it kills the exterior ideal generated by these mixed products. Its degree-$g-2$ image is therefore a quotient of
\[
 \bigwedge^{g-2}H_1\oplus\bigwedge^{g-2}H_2,
\]
of dimension $2\binom g2$. Equation~\eqref{eq:pure-weil-dual} says that this image is all of $R^{g-2}$. Invariant Serre duality gives $\dim R^2=\dim R^{g-2}\le2\binom g2$, and the reverse inequality was already proved. This proves necessity and the theorem.
\end{proof}

Strong semiregularity implies the weak equivariant condition in Perry's theorem: the forgetful semiregularity square commutes, so an element killed by the equivariant trace is also killed by the injective ordinary trace on invariants. No general equality of the weak and strong definitions is needed here. For a pure Weil character on an eightfold, the exact strong criterion is $\dim\operatorname{Ext}^2(E,E)^G=56$.


\subsection{Comparison with semiregularity and analytic deformation results}

Bhattacharjee--Bhattacharya give the untwisted, nonequivariant numerical
criterion in \cite[Theorem~6]{BBSecantObjects}; their Theorem~10 also
asserts nonexistence of pure semiregular perfect complexes. These are
logically separate assertions: the dimension criterion does not assert
that equality is attained. The form proved here includes a rational
$B$-field and invariants for a coherent finite translation--tensor action.

An analytic nonexistence argument requires actual perfect objects on
nearby unpolarized tori. Buchweitz--Flenner define semiregularity for
perfect complexes, but their analytic propagation theorems are stated
for coherent sheaves \cite[Definition~4.1, pp.~165--166;
Theorems~5.1 and~7.3, pp.~174--179 and~189]{BF}.
Pridham proves all-obstruction semiregularity in the algebraic derived
setting \cite[Corollaries~2.24--2.25, pp.~18--20]{Pridham}.
Wei constructs convergent Kuranishi theory for cohesive modules on
a fixed compact complex manifold \cite[Theorem~3.11 and \S5]{WeiCohesive},
while Chan--Suen treat simultaneous deformations of a manifold and a
vector bundle \cite[\S\S3,5--6]{ChanSuen}.
These precise statements do not by themselves identify a convergent
relative deformation theory for arbitrary perfect complexes on varying
tori. Goldman--Millson's formal comparison
\cite[\S1.1, p.~49; \S\S2--3]{GoldmanMillson} likewise does not state
that particular relative theorem. A relative cohesive-module argument
and its comparison with formal perfect deformations would supply the
additional bridge. We do not infer nonexistence in the twisted
equivariant case from the untwisted assertion. The companion gives the
coherent-sheaf exclusion and explains why global strict vector-bundle
resolutions cannot be assumed on nonprojective fibres.

\section{Detecting components of complete intersections}
\label{sec:complete-intersections}
This section makes the geometric question concrete: one may ask for a detecting component in one fixed multiple of the polarization.

Let $(A,L)$ be of Weil type of dimension $2n$, with compatible polarization $\theta=c_1(L)$, and put $\Delta=\chi(L)$. Define $m_W(A,L)$ to be the least positive integer $m$ such that $n$ effective divisors in $|mL|$ meet properly and their intersection cycle has an irreducible detecting component. Set $m_W=+\infty$ when no such intersection exists.

\begin{proposition}\label{prop:complete-intersections}
The numbers $m_W$ and $\delta_W$ are finite simultaneously. When finite,
\[
 \delta_W\leq\frac{(2n)!\Delta}{2}\,m_W^n,
 \qquad m_W\leq3^{n+1}\delta_W.
\]
Consequently a uniform detector-degree bound in a fixed polarized family is equivalent to a uniform bound on $m_W$.
\end{proposition}

\begin{proof}
Compatibility of the polarization gives $e_W(\theta^n)=0$: with respect to the two $K$-eigenspaces, $\theta^n$ has exterior bidegree $(n,n)$, whereas the Weil lines have bidegrees $(2n,0)$ and $(0,2n)$. If a proper intersection of $n$ members of $|mL|$ has one detecting component, it must have at least two distinct detecting components. Indeed its total projected class is zero, and one nonzero projected component with positive multiplicity cannot sum to zero by itself. The total degree is $m^n(2n)!\Delta$. The smaller detecting component has degree at most half this number, proving the first inequality.

For the converse we give an elementary bound on equations. If an integral $r$-fold $Y\subset\mathbf P^N$ has degree $d$, then degree-$d$ hypersurfaces through $Y$ have set-theoretic base locus exactly $Y$. Fix $x\notin Y$. There is a linear $\mathbf P^{N-r-1}$ through $x$ disjoint from $Y$: project $Y$ from $x$ to $\mathbf P^{N-1}$, take a general $\mathbf P^{N-r-2}$ avoiding its image, and take the cone with vertex $x$. Choose a hyperplane $K$ in this linear space not containing $x$. Projection from $K$ to $\mathbf P^{r+1}$ is a morphism on $Y$ with finite fibres, since a positive-dimensional projective fibre would meet $K$. Its image is an integral hypersurface $Y'$ not containing the image of $x$. If its degree is $d'$ and the finite map has degree $e$, then $d=ed'$, so $d'\leq d$. Pull back its equation and multiply by a power of a linear form nonzero at $x$ to obtain a degree-$d$ equation vanishing on $Y$ but not at $x$. When $N=r+1$ use the equation of $Y$ itself. This proves the base-locus assertion; it does not assert generation of the ideal sheaf.

The bundle $3L$ is very ample \cite[Remark~7.3, p.~12]{MilneAV}. Embed $A$ by it and choose a detector $Y$ of degree $\delta_W$. Its embedded degree is $d=3^n\delta_W$. Degree-$d$ hypersurfaces through $Y$ restrict to a linear system on $A$ with base locus exactly $Y$. Choose $n$ members successively so that the next contains no irreducible component of the previous intersection. This is possible until the $n$th cut: the preceding components have dimension greater than $n$ and cannot be contained in $Y$. There are finitely many components, and avoiding each is the complement of a proper linear subspace of the space of sections.

More precisely, the successive intersections are local complete intersections in the smooth variety $A$, hence Cohen--Macaulay and without embedded associated primes. Avoiding their components makes the next equation a nonzerodivisor and preserves the expected pure dimension. The final intersection has pure dimension $n$ and contains $Y$, so $Y$ occurs in its fundamental cycle with positive multiplicity. Each divisor belongs to $|3dL|$. Thus $m_W\leq3d=3^{n+1}\delta_W$.
\end{proof}

The proposition replaces a free search for subvarieties by a search for components of fixed complete intersections. It supplies no fixed value of $m$ in an unproved family. The total complete intersection always has zero Weil projection; the required phenomenon is a splitting into components with nonzero projections that cancel.

\section{An optional Galois-orbit criterion}
\label{sec:galois}
This section gives a sufficient arithmetic condition for a bounded detector locus to fill a full unitary family. It uses Andr\'e--Oort and is separate from the proof of Theorem~\ref{thm:global}.

Fix a geometrically connected fine unitary Shimura base $S$ over a number field $F$, its polarized abelian scheme, and a labelled integral order action in $K$, all over that same $F$. Suppose that $G^{\rm ad}$ is $\Q$-simple of type $A$, as for the full signature-$(n,n)$ unitary datum. A special point is \emph{relative Weyl} if its defining torus has Galois action containing the Weyl group of this datum \cite[Definitions~5.3--5.4]{ChaiOort}. This refers to the unitary group; a full Weyl CM field for the ambient Siegel variety is not required.

Choose a locally closed embedding $S\hookrightarrow\mathbf P^N$. For a reduced closed set $V$ let $\operatorname{tdeg}(\overline V)$ be the sum of the degrees of all geometric irreducible components of its projective closure, in every dimension. This convention matters: a line with arbitrarily many additional isolated points has fixed top-dimensional degree.

Andr\'e--Oort holds for every Cartesian power of $S$ by
Pila--Shankar--Tsimerman, with Esnault--Groechenig
\cite[Theorem~1.1]{PSTFull}. Scanlon's automatic uniformity theorem \cite[Theorem~4.2]{Scanlon} gives a function $B_S(H)$ bounding the number of special points on a closed locus of total degree at most $H$ that lie on no positive-dimensional special subvariety contained in that locus. To obtain this form, parameterize bounded-degree reduced closed sets by finitely many Chow families. Automatic uniformity bounds the families of Zariski closures of their special points, and hence bounds the numbers of their zero-dimensional components. Andr\'e--Oort identifies all the positive-dimensional components as special. The resulting function is qualitative. Jaswal's effective refinement
\cite[Theorem~1.1(ii)]{JaswalEffective} is a separate recent input;
it requires its stated zero-free half-plane and degree data. No
numerical rate for this Scanlon function is asserted.

\begin{proposition}\label{prop:galois-witness}
Let $\mathcal D_D=\{s:\delta_W(s)\leq D\}$ and choose any integer $H(D)$ at least $\operatorname{tdeg}(\overline{\mathcal D_D})$. Under the preceding hypotheses, a relative-Weyl special point $x$ satisfying
\[
 \delta_W(x)\leq D,\qquad
 |\operatorname{Gal}(\overline F/F)\cdot x|>B_S(H(D))
\]
forces $\mathcal D_D=S$. More generally, if $\mathcal D_D$ is proper, it contains at most $B_S(H(D))$ relative-Weyl special points.
\end{proposition}

\begin{proof}
First, the locus $\mathcal D_D$ descends over the fixed field $F$, independently of $D$. Choose one global order endomorphism $a$ with $a/\bar a$ not a root of unity. The integral operator $P(a^*+\bar a^*)$ in the proof of Theorem~\ref{thm:extraction} is defined over $F$ and detects exactly the Weil projection. On a projective parameter space for effective cycles of degree at most $D$, apply this operator to the universal class in one fixed $\ell$-adic cohomology theory. Its nonzero locus is a finite union of geometric connected components, stable under $\operatorname{Gal}(\overline F/F)$ because the operator and the class construction are defined over $F$. This open and closed subscheme descends over $F$; its proper image is precisely $\mathcal D_D$. A detecting effective cycle has a detecting irreducible component of no greater degree; conversely that irreducible component itself occurs in the parameter space. This is an equivalence of existence statements, not an assertion that detecting components can never cancel. Thus the entire $F$-orbit of $x$ is contained in $\mathcal D_D$.

A relative-Weyl special point lies on no proper positive-dimensional special subvariety of $S$. Its Mumford--Tate torus maps onto the full adjoint Weyl torus \cite[Remark~5.8]{ChaiOort}. The group of a special subvariety through the point contains that torus. Its roots are stable under Galois and hence under the full Weyl group. In type $A$ the roots of each absolute simple factor form one Weyl orbit; $\Q$-simplicity makes Galois transitive on the factors. The subgroup therefore has either no roots, giving a zero-dimensional domain, or all roots, giving the whole connected domain. This is the calculation in \cite[Lemma~6.3]{ChaiOort}.

The same property holds for Galois conjugates, because Galois conjugates of special subvarieties are special. If $\mathcal D_D$ were proper, every point in the indicated orbit would therefore be one of its isolated maximal special points. The bound $B_S(H(D))$ contradicts the assumed orbit size. The same reasoning counts all relative-Weyl points in a proper $\mathcal D_D$.
\end{proof}

The detector degree $D$, the degree $H(D)$ of its locus in the base, and the special-point bound $B_S(H(D))$ are different quantities. Closedness supplies no growth estimate for $H(D)$. Automatic uniformity supplies no effective estimate for $B_S$. Most importantly, an orbit lower bound supplies no detector membership. A pointwise detector of degree $D(x)$ is useful here only if the comparison with $B_S(H(D(x)))$ is proved.

Without the relative-Weyl restriction, orbit size alone is insufficient. In $Y(1)^2$, the fixed special line $\{0\}\times Y(1)$ contains the full Galois orbits of $(0,j)$ for singular moduli $j$, with unbounded orbit sizes. On a curve base, in contrast, every proper closed locus is zero-dimensional, and its total degree directly bounds the size of an orbit it contains.
\section{What one finite-field detector would have to certify}
\label{sec:one-prime}
This section distinguishes flatness over the arithmetic base from the relative flatness that would make one finite-field cycle propagate through a whole family.

Let $\mathcal S$ be an integral scheme of finite type, flat over a localization of the ring of integers of a number field. Let $\mathcal A\to\mathcal S$ be a polarized abelian scheme with a fixed integral quadratic action whose characteristic-zero family is of Weil type. Fix an auxiliary prime $\ell$ invertible on $\mathcal S$. The integral polynomial projector of Section~\ref{sec:weil} defines the Weil summand in the corresponding lisse $\Q_\ell$-sheaf. For $D\ge1$, let $H_D\to\mathcal S$ be a projective parameter space for effective middle-dimensional cycles of degree at most $D$. Such a space is obtained by the construction of Lemma~\ref{lem:parameters}, using single cycles instead of pairs.

\begin{proposition}\label{prop:flat-detector-witness}
Suppose a geometric point $h\in H_D$ parametrizes a cycle with nonzero Weil projection, and the morphism $H_D\to\mathcal S$ is flat at $h$. Then every complex fibre of the family has a detector of degree at most $D$. The point $h$ may have positive characteristic.
\end{proposition}

\begin{proof}
Applying the projector to the universal cycle class gives a lisse section on $H_D$. Its nonzero locus $H_D^{W}$ is a union of geometric connected components. It is open and closed, descends over the base by Galois invariance, and is proper over $\mathcal S$. A flat morphism of finite presentation is open near a point where it is flat: restrict to the open flat locus \cite[Theorems~2.4.6 and 11.1.1]{EGAIV}. Thus the image of $H_D^{W}$ contains a nonempty open subset of $\mathcal S$. Properness makes the image closed, and integrality of $\mathcal S$ makes that image all of $\mathcal S$. In each complex fibre, nonzero projection of an effective cycle implies nonzero projection of at least one of its irreducible components, of degree no larger than $D$.
\end{proof}

For example, if a local-complete-intersection detector $Y$ on a geometric finite-field fibre has $H^1(Y,N_{Y/\mathcal A_s})=0$, the usual embedded deformation theory makes the relative Hilbert scheme smooth at $[Y]$. This gives the proposition's flatness hypothesis. The vanishing is sufficient, not necessary; relative flatness can hold even when this obstruction space is nonzero. No such detector with the required relative deformation property is constructed here in an unresolved eightfold family.

Flatness only over the ring of integers does not suffice. Consider the proper closed immersion
\[
 H=\{x=0\}\longrightarrow\mathcal S=\mathbf A^1_\Z.
\]
The scheme $H\cong\Spec\Z$ is flat over $\Z$ and has no vertical components. For any rational prime $p$, the section $s_p:\Spec\Z\to\mathcal S$ defined by $x=p$ has fibre product
\[
 H\times_{\mathcal S,s_p}\Spec\Z\cong\Spec\mathbf F_p.
\]
It has a point at $p$ and no characteristic-zero point. Since $p$ can be arbitrarily large, discarding a fixed set of vertical primes of the universal parameter space does not solve this problem. The obstruction is failure of flatness over the moduli base along the chosen section.

For a fixed CM fibre, Proposition~\ref{prop:lift} still gives a finite exceptional set depending on that fibre and on $D$. A degree-$D$ detector at a place outside that set suffices for that fibre. Proposition~\ref{prop:flat-detector-witness} is a different sufficient condition: a single detector together with relative flatness makes an entire parameter component dominate the moduli base. A construction must establish one of these properties; flatness over the arithmetic base alone establishes neither.

\section{What remains open}\label{sec:open}
The equivalence leaves a precise existence problem: bound detector
degrees uniformly on CM points of the principal hyperbolic family
$U^n$, for each imaginary quadratic field and each dimension.
The first unresolved case beyond the cited constructions is a split
eightfold over a field other than $\Q(\sqrt{-3})$, for example $\Q(i)$.
The Eisenstein eightfold estimate uses the dominance theorem proved
above and the cyclic-cover transfer. No case in dimension at least ten follows
from these arguments.

The pure-character method requires an equivariant gluable perfect
object with invariant second self-Ext dimension $56$ in dimension
eight. The exact dimension criterion does not construct one. The
complete-intersection reformulation requires a fixed multiple $mL$
with a detecting component on a dense set; the total intersection
has zero Weil projection, so the needed splitting is substantive.
The one-prime criterion requires a detector outside the finite
exceptional set for that fibre. A uniform estimate of those sets and
a finite-field construction meeting it remain separate missing inputs.

The companion obstructions concern specified quadric, support, bundle
and finite-cover constructions. They do not exclude all cycles.
Pointwise algebraicity at CM points does not replace the uniform bound,
and neither the Hodge conjecture nor its reduction to a finite
computation is proved here.

\section*{Method, checks and records}
The mathematics and text were produced by GPT-6 Astra (ChatGPT Pro
mode). Separate re-proofs by the same model are distinguished in the
records from targeted proof checks and exact computations; the interface
did not expose the reasoning-effort setting. Earlier proofs received
an outside reading by Claude Opus 5.5 at effort high, which found no
error in the statements identified in its records. That reading does
not certify every extension made here. No human has checked this work.
The text and original programs are dedicated to the public domain
under CC0 1.0.

The programs, exact commands, source locators and check records are
collected in the \emph{Algebraic cycles and motives workbench},
\url{\workbenchurl}. Their permanent entry point is
\url{\recordurl}. The records distinguish conditional conclusions
from open existence problems.
\begin{thebibliography}{99}
\bibitem{BS} L.~Br\"unjes and C.~Serp\'e, \emph{Uniform bounds and ultraproducts of cycles}, arXiv:0901.4853 (2009), Definition~2.3(ii), p.~4, and Theorem~3.4, p.~8. \url{https://arxiv.org/abs/0901.4853}.
\bibitem{FGA} A.~Grothendieck, \emph{Techniques de construction et th\'eor\`emes d'existence en g\'eom\'etrie alg\'ebrique IV: les sch\'emas de Hilbert}, S\'eminaire Bourbaki, Expos\'e~221 (1961), in \emph{Fondements de la g\'eom\'etrie alg\'ebrique}, Theorem~3.2, p.~260, and pp.~265--266.
\bibitem{SGA4} M.~Artin, A.~Grothendieck and J.-L.~Verdier (eds.), \emph{Th\'eorie des topos et cohomologie \'etale des sch\'emas}, SGA~4, Lecture Notes in Mathematics 269, 270, 305, Springer (1972--1973): XI, Theorem~4.4; XVI, Corollary~2.2; XVIII, Theorems~2.9 and 3.2.5.
\bibitem{SGA45} P.~Deligne, \emph{Cohomologie \'etale}, SGA~$4\frac12$, Lecture Notes in Mathematics 569, Springer (1977), expos\'e \emph{Cycle}, \S2.3.1--2.3.8; erratum to SGA~4 XVI, Corollary~2.2.
\bibitem{SGA6} P.~Berthelot, A.~Grothendieck and L.~Illusie (eds.), \emph{Th\'eorie des intersections et th\'eor\`eme de Riemann--Roch}, SGA~6, Lecture Notes in Mathematics 225, Springer (1971), Expos\'e~XIII (S.~Kleiman), Definition~6.9 and Corollary~6.11, pp.~658--660.
\bibitem{MilneSV} J.~S.~Milne, \emph{Introduction to Shimura varieties}, in \emph{Harmonic Analysis, the Trace Formula, and Shimura Varieties}, Clay Mathematics Proceedings 4 (2005), 265--378; revised notes (2017), Theorem~5.4 and \S\S5--7. \url{https://www.jmilne.org/math/xnotes/svi.pdf}.
\bibitem{CDK} E.~Cattani, P.~Deligne and A.~Kaplan, \emph{On the locus of Hodge classes}, J. Amer. Math. Soc. 8 (1995), 483--506, Introduction and Theorem~1.1. \url{https://arxiv.org/abs/alg-geom/9402009}.
\bibitem{Voisin07} C.~Voisin, \emph{Hodge loci and absolute Hodge classes}, Compos. Math. 143 (2007), 945--958, Propositions~0.2 and 0.5, Theorem~0.6, pp.~946--947. \url{https://doi.org/10.1112/S0010437X07002837}.
\bibitem{BBDense} D.~Bhattacharjee and U.~Bhattacharya, \emph{Dense Algebraic Loci of Hodge Classes}, Hexagon 2610.00315v2 (2026), Lemmas~4.2--4.3 and Theorem~4.5, pp.~10--12. Unrefereed preprint. \url{https://hexagonmath.org/2610.00315v2}.
\bibitem{BBAlgebraic} D.~Bhattacharjee and U.~Bhattacharya, \emph{Algebraic Loci of Weil Classes on Abelian Varieties}, companion manuscript (2026), Theorems~16.9 and 16.17. Unrefereed. \url{https://doi.org/10.5281/zenodo.23885230}.
\bibitem{OpenAICM} OpenAI, \emph{The rational Hodge conjecture for CM abelian varieties}, OpenAI Math Release, 6 October 2026, 53 pp., Theorem~1.1; Lemma~2.2; Proposition~3.1; Corollary~6.7. Unrefereed preprint. \url{https://github.com/openai/math/tree/main/preprints/The-rational-Hodge-conjecture-for-CM-abelian-varieties-October-6-2026}.
\bibitem{vHS} P.~van Hoften and J.~Sempliner, \emph{The Hodge Conjecture for CM abelian varieties}, Hexagon 2610.00259v2, 15 October 2026, 25 pp., Theorem~I; Lemma~2.3.2; Theorem~2.3.4. Unrefereed preprint. \url{https://hexagonmath.org/2610.00259v2}.
\bibitem{Deligne82} P.~Deligne, notes by J.~S.~Milne, \emph{Hodge cycles on abelian varieties}, in \emph{Hodge Cycles, Motives, and Shimura Varieties}, LNM 900, Springer (1982), 9--100. Revised transcription (2018), Proposition~4.1 and Corollary~4.2, p.~28; Proposition~6.1, pp.~41--42. \url{https://www.jmilne.org/math/Documents/Deligne82.pdf}.
\bibitem{MilneAV} J.~S.~Milne, \emph{Abelian varieties}, in \emph{Arithmetic Geometry}, Springer (1986), 103--150; corrected notes (2022), Remark~7.3, p.~12; Theorem~13.3(b), p.~25; \S17. \url{https://www.jmilne.org/math/xnotes/AVs.pdf}.
\bibitem{Ballmann} W.~Ballmann, \emph{Lectures on K\"ahler Manifolds}, ESI Lectures in Mathematics and Physics, EMS (2006), Theorem~5.42 and equation~(5.43), pp.~70--71. \url{https://ems.press/books/esi/19}.
\bibitem{Schoen98} C.~Schoen, \emph{Addendum to: ``Hodge classes on self-products of a variety with an automorphism''}, Compos. Math. 114 (1998), 329--336, Proposition~10, pp.~332--333. \url{https://doi.org/10.1023/A:1000566205021}.
\bibitem{Schoen07} C.~Schoen, \emph{A family of surfaces constructed from genus 2 curves}, Internat. J. Math. 18 (2007), 585--612. \url{https://doi.org/10.1142/S0129167X07004175}.
\bibitem{vG96} B.~van Geemen, \emph{Theta functions and cycles on some abelian fourfolds}, Math. Z. 221 (1996), 617--631, \S2.1; Lemma~3.3; Proposition~3.6; Theorem~3.7. \url{https://sites.unimi.it/vangeemen/ThetafctsCycles4folds_MathZ1996.pdf}.
\bibitem{vG22} B.~van Geemen, \emph{Weil Classes and Decomposable Abelian Fourfolds}, SIGMA 18 (2022), 097, Introduction and \S\S6.1--6.3. \url{https://doi.org/10.3842/SIGMA.2022.097}.
\bibitem{RRS} C.~Rito, X.~Roulleau and A.~Sarti, \emph{Explicit Schoen surfaces}, Algebraic Geometry 6 (2019), 410--426, Introduction and Theorem~1. \url{https://arxiv.org/abs/1609.02235}.
\bibitem{Beauville} A.~Beauville, \emph{A tale of two surfaces}, in \emph{Higher Dimensional Algebraic Geometry}, Adv. Stud. Pure Math. 74 (2017), 1--10, \S3. \url{https://math.univ-cotedazur.fr/u/beauvill/pubs/Tale2.pdf}.
\bibitem{VoisinB} C.~Voisin, \emph{La conjecture de Hodge pour les vari\'et\'es ab\'eliennes de dimension au plus 5}, S\'eminaire Bourbaki, Expos\'e 1248, January 2026, Th\'eor\`emes~1.3--1.4, p.~1248--03; \S2.3; Lemma~2.9 and Corollary~2.10, p.~1248--16. Official seminar manuscript. \url{https://www.bourbaki.fr/TEXTES/Exp1248-Voisin.pdf}.
\bibitem{MarkmanICM} E.~Markman, \emph{Secant Sheaves and Weil Classes on Abelian Varieties}, Proceedings of the International Congress of Mathematicians 2026, Vol.~3, 586--602, Theorem~1.2 and Corollary~1.3, p.~587. \url{https://doi.org/10.1137/25M1803796}. Detailed proof: \emph{Cycles on abelian $2n$-folds of Weil type from secant sheaves on abelian $n$-folds}, arXiv:2502.03415v2 (2025), Theorem~1.5.1. \url{https://arxiv.org/abs/2502.03415}.
\bibitem{Andre96} Y.~Andr\'e, \emph{Pour une th\'eorie inconditionnelle des motifs}, Publ. Math. IHES 83 (1996), 5--49, Theorems~0.5 and 0.6.2, pp.~8--9; Lemma~6.3.2, p.~32. \url{https://doi.org/10.1007/BF02698643}.
\bibitem{Scanlon} T.~Scanlon, \emph{Automatic uniformity}, Internat. Math. Res. Notices 2004, no.~62, 3317--3326, Definition~2.2 and Theorem~4.2. \url{https://math.berkeley.edu/~scanlon/papers/autouni.pdf}.
\bibitem{ChaiOort} C.-L.~Chai and F.~Oort, \emph{Abelian varieties isogenous to a Jacobian}, Ann. of Math. 176 (2012), 589--635, Definitions~5.3--5.4, Remark~5.8 and Lemma~6.3, pp.~607,610 and 621--622. \url{https://doi.org/10.4007/annals.2012.176.1.11}.
\bibitem{Perry} A.~Perry, \emph{The semiregularity theorem for equivariant noncommutative varieties}, arXiv:2604.00511v2 (24 June 2026), Theorem~1.2; Definition~2.6; Lemma~3.3; \S6.2. Unrefereed preprint. \url{https://arxiv.org/abs/2604.00511v2}.
\bibitem{BF} R.-O.~Buchweitz and H.~Flenner, \emph{A semiregularity map for modules and applications to deformations}, Compos. Math. 137 (2003), 135--210, pp.~136--137; Theorem~5.1, pp.~174--175, proof p.~179; Proposition~5.9 and Lemma~5.10, pp.~178--179; \S8. \url{https://doi.org/10.1023/A:1023999012081}.
\bibitem{Schoen88} C.~Schoen, \emph{Hodge classes on self-products of a variety with an automorphism}, Compos. Math. 65 (1988), 3--32, Theorem~2.0, p.~11; Lemmas~2.2 and 2.6, pp.~13 and 16; \S3, pp.~24--25. \url{https://www.numdam.org/article/CM_1988__65_1_3_0.pdf}.
\bibitem{Koike} K.~Koike, \emph{Algebraicity of some Weil Hodge Classes}, Canad. Math. Bull. 47 (2004), 566--572, Theorem~2.1, p.~568. \url{https://doi.org/10.4153/CMB-2004-055-x}.
\bibitem{EGAIV} A.~Grothendieck and J.~Dieudonn\'e, \emph{\'El\'ements de g\'eom\'etrie alg\'ebrique IV}, Publ. Math. IHES 24 (1965) and 28 (1966), Theorems~2.4.6 and 11.1.1.
\bibitem{Mukai} S.~Mukai, \emph{Duality between $D(X)$ and $D(\widehat X)$ with its application to Picard sheaves}, Nagoya Math. J. 81 (1981), 153--175, Theorem~2.2, p.~156; \S3. \url{https://doi.org/10.1017/S002776300001922X}.
\bibitem{BFHH} R.-O.~Buchweitz and H.~Flenner, \emph{The global decomposition theorem for Hochschild (co-)homology of singular spaces via the Atiyah--Chern character}, Adv. Math. 217 (2008), 243--281, Proposition~6.2.1, Corollary~6.3.2 and Theorem~6.4.2; author version, pp.~29--31. \url{https://arxiv.org/abs/math/0606730}.
\bibitem{Hotchkiss} J.~Hotchkiss, \emph{Hodge theory of twisted derived categories and the period-index problem}, arXiv:2212.10638v1 (20 December 2022), Definitions~4.5 and 4.16, Lemma~4.6, Remark~4.13 and Theorem~4.17, pp.~15--18. Unrefereed author version. \url{https://arxiv.org/abs/2212.10638v1}.
\bibitem{DettweilerWewersPoincare}
M.~Dettweiler and S.~Wewers, \emph{Variation of parabolic cohomology
and Poincar\'e duality}, S\'eminaires et Congr\`es \textbf{13}
(2006), 145--164. Published Theorem~2.6, p.~158;
\url{https://www.emis.de/ft/47210}.
The author version, arXiv:math/0411119v1, numbers this formula
Theorem~2.5, p.~12. \url{https://arxiv.org/abs/math/0411119v1}.

\bibitem{MarkmanSecant} E.~Markman, \emph{Cycles on abelian $2n$-folds of Weil type from secant sheaves on abelian $n$-folds}, arXiv:2502.03415v2 (2025), Theorem~1.5.1 and \S\S8--9. Unrefereed detailed preprint. \url{https://arxiv.org/abs/2502.03415v2}.
\bibitem{Severn} B.~Severn, \emph{Weil classes on abelian sixfolds and eightfolds of Weil type for $\Q(\sqrt{-3})$}, Hexagon 2610.00387v2, served manuscript dated 21 October 2026, internal version 1.5, 31 pp. Lemma~4.4, pp.~9--10; Proposition~4.6, pp.~11--12; Theorem~5.1, pp.~12--13; Proposition~8.5, pp.~20--21; \S9 and Appendix~A. Unrefereed preprint. \url{https://hexagonmath.org/2610.00387v2}.
\bibitem{MilneLefschetz} J.~S.~Milne, \emph{Lefschetz classes on abelian varieties}, Duke Math. J. 96 (1999), 639--675. Author version, Lemma~3.5, pp.~15--16; Proposition~3.6 and Remark~3.7, pp.~17--18. \url{https://jmilne.org/math/articles/1999aP.pdf}.
\bibitem{BBSecantObjects} D.~Bhattacharjee and U.~Bhattacharya,
\emph{Secant Objects for Weil Classes on Abelian Varieties},
Preprints.org 202611.0672, version~1 (6 November 2026),
\url{https://www.preprints.org/manuscript/202611.0672}.
Unrefereed. Lemmas~2--3 and Proposition~2; Theorem~10 and its
analytic deformation step are discussed in Section~\compref{sec:analytic-weil-sheaves}.
\bibitem{CompanionB} \emph{Obstructions to cycle constructions and criteria for Mumford's exceptional classes}, companion paper, numbered references prefixed B. \url{\workbenchurl}.
\bibitem{OpenAIEightWithdrawn} OpenAI, \emph{Algebraicity of Weil classes on split abelian eightfolds}, withdrawal notice, 6 October 2026. The notice identifies a sign error in the stabilization-trace argument; the claimed algebraicity theorem is not established by that manuscript. \url{https://github.com/openai/math/blob/main/preprints/Algebraicity-of-Weil-classes-on-split-abelian-eightfolds-September-18-2026/README.md}.

\bibitem{SevernGaussian} B.~Severn,
\emph{Weil classes on abelian sixfolds of Weil type for $\Q(i)$ with
discriminant $-3$}, Hexagon 2610.00477v1, 22 October 2026, 53 pp.
Theorems~A--C, Sections~3--7 and Appendices~A--B. Unrefereed preprint.
\url{https://hexagonmath.org/2610.00477v1}; certificate deposit
\url{https://doi.org/10.5281/zenodo.23549514}.
\bibitem{ShiodaFermat} T.~Shioda, \emph{The Hodge conjecture for Fermat
varieties}, Math. Ann. \textbf{245} (1979), 175--184. Theorems~I--IV;
condition $(P_{12})$, p.~182. \url{https://doi.org/10.1007/BF01428804}.
\bibitem{LooijengaFermat} E.~Looijenga, \emph{Fermat varieties and the
periods of some hypersurfaces}, Adv. Stud. Pure Math. \textbf{58}
(2010), 47--67. Corollary~2.2.
\url{https://webspace.science.uu.nl/~looij101/fermatperiod4.pdf}.
\bibitem{TruongMoving} T.~T.~Truong, \emph{Relative dynamical degrees
of correspondences over a field of arbitrary characteristic},
Lemma~2.2, author preprint \url{https://arxiv.org/abs/1605.05049}.
\bibitem{CaldararuWillerton} A.~C\u{a}ld\u{a}raru and S.~Willerton,
\emph{The Mukai pairing. I. A categorical approach},
New York J. Math. \textbf{16} (2010), 61--98. Proposition~11,
p.~83; \S7.3, pp.~88--89.
\url{https://nyjm.albany.edu/j/2010/16-6.html}.
\bibitem{Pridham} J.~P.~Pridham, \emph{Semiregularity as a consequence
of Goodwillie's theorem}, Forum Math. Sigma \textbf{12} (2024), e126.
Corollaries~2.24--2.25. \url{https://doi.org/10.1017/fms.2024.132}.
\bibitem{WeiCohesive} Z.~Wei, \emph{Deformations of cohesive modules
on compact complex manifolds}, Differential Geom. Appl.
\textbf{101} (2025), 102306. Theorem~3.11 and \S5, author version
arXiv:2308.01266v2. \url{https://doi.org/10.1016/j.difgeo.2025.102306}.
\bibitem{ChanSuen} K.~Chan and Y.-H.~Suen,
\emph{A differential-geometric approach to deformations of pairs
$(X,E)$}, author version arXiv:1406.6753v4, \S\S3,5--6.
\url{https://arxiv.org/abs/1406.6753v4}.
\bibitem{GoldmanMillson} W.~M.~Goldman and J.~J.~Millson,
\emph{The deformation theory of representations of fundamental groups
of compact K\"ahler manifolds}, Publ. Math. IHES \textbf{67}
(1988), 43--96. \S1.1, p.~49; \S\S2--3.
\url{https://doi.org/10.1007/BF02699127}.
\bibitem{GaussianLandherr} W.~Landherr,
\emph{\"Aquivalenz Hermitescher Formen \"uber einem beliebigen
algebraischen Zahlk\"orper}, Abh. Math. Sem. Univ. Hamburg
\textbf{11} (1935), 245--248, especially pp.~246--247.
\url{https://webhomes.maths.ed.ac.uk/~v1ranick/papers/landherr.pdf}.
\bibitem{GaussianKirschmer} M.~Kirschmer,
\emph{Determinant groups of hermitian lattices over local fields},
Arch. Math. \textbf{113} (2019), 337--347.
Theorem~4.3 and Remark~4.8(2), author pp.~7--9.
\url{https://doi.org/10.1007/s00013-019-01348-z};
\url{https://www.math.rwth-aachen.de/~Markus.Kirschmer/papers/det.pdf}.
\bibitem{PSTFull} J.~Pila, A.~Shankar and J.~Tsimerman, with an appendix by H.~Esnault and M.~Groechenig, \emph{Canonical heights on Shimura varieties and the Andr\'e--Oort conjecture}, arXiv:2109.08788, Theorem~1.1. \url{https://arxiv.org/abs/2109.08788}.
\bibitem{JaswalEffective} M.~Jaswal, \emph{An effective Andr\'e--Oort theorem for Shimura varieties}, Hexagon 2611.00339v3, Theorem~1.1(ii), Corollary~1.2. Unrefereed. \url{https://hexagonmath.org/2611.00339v3}.

\bibitem{Davis08}
D.~M. Davis, \emph{Homotopy type and $v_1$-periodic homotopy groups of
$p$-compact groups}, Topology Appl. \textbf{156} (2008), 300--321,
Theorem~5.1 and its proof, pp.~312--313.
\url{https://www.lehigh.edu/~dmd1/TOPOL3486.pdf}.

\bibitem{FredianiGhigiPenegini15}
P.~Frediani, A.~Ghigi and M.~Penegini, \emph{Shimura varieties in the
Torelli locus via Galois coverings}, Int. Math. Res. Not. IMRN
\textbf{2015}, no.~20, 10595--10623, Theorem~2.10;
arXiv:1402.0973v3, p.~8.
\url{https://doi.org/10.1093/imrn/rnu272}.
\url{https://arxiv.org/abs/1402.0973v3}.

\bibitem{DettweilerWewers04}
M.~Dettweiler and S.~Wewers, \emph{Variation of local systems and
parabolic cohomology}, Israel J. Math. \textbf{156} (2006), 157--185;
arXiv:math/0310139v2, Lemma~2.2, p.~7, and Theorem~2.5, p.~12.
\url{https://doi.org/10.1007/BF02773830}.
\url{https://arxiv.org/abs/math/0310139v2}.

\bibitem{Schmid73}
W.~Schmid, \emph{Variation of Hodge structure: the singularities of
the period mapping}, Invent. Math. \textbf{22} (1973), 211--319,
Theorems~(4.12), p.~233, and~(6.16), p.~255.
\url{https://doi.org/10.1007/BF01389674}.

\bibitem{SevernData}
B.~Severn, \emph{Weil classes on abelian sixfolds and eightfolds of
Weil type for $\mathbf Q(\sqrt{-3})$}, Hexagon 2610.00387v2,
internal version~1.5, 21 October 2026, Appendix~A.
Unrefereed manuscript; cited for the numerical data only.
The portal heading mentions sixfolds; the title here is the actual
PDF and TeX title.
\url{https://hexagonmath.org/2610.00387v2}.



\bibitem{SerreRepresentations} J.-P.~Serre,
\emph{Linear representations of finite groups}, GTM~42,
Springer, 1977, \S10.5, Theorem~20, p.~78.
\url{https://doi.org/10.1007/978-1-4684-9458-7}.

\end{thebibliography}


\end{document}
